% 代数与几何课程笔记 —— 主文档（正文排版集中在 0.样式.tex，章节文件只放内容）
% 编译：xelatex 1.main.tex 两遍；成品输出见 build.ps1
% 章序对齐教材：同济《工程数学·线性代数》第七版（2023），第 1–6 章 + 附录
\documentclass[11pt,a4paper,oneside,openany]{book}
\usepackage[UTF8,fontset=windows,scheme=plain]{ctex}
\usepackage{fontspec}
\setmainfont{Times New Roman}
\setsansfont{Microsoft YaHei}

\newcommand{\YHBbooktrue}{}
\newcommand{\YHBcourse}{代数与几何}
\newcommand{\YHBversion}{v1.2}
\newcommand{\YHBquote}{以矩阵刻画变换，以向量描摹空间。}
\newcommand{\YHBcoverSize}{38pt}
\newcommand{\YHBcoverLead}{46pt}
\usepackage{xcolor}
\definecolor{yhbAccent}{HTML}{3F3350}
\definecolor{yhbAccent2}{HTML}{7A6386}
% ============================================================
% 0.样式.tex —— 课程笔记统一正文排版
% 用法：主文件先定义课程信息与主色，再 % ============================================================
% 0.样式.tex —— 课程笔记统一正文排版
% 用法：主文件先定义课程信息与主色，再 % ============================================================
% 0.样式.tex —— 课程笔记统一正文排版
% 用法：主文件先定义课程信息与主色，再 \input{0.样式.tex}
%   \newcommand{\YHBcourse}{微积分A}      % 课程名（页眉与书签用）
%   \newcommand{\YHBversion}{v1.2}        % 版本号（封面用）
%   \newcommand{\YHBquote}{……}            % 封面题句
%   \definecolor{yhbAccent}{HTML}{24465C} % 主色（深）
%   \definecolor{yhbAccent2}{HTML}{347A80}% 辅色（浅）
%   \YHBbooktrue / \YHBbookfalse          % 是否 book 文档类（有 \chapter）
% 章节文件只写内容，不要重复 \usepackage、不要改字号。
% ============================================================

% ---------- 页面与字体 ----------
\usepackage{geometry}
\geometry{a4paper,left=24mm,right=24mm,top=23mm,bottom=24mm,headheight=15pt,headsep=7mm,footskip=13mm}
\usepackage{microtype}
\usepackage{amsmath,amssymb,amsfonts,mathtools}
\usepackage{esint}
\usepackage{graphicx}
\usepackage{adjustbox}
\usepackage{array,booktabs,multirow,makecell,tabularx,longtable}
\usepackage{xcolor,colortbl}
\usepackage{enumitem}
\usepackage{needspace,etoolbox}
\usepackage{titlesec,titletoc}
\usepackage{fancyhdr}
\usepackage{caption}
\usepackage[most]{tcolorbox}

% ---------- 颜色别名（章节里用的旧名继续有效） ----------
\colorlet{mainblue}{yhbAccent}
\colorlet{secondcyan}{yhbAccent2}
\colorlet{thirdteal}{yhbAccent2}
\colorlet{highlight}{yhbAccent!6!white}
\colorlet{tablehead}{yhbAccent!12!white}
\definecolor{textgray}{HTML}{4A5057}

% ---------- 课程标签 ----------
\providecommand{\YHBdoclabel}{课程笔记}
% ---------- 正文节奏 ----------
\linespread{1.3}
\setlength{\parindent}{0pt}
\setlength{\parskip}{6pt plus 1.5pt minus 1.5pt}
\setlength{\emergencystretch}{2.5em}
\raggedbottom
\widowpenalty=10000
\clubpenalty=10000
\setlength{\tabcolsep}{7pt}
\renewcommand{\arraystretch}{1.3}
\allowdisplaybreaks[2]
\AtBeginDocument{%
  \setlength{\abovedisplayskip}{10pt plus 2pt minus 2pt}%
  \setlength{\belowdisplayskip}{10pt plus 2pt minus 2pt}%
  \setlength{\abovedisplayshortskip}{6pt}%
  \setlength{\belowdisplayshortskip}{6pt}%
}
\captionsetup{font=small,labelfont=bf,skip=6pt}
\renewcommand{\tablename}{表}
\renewcommand{\figurename}{图}

% ---------- 列表 ----------
\setlist{leftmargin=2em,labelsep=0.6em,itemsep=4pt,parsep=2pt,topsep=6pt,partopsep=0pt}
\setlist[itemize,2]{label=\textbullet,leftmargin=1.6em}

% ---------- 标题层级 ----------
\ifdefined\YHBbooktrue
  \titleformat{\chapter}[display]
    {\sffamily\bfseries\color{yhbAccent}}
    {\normalsize\sffamily 第\,\thechapter\,章}
    {6pt}
    {\LARGE}
    [{\vspace{5pt}{\color{yhbAccent!55}\titlerule[1pt]}}]
  \titlespacing*{\chapter}{0pt}{14pt}{24pt}
  \titleformat{\section}{\Large\sffamily\bfseries\color{yhbAccent}}{\thesection}{0.7em}{}
  \titlespacing*{\section}{0pt}{18pt plus 3pt minus 2pt}{8pt}
\else
  \titleformat{\section}{\Large\sffamily\bfseries\color{yhbAccent}}{\thesection}{0.7em}{}[{\vspace{3pt}{\color{yhbAccent!45}\titlerule[0.7pt]}}]
  \titlespacing*{\section}{0pt}{18pt plus 3pt minus 2pt}{9pt}
\fi
\titleformat{\subsection}{\large\sffamily\bfseries\color{yhbAccent!85!black}}{\thesubsection}{0.7em}{}
\titleformat{\subsubsection}{\normalsize\bfseries\color{black!85}}{\thesubsubsection}{0.7em}{}
\titlespacing*{\subsection}{0pt}{12pt plus 2pt minus 1pt}{5pt}
\titlespacing*{\subsubsection}{0pt}{9pt plus 2pt minus 1pt}{4pt}
\pretocmd{\section}{\Needspace{6\baselineskip}}{}{}
\pretocmd{\subsection}{\Needspace{5\baselineskip}}{}{}
\pretocmd{\subsubsection}{\Needspace{4\baselineskip}}{}{}

\ifdefined\YHBbooktrue
  \setcounter{secnumdepth}{2}
  \setcounter{tocdepth}{1}
\else
  \setcounter{secnumdepth}{3}
  \setcounter{tocdepth}{2}
\fi
\renewcommand{\contentsname}{目录}
% 目录编号颜色由 titletoc 的 titlecontents 控制，不再重定义 \numberline

% ---------- 目录样式 ----------
\ifdefined\YHBbooktrue
  \titlecontents{chapter}[0pt]
    {\addvspace{6pt}\sffamily\bfseries\color{yhbAccent}}
    {\contentslabel{2.2em}}{}{\titlerule*[0.6em]{.}\contentspage}
  \titlecontents{section}[2.2em]
    {\addvspace{2pt}}
    {\contentslabel{2.2em}}{}{\titlerule*[0.6em]{.}\contentspage}
\else
  \titlecontents{section}[0pt]
    {\addvspace{5pt}\sffamily\bfseries\color{yhbAccent}}
    {\contentslabel{2.2em}}{}{\titlerule*[0.6em]{.}\contentspage}
  \titlecontents{subsection}[2.2em]
    {\addvspace{2pt}}
    {\contentslabel{2.4em}}{}{\titlerule*[0.6em]{.}\contentspage}
\fi

% ---------- 页眉页脚 ----------
\pagestyle{fancy}
\fancyhf{}
\fancyhead[L]{\footnotesize\sffamily\color{textgray}\YHBcourse\ \YHBdoclabel}
\ifdefined\YHBbooktrue
  \fancyhead[R]{\footnotesize\sffamily\color{textgray}\nouppercase{\leftmark}}
  \renewcommand{\chaptermark}[1]{\markboth{#1}{}}
\else
  \fancyhead[R]{\footnotesize\sffamily\color{textgray}\nouppercase{\rightmark}}
  \renewcommand{\sectionmark}[1]{\markright{#1}}
\fi
\fancyfoot[C]{\footnotesize\color{textgray}\thepage}
\renewcommand{\headrulewidth}{0.35pt}
\renewcommand{\headrule}{{\color{yhbAccent!45}\hrule width\headwidth height\headrulewidth}}
\fancypagestyle{plain}{%
  \fancyhf{}%
  \fancyfoot[C]{\footnotesize\color{textgray}\thepage}%
  \renewcommand{\headrulewidth}{0pt}%
}

% ---------- 公式与考点框 ----------
\newcommand{\keyformula}[1]{%
  \begin{tcolorbox}[enhanced,breakable,colback=highlight,colframe=yhbAccent!30!white,boxrule=0.5pt,arc=2.5pt,left=11pt,right=11pt,top=9pt,bottom=9pt,boxsep=0pt,before skip=11pt,after skip=11pt,halign=center,width=\linewidth]
  \adjustbox{max width=\linewidth-2em}{$\displaystyle #1$}
  \end{tcolorbox}}
\newcommand{\corepoint}[1]{%
  \begin{tcolorbox}[enhanced,breakable,colback=highlight,colframe=yhbAccent2,boxrule=0pt,borderline west={2.4pt}{0pt}{yhbAccent2},arc=0pt,left=11pt,right=11pt,top=8pt,bottom=8pt,before skip=11pt,after skip=11pt]
  \textbf{\color{yhbAccent}核心考点}\quad #1
  \end{tcolorbox}}
\newcommand{\symb}[2]{\textbf{#1}：#2\par}

% ---------- 超链接（务必最后加载） ----------
\usepackage[unicode=true,colorlinks=true,linkcolor=yhbAccent,citecolor=yhbAccent,urlcolor=yhbAccent!85!black,pdfborder={0 0 0},bookmarksnumbered=true,pdfstartview=FitH]{hyperref}
\usepackage{hypcap}
\hypersetup{pdftitle={\YHBcourse 课程笔记},pdfauthor={浮萍},pdfsubject={课程学习与复习笔记},pdfcreator={XeLaTeX}}

%   \newcommand{\YHBcourse}{微积分A}      % 课程名（页眉与书签用）
%   \newcommand{\YHBversion}{v1.2}        % 版本号（封面用）
%   \newcommand{\YHBquote}{……}            % 封面题句
%   \definecolor{yhbAccent}{HTML}{24465C} % 主色（深）
%   \definecolor{yhbAccent2}{HTML}{347A80}% 辅色（浅）
%   \YHBbooktrue / \YHBbookfalse          % 是否 book 文档类（有 \chapter）
% 章节文件只写内容，不要重复 \usepackage、不要改字号。
% ============================================================

% ---------- 页面与字体 ----------
\usepackage{geometry}
\geometry{a4paper,left=24mm,right=24mm,top=23mm,bottom=24mm,headheight=15pt,headsep=7mm,footskip=13mm}
\usepackage{microtype}
\usepackage{amsmath,amssymb,amsfonts,mathtools}
\usepackage{esint}
\usepackage{graphicx}
\usepackage{adjustbox}
\usepackage{array,booktabs,multirow,makecell,tabularx,longtable}
\usepackage{xcolor,colortbl}
\usepackage{enumitem}
\usepackage{needspace,etoolbox}
\usepackage{titlesec,titletoc}
\usepackage{fancyhdr}
\usepackage{caption}
\usepackage[most]{tcolorbox}

% ---------- 颜色别名（章节里用的旧名继续有效） ----------
\colorlet{mainblue}{yhbAccent}
\colorlet{secondcyan}{yhbAccent2}
\colorlet{thirdteal}{yhbAccent2}
\colorlet{highlight}{yhbAccent!6!white}
\colorlet{tablehead}{yhbAccent!12!white}
\definecolor{textgray}{HTML}{4A5057}

% ---------- 课程标签 ----------
\providecommand{\YHBdoclabel}{课程笔记}
% ---------- 正文节奏 ----------
\linespread{1.3}
\setlength{\parindent}{0pt}
\setlength{\parskip}{6pt plus 1.5pt minus 1.5pt}
\setlength{\emergencystretch}{2.5em}
\raggedbottom
\widowpenalty=10000
\clubpenalty=10000
\setlength{\tabcolsep}{7pt}
\renewcommand{\arraystretch}{1.3}
\allowdisplaybreaks[2]
\AtBeginDocument{%
  \setlength{\abovedisplayskip}{10pt plus 2pt minus 2pt}%
  \setlength{\belowdisplayskip}{10pt plus 2pt minus 2pt}%
  \setlength{\abovedisplayshortskip}{6pt}%
  \setlength{\belowdisplayshortskip}{6pt}%
}
\captionsetup{font=small,labelfont=bf,skip=6pt}
\renewcommand{\tablename}{表}
\renewcommand{\figurename}{图}

% ---------- 列表 ----------
\setlist{leftmargin=2em,labelsep=0.6em,itemsep=4pt,parsep=2pt,topsep=6pt,partopsep=0pt}
\setlist[itemize,2]{label=\textbullet,leftmargin=1.6em}

% ---------- 标题层级 ----------
\ifdefined\YHBbooktrue
  \titleformat{\chapter}[display]
    {\sffamily\bfseries\color{yhbAccent}}
    {\normalsize\sffamily 第\,\thechapter\,章}
    {6pt}
    {\LARGE}
    [{\vspace{5pt}{\color{yhbAccent!55}\titlerule[1pt]}}]
  \titlespacing*{\chapter}{0pt}{14pt}{24pt}
  \titleformat{\section}{\Large\sffamily\bfseries\color{yhbAccent}}{\thesection}{0.7em}{}
  \titlespacing*{\section}{0pt}{18pt plus 3pt minus 2pt}{8pt}
\else
  \titleformat{\section}{\Large\sffamily\bfseries\color{yhbAccent}}{\thesection}{0.7em}{}[{\vspace{3pt}{\color{yhbAccent!45}\titlerule[0.7pt]}}]
  \titlespacing*{\section}{0pt}{18pt plus 3pt minus 2pt}{9pt}
\fi
\titleformat{\subsection}{\large\sffamily\bfseries\color{yhbAccent!85!black}}{\thesubsection}{0.7em}{}
\titleformat{\subsubsection}{\normalsize\bfseries\color{black!85}}{\thesubsubsection}{0.7em}{}
\titlespacing*{\subsection}{0pt}{12pt plus 2pt minus 1pt}{5pt}
\titlespacing*{\subsubsection}{0pt}{9pt plus 2pt minus 1pt}{4pt}
\pretocmd{\section}{\Needspace{6\baselineskip}}{}{}
\pretocmd{\subsection}{\Needspace{5\baselineskip}}{}{}
\pretocmd{\subsubsection}{\Needspace{4\baselineskip}}{}{}

\ifdefined\YHBbooktrue
  \setcounter{secnumdepth}{2}
  \setcounter{tocdepth}{1}
\else
  \setcounter{secnumdepth}{3}
  \setcounter{tocdepth}{2}
\fi
\renewcommand{\contentsname}{目录}
% 目录编号颜色由 titletoc 的 titlecontents 控制，不再重定义 \numberline

% ---------- 目录样式 ----------
\ifdefined\YHBbooktrue
  \titlecontents{chapter}[0pt]
    {\addvspace{6pt}\sffamily\bfseries\color{yhbAccent}}
    {\contentslabel{2.2em}}{}{\titlerule*[0.6em]{.}\contentspage}
  \titlecontents{section}[2.2em]
    {\addvspace{2pt}}
    {\contentslabel{2.2em}}{}{\titlerule*[0.6em]{.}\contentspage}
\else
  \titlecontents{section}[0pt]
    {\addvspace{5pt}\sffamily\bfseries\color{yhbAccent}}
    {\contentslabel{2.2em}}{}{\titlerule*[0.6em]{.}\contentspage}
  \titlecontents{subsection}[2.2em]
    {\addvspace{2pt}}
    {\contentslabel{2.4em}}{}{\titlerule*[0.6em]{.}\contentspage}
\fi

% ---------- 页眉页脚 ----------
\pagestyle{fancy}
\fancyhf{}
\fancyhead[L]{\footnotesize\sffamily\color{textgray}\YHBcourse\ \YHBdoclabel}
\ifdefined\YHBbooktrue
  \fancyhead[R]{\footnotesize\sffamily\color{textgray}\nouppercase{\leftmark}}
  \renewcommand{\chaptermark}[1]{\markboth{#1}{}}
\else
  \fancyhead[R]{\footnotesize\sffamily\color{textgray}\nouppercase{\rightmark}}
  \renewcommand{\sectionmark}[1]{\markright{#1}}
\fi
\fancyfoot[C]{\footnotesize\color{textgray}\thepage}
\renewcommand{\headrulewidth}{0.35pt}
\renewcommand{\headrule}{{\color{yhbAccent!45}\hrule width\headwidth height\headrulewidth}}
\fancypagestyle{plain}{%
  \fancyhf{}%
  \fancyfoot[C]{\footnotesize\color{textgray}\thepage}%
  \renewcommand{\headrulewidth}{0pt}%
}

% ---------- 公式与考点框 ----------
\newcommand{\keyformula}[1]{%
  \begin{tcolorbox}[enhanced,breakable,colback=highlight,colframe=yhbAccent!30!white,boxrule=0.5pt,arc=2.5pt,left=11pt,right=11pt,top=9pt,bottom=9pt,boxsep=0pt,before skip=11pt,after skip=11pt,halign=center,width=\linewidth]
  \adjustbox{max width=\linewidth-2em}{$\displaystyle #1$}
  \end{tcolorbox}}
\newcommand{\corepoint}[1]{%
  \begin{tcolorbox}[enhanced,breakable,colback=highlight,colframe=yhbAccent2,boxrule=0pt,borderline west={2.4pt}{0pt}{yhbAccent2},arc=0pt,left=11pt,right=11pt,top=8pt,bottom=8pt,before skip=11pt,after skip=11pt]
  \textbf{\color{yhbAccent}核心考点}\quad #1
  \end{tcolorbox}}
\newcommand{\symb}[2]{\textbf{#1}：#2\par}

% ---------- 超链接（务必最后加载） ----------
\usepackage[unicode=true,colorlinks=true,linkcolor=yhbAccent,citecolor=yhbAccent,urlcolor=yhbAccent!85!black,pdfborder={0 0 0},bookmarksnumbered=true,pdfstartview=FitH]{hyperref}
\usepackage{hypcap}
\hypersetup{pdftitle={\YHBcourse 课程笔记},pdfauthor={浮萍},pdfsubject={课程学习与复习笔记},pdfcreator={XeLaTeX}}

%   \newcommand{\YHBcourse}{微积分A}      % 课程名（页眉与书签用）
%   \newcommand{\YHBversion}{v1.2}        % 版本号（封面用）
%   \newcommand{\YHBquote}{……}            % 封面题句
%   \definecolor{yhbAccent}{HTML}{24465C} % 主色（深）
%   \definecolor{yhbAccent2}{HTML}{347A80}% 辅色（浅）
%   \YHBbooktrue / \YHBbookfalse          % 是否 book 文档类（有 \chapter）
% 章节文件只写内容，不要重复 \usepackage、不要改字号。
% ============================================================

% ---------- 页面与字体 ----------
\usepackage{geometry}
\geometry{a4paper,left=24mm,right=24mm,top=23mm,bottom=24mm,headheight=15pt,headsep=7mm,footskip=13mm}
\usepackage{microtype}
\usepackage{amsmath,amssymb,amsfonts,mathtools}
\usepackage{esint}
\usepackage{graphicx}
\usepackage{adjustbox}
\usepackage{array,booktabs,multirow,makecell,tabularx,longtable}
\usepackage{xcolor,colortbl}
\usepackage{enumitem}
\usepackage{needspace,etoolbox}
\usepackage{titlesec,titletoc}
\usepackage{fancyhdr}
\usepackage{caption}
\usepackage[most]{tcolorbox}

% ---------- 颜色别名（章节里用的旧名继续有效） ----------
\colorlet{mainblue}{yhbAccent}
\colorlet{secondcyan}{yhbAccent2}
\colorlet{thirdteal}{yhbAccent2}
\colorlet{highlight}{yhbAccent!6!white}
\colorlet{tablehead}{yhbAccent!12!white}
\definecolor{textgray}{HTML}{4A5057}

% ---------- 课程标签 ----------
\providecommand{\YHBdoclabel}{课程笔记}
% ---------- 正文节奏 ----------
\linespread{1.3}
\setlength{\parindent}{0pt}
\setlength{\parskip}{6pt plus 1.5pt minus 1.5pt}
\setlength{\emergencystretch}{2.5em}
\raggedbottom
\widowpenalty=10000
\clubpenalty=10000
\setlength{\tabcolsep}{7pt}
\renewcommand{\arraystretch}{1.3}
\allowdisplaybreaks[2]
\AtBeginDocument{%
  \setlength{\abovedisplayskip}{10pt plus 2pt minus 2pt}%
  \setlength{\belowdisplayskip}{10pt plus 2pt minus 2pt}%
  \setlength{\abovedisplayshortskip}{6pt}%
  \setlength{\belowdisplayshortskip}{6pt}%
}
\captionsetup{font=small,labelfont=bf,skip=6pt}
\renewcommand{\tablename}{表}
\renewcommand{\figurename}{图}

% ---------- 列表 ----------
\setlist{leftmargin=2em,labelsep=0.6em,itemsep=4pt,parsep=2pt,topsep=6pt,partopsep=0pt}
\setlist[itemize,2]{label=\textbullet,leftmargin=1.6em}

% ---------- 标题层级 ----------
\ifdefined\YHBbooktrue
  \titleformat{\chapter}[display]
    {\sffamily\bfseries\color{yhbAccent}}
    {\normalsize\sffamily 第\,\thechapter\,章}
    {6pt}
    {\LARGE}
    [{\vspace{5pt}{\color{yhbAccent!55}\titlerule[1pt]}}]
  \titlespacing*{\chapter}{0pt}{14pt}{24pt}
  \titleformat{\section}{\Large\sffamily\bfseries\color{yhbAccent}}{\thesection}{0.7em}{}
  \titlespacing*{\section}{0pt}{18pt plus 3pt minus 2pt}{8pt}
\else
  \titleformat{\section}{\Large\sffamily\bfseries\color{yhbAccent}}{\thesection}{0.7em}{}[{\vspace{3pt}{\color{yhbAccent!45}\titlerule[0.7pt]}}]
  \titlespacing*{\section}{0pt}{18pt plus 3pt minus 2pt}{9pt}
\fi
\titleformat{\subsection}{\large\sffamily\bfseries\color{yhbAccent!85!black}}{\thesubsection}{0.7em}{}
\titleformat{\subsubsection}{\normalsize\bfseries\color{black!85}}{\thesubsubsection}{0.7em}{}
\titlespacing*{\subsection}{0pt}{12pt plus 2pt minus 1pt}{5pt}
\titlespacing*{\subsubsection}{0pt}{9pt plus 2pt minus 1pt}{4pt}
\pretocmd{\section}{\Needspace{6\baselineskip}}{}{}
\pretocmd{\subsection}{\Needspace{5\baselineskip}}{}{}
\pretocmd{\subsubsection}{\Needspace{4\baselineskip}}{}{}

\ifdefined\YHBbooktrue
  \setcounter{secnumdepth}{2}
  \setcounter{tocdepth}{1}
\else
  \setcounter{secnumdepth}{3}
  \setcounter{tocdepth}{2}
\fi
\renewcommand{\contentsname}{目录}
% 目录编号颜色由 titletoc 的 titlecontents 控制，不再重定义 \numberline

% ---------- 目录样式 ----------
\ifdefined\YHBbooktrue
  \titlecontents{chapter}[0pt]
    {\addvspace{6pt}\sffamily\bfseries\color{yhbAccent}}
    {\contentslabel{2.2em}}{}{\titlerule*[0.6em]{.}\contentspage}
  \titlecontents{section}[2.2em]
    {\addvspace{2pt}}
    {\contentslabel{2.2em}}{}{\titlerule*[0.6em]{.}\contentspage}
\else
  \titlecontents{section}[0pt]
    {\addvspace{5pt}\sffamily\bfseries\color{yhbAccent}}
    {\contentslabel{2.2em}}{}{\titlerule*[0.6em]{.}\contentspage}
  \titlecontents{subsection}[2.2em]
    {\addvspace{2pt}}
    {\contentslabel{2.4em}}{}{\titlerule*[0.6em]{.}\contentspage}
\fi

% ---------- 页眉页脚 ----------
\pagestyle{fancy}
\fancyhf{}
\fancyhead[L]{\footnotesize\sffamily\color{textgray}\YHBcourse\ \YHBdoclabel}
\ifdefined\YHBbooktrue
  \fancyhead[R]{\footnotesize\sffamily\color{textgray}\nouppercase{\leftmark}}
  \renewcommand{\chaptermark}[1]{\markboth{#1}{}}
\else
  \fancyhead[R]{\footnotesize\sffamily\color{textgray}\nouppercase{\rightmark}}
  \renewcommand{\sectionmark}[1]{\markright{#1}}
\fi
\fancyfoot[C]{\footnotesize\color{textgray}\thepage}
\renewcommand{\headrulewidth}{0.35pt}
\renewcommand{\headrule}{{\color{yhbAccent!45}\hrule width\headwidth height\headrulewidth}}
\fancypagestyle{plain}{%
  \fancyhf{}%
  \fancyfoot[C]{\footnotesize\color{textgray}\thepage}%
  \renewcommand{\headrulewidth}{0pt}%
}

% ---------- 公式与考点框 ----------
\newcommand{\keyformula}[1]{%
  \begin{tcolorbox}[enhanced,breakable,colback=highlight,colframe=yhbAccent!30!white,boxrule=0.5pt,arc=2.5pt,left=11pt,right=11pt,top=9pt,bottom=9pt,boxsep=0pt,before skip=11pt,after skip=11pt,halign=center,width=\linewidth]
  \adjustbox{max width=\linewidth-2em}{$\displaystyle #1$}
  \end{tcolorbox}}
\newcommand{\corepoint}[1]{%
  \begin{tcolorbox}[enhanced,breakable,colback=highlight,colframe=yhbAccent2,boxrule=0pt,borderline west={2.4pt}{0pt}{yhbAccent2},arc=0pt,left=11pt,right=11pt,top=8pt,bottom=8pt,before skip=11pt,after skip=11pt]
  \textbf{\color{yhbAccent}核心考点}\quad #1
  \end{tcolorbox}}
\newcommand{\symb}[2]{\textbf{#1}：#2\par}

% ---------- 超链接（务必最后加载） ----------
\usepackage[unicode=true,colorlinks=true,linkcolor=yhbAccent,citecolor=yhbAccent,urlcolor=yhbAccent!85!black,pdfborder={0 0 0},bookmarksnumbered=true,pdfstartview=FitH]{hyperref}
\usepackage{hypcap}
\hypersetup{pdftitle={\YHBcourse 课程笔记},pdfauthor={浮萍},pdfsubject={课程学习与复习笔记},pdfcreator={XeLaTeX}}

\usepackage{tikz}
\usepackage{pgfplots}
\pgfplotsset{compat=1.18}
\usepackage{amsthm}
\newtheorem{theorem}{定理}[section]
\newtheorem{lemma}[theorem]{引理}
\newtheorem{corollary}[theorem]{推论}
\newtheorem{definition}[theorem]{定义}
\newtheorem{property}[theorem]{性质}
\begin{document}
\newcommand{\makeCover}
{
  \begin{titlepage}
    \newgeometry{left=2cm, right=2cm, top=2cm, bottom=2cm}

    \centering
    \vspace*{1cm}

    {\Large \itshape 以矩阵刻画变换，以向量描摹空间}\\[0.5cm]
    \rule{\textwidth}{1pt}\\[0.5cm]

    {\Huge \bfseries \color{mainblue} 代数与几何}\\[8pt]
    {\Huge \bfseries \color{mainblue} 核心公式与考点汇总}\\[12pt]
    {\Large \itshape —— 简洁版复习讲义 ——}\\[1cm]

    \rule{0.6\textwidth}{0.5pt}\\[1cm]

    {\Large \textbf{内容概要：}行列式 | 矩阵 | 线性方程组 | 空间解析几何 | 特征值与二次型}\\[8pt]
    {\Large \textbf{适用场景：}课程复习 · 期末备考 · 公式速查}\\[1.5cm]

    {\Large \textbf{整理者：}浮萍}\\[6pt]
    {\large \textbf{联系邮箱：}\href{mailto:3071974740@qq.com}{3071974740@qq.com}}\\[6pt]
    {\large \textbf{GitHub：}\url{https://github.com/Duckweed-yhb}}\\[6pt]
    {\Large 整理日期：\today}\\[1cm]

    \rule{0.6\textwidth}{0.5pt}\\[1cm]

    {\large \textbf{版本：} v1.0 简洁版}\\[6pt]
    {\small \itshape Copyright © \the\year\ 浮萍 · 禁止商用 · 欢迎交流}\\[0.5cm]

    \rule{\textwidth}{1pt}
    \vfill
    \restoregeometry
  \end{titlepage}
}
\makeCover
\frontmatter
\thispagestyle{fancy}
\tableofcontents
\mainmatter

% ---------------- 第 1 章 行列式 ----------------
% 教材 §1 二阶与三阶行列式 / §2 全排列和对换 / §3 n 阶行列式的定义
%      §4 行列式的性质 / §5 行列式按行（列）展开
\chapter{行列式}
\section{二阶与三阶行列式}

\subsection{二元线性方程组与二阶行列式}

用消元法解二元线性方程组
\[
\begin{cases}
a_{11}x_1+a_{12}x_2=b_1,\\
a_{21}x_1+a_{22}x_2=b_2,
\end{cases}
\]
以$a_{22}$与$a_{12}$分别乘两方程两端再相减消去$x_2$，同理消去$x_1$，得
\[
(a_{11}a_{22}-a_{12}a_{21})x_1=b_1a_{22}-a_{12}b_2,\qquad
(a_{11}a_{22}-a_{12}a_{21})x_2=a_{11}b_2-b_1a_{21}.
\]
当$a_{11}a_{22}-a_{12}a_{21}\neq 0$时，方程组的解由四个系数与两个常数项按“两对相乘再相减”的固定格式组成。把这四个数按原位置排成二行二列的数表，并规定

\begin{definition}
由$2^2$个数排成的数表所确定的表达式
\[
\begin{vmatrix}
a_{11} & a_{12}\\
a_{21} & a_{22}
\end{vmatrix}
=a_{11}a_{22}-a_{12}a_{21}
\]
称为\textbf{二阶行列式}。数$a_{ij}$称为行列式的\textbf{元素}：第一个下标$i$为\textbf{行标}（位于第$i$行），第二个下标$j$为\textbf{列标}（位于第$j$列）；位于第$i$行第$j$列的元素称为$(i,j)$元。
\end{definition}

从$a_{11}$到$a_{22}$的实连线称为\textbf{主对角线}，从$a_{12}$到$a_{21}$的连线称为\textbf{副对角线}。于是二阶行列式等于“主对角线上两元素之积减去副对角线上两元素之积”，这一记忆法称为\textbf{对角线法则}。

记方程组的\textbf{系数行列式}$D$，把$D$的第1列换成常数项得$D_1$，把第2列换成常数项得$D_2$：
\[
D=\begin{vmatrix}a_{11}&a_{12}\\a_{21}&a_{22}\end{vmatrix},\qquad
D_1=\begin{vmatrix}b_1&a_{12}\\b_2&a_{22}\end{vmatrix},\qquad
D_2=\begin{vmatrix}a_{11}&b_1\\a_{21}&b_2\end{vmatrix},
\]
当$D\neq0$时，方程组的解可直接写出：

\keyformula{x_1=\frac{D_1}{D},\qquad x_2=\frac{D_2}{D}}

这就是克拉默法则的二阶情形（一般情形见第2章§4）。注意分子$D_j$是把$D$的\textbf{第$j$列}换成常数项列所得的行列式，列的位置不能换错。

\subsection{三阶行列式}

\begin{definition}
设有$3^2$个数排成3行3列的数表，规定
\[
\begin{vmatrix}
a_{11}&a_{12}&a_{13}\\
a_{21}&a_{22}&a_{23}\\
a_{31}&a_{32}&a_{33}
\end{vmatrix}
=a_{11}a_{22}a_{33}+a_{12}a_{23}a_{31}+a_{13}a_{21}a_{32}
-a_{13}a_{22}a_{31}-a_{11}a_{23}a_{32}-a_{12}a_{21}a_{33},
\]
称为该数表所确定的\textbf{三阶行列式}。
\end{definition}

三阶行列式含$6$项，每项都是不同行、不同列的三个元素之积再冠以正负号，规律遵循\textbf{对角线法则}：平行于主对角线的三条“实线”上三元素之积冠正号，平行于副对角线的三条“虚线”上三元素之积冠负号。

\corepoint{
对角线法则\textbf{只适用于二阶与三阶}行列式，四阶及以上不成立（四阶行列式有$4!=24$项，而“对角线”只能给出$8$项）。三阶以上必须用§3的定义或§5的展开法则计算。
}

与二元方程组类似，对三元线性方程组，若其系数行列式$D\neq0$，则未知数可用三阶行列式表示：$x_j=D_j/D$，其中$D_j$是把$D$的第$j$列换成常数项列所得的三阶行列式（见第2章§4克拉默法则）。

\section{全排列与对换}

\subsection{排列及其逆序数}

把$n$个不同的元素排成一行，称为这$n$个元素的\textbf{全排列}（简称\textbf{排列}）。所有$n$阶排列构成的集合记为$S_n$，其中元素个数即$n$阶排列的个数，记为$|S_n|$：从$n$个元素中任取一个放第一个位置有$n$种取法，再从剩下的$n-1$个中取一个放第二个位置有$n-1$种取法……最后只剩一个元素放在第$n$个位置，于是
\[
|S_n|=n(n-1)\cdots 2\cdot 1=n!.
\]

对$n$个不同元素先规定一个\textbf{标准次序}（$n$个不同自然数通常规定由小到大为标准次序）。在任一排列中，当某一对元素的先后次序与标准次序不同时，就说它构成1个\textbf{逆序}；一个排列中所有逆序的总数叫做这个排列的\textbf{逆序数}，排列$p_1p_2\cdots p_n$的逆序数记为$t(p_1p_2\cdots p_n)$。

\begin{definition}
逆序数为奇数的排列称为\textbf{奇排列}，逆序数为偶数的排列称为\textbf{偶排列}。自然排列$12\cdots n$的逆序数为$0$，是偶排列。
\end{definition}

\textbf{逆序数的计算}：不妨设元素为$1,2,\cdots,n$且由小到大为标准次序。对排列中的每个元素$a_i=p_i$，数出排在它\textbf{前面}且\textbf{比它大}的元素个数$t_i$（即$a_i$的逆序数），则
\[
t(p_1p_2\cdots p_n)=t_1+t_2+\cdots+t_n.
\]
等价地，也可数“排在它后面且比它小”的元素个数，逐元素累加；两种数法结果相同。

\subsection{对换}

在排列中，将任意两个元素对调、其余元素不动，这种作出新排列的操作叫做\textbf{对换}；将相邻两个元素对换，叫做\textbf{相邻对换}。记对换元素$a,b$为$(a\ b)$。

\begin{theorem}
一个排列中的任意两个元素对换，排列改变奇偶性。
\end{theorem}

证明思路：先证相邻对换——设排列为$\cdots ab\cdots$，对换$a,b$后，除$a,b$外其余元素的逆序数不变；当$a<b$时对换后$a$的逆序数增加1而$b$的不变，当$a>b$时$a$的不变而$b$的减少1，故排列奇偶性改变。一般对换可归结为奇数次相邻对换（先经$m$次相邻对换把$a$移到$b$左侧，再经$m+1$次把$b$移到$a$的位置，共$2m+1$次），而每次相邻对换都改变奇偶性，共改变奇数次，故奇偶性改变。

\begin{corollary}
奇排列对换成标准排列的对换次数为奇数，偶排列对换成标准排列的对换次数为偶数。
\end{corollary}

\corepoint{
本节的结论只为§3给出$n$阶行列式的一般定义、并在§4证明行列式性质时使用，属于“非必读”的理论准备。考试中的常考形式是：给定排列求逆序数、判断奇偶性；或由“对换改变奇偶性”推出“$n\ge2$时奇、偶排列各占$n!/2$个”。
}

\section{$n$阶行列式的定义}

先看三阶行列式的结构，以便推广。三阶行列式的6项中，每项都是不同行、不同列的三个元素之积，故除符号外每项可写成$a_{1p_1}a_{2p_2}a_{3p_3}$：第一个下标（行标）排成标准次序$123$，第二个下标（列标）$p_1p_2p_3$是$1,2,3$的某个排列。这样的排列共$3!=6$种，正好对应6项。再看符号：带正号的三项列标排列为$123,231,312$，都是偶排列；带负号的三项列标排列为$132,213,321$，都是奇排列。因此各项所带符号可统一写成$(-1)^{t}$，$t$为列标排列的逆序数，即
\[
\begin{vmatrix}
a_{11}&a_{12}&a_{13}\\
a_{21}&a_{22}&a_{23}\\
a_{31}&a_{32}&a_{33}
\end{vmatrix}
=\sum_{p_1p_2p_3}(-1)^{t(p_1p_2p_3)}a_{1p_1}a_{2p_2}a_{3p_3}.
\]

\begin{definition}
设有$n^2$个数排成$n$行$n$列的数表，作出表中位于不同行不同列的$n$个数的乘积，并冠以符号$(-1)^{t}$，得到形如
\[
(-1)^{t(p_1p_2\cdots p_n)}a_{1p_1}a_{2p_2}\cdots a_{np_n}
\]
的项，其中$p_1p_2\cdots p_n$为自然数$1,2,\cdots,n$的一个排列，$t(p_1p_2\cdots p_n)$为该排列的逆序数。这样的项共有$n!$个，所有这$n!$项的代数和
\[
D=\sum_{p_1p_2\cdots p_n}(-1)^{t(p_1p_2\cdots p_n)}a_{1p_1}a_{2p_2}\cdots a_{np_n}
\]
称为\textbf{$n$阶行列式}，记作
\[
D=\begin{vmatrix}
a_{11}&a_{12}&\cdots&a_{1n}\\
a_{21}&a_{22}&\cdots&a_{2n}\\
\vdots&\vdots&&\vdots\\
a_{n1}&a_{n2}&\cdots&a_{nn}
\end{vmatrix},
\]
简记作$\det(a_{ij})$，数$a_{ij}$称为行列式$D$的$(i,j)$元。
\end{definition}

按此定义作出的二阶、三阶行列式与§1用对角线法则定义的完全一致。当$n=1$时，一阶行列式$|a_{11}|=a_{11}$，注意不要与绝对值记号混淆。

由于行标也可以取遍所有排列，定义还有等价形式（行标为排列、列标为标准次序）：
$D=\sum_{q_1q_2\cdots q_n}(-1)^{t(q_1q_2\cdots q_n)}a_{q_11}a_{q_22}\cdots a_{q_nn}$。

\begin{definition}
主对角线以下（上）的元素全为$0$的行列式称为\textbf{上（下）三角形行列式}；主对角线以下和以上的元素全为$0$的行列式称为\textbf{对角行列式}。
\end{definition}

\begin{theorem}
上三角形、下三角形与对角行列式都等于其主对角线上各元素之积：
\[
\begin{vmatrix}
a_{11}&a_{12}&\cdots&a_{1n}\\
0&a_{22}&\cdots&a_{2n}\\
\vdots&\vdots&\ddots&\vdots\\
0&0&\cdots&a_{nn}
\end{vmatrix}
=\begin{vmatrix}
a_{11}&0&\cdots&0\\
a_{21}&a_{22}&\cdots&0\\
\vdots&\vdots&\ddots&\vdots\\
a_{n1}&a_{n2}&\cdots&a_{nn}
\end{vmatrix}
=a_{11}a_{22}\cdots a_{nn}.
\]
\end{theorem}

\keyformula{\begin{vmatrix}\lambda_1&&\\&\ddots&\\&&\lambda_n\end{vmatrix}=\lambda_1\lambda_2\cdots\lambda_n}

证明思路（以下三角为例）：当$i<j$时$a_{ij}=0$，故$D$中可能不为零的元素$a_{ip_i}$须满足$i\ge p_i$，又$p_1p_2\cdots p_n$是排列，只能有$p_i=i$，即可能不为零的项只有$a_{11}a_{22}\cdots a_{nn}$一项，其符号$(-1)^{t(12\cdots n)}=1$。

\corepoint{
由定义求$n$阶行列式要枚举$n!$项，只有“非零项极少”的特殊行列式（三角形、对角、副对角线型）才直接用定义。一般情形用§4的性质化三角形、或用§5的展开法则降阶。
}

\section{行列式的性质}

记$D^T$为把$D$的行换成同序数的列所得的行列式，称为$D$的\textbf{转置行列式}，即$D^T=\det(a_{ji})$。约定记号：$r_i\leftrightarrow r_j$表示对换$i,j$两行，$r_i(k)$表示第$i$行乘$k$，$r_i+kr_j$表示第$j$行的$k$倍加到第$i$行；列变换相应记作$c_i\leftrightarrow c_j$、$c_i(k)$、$c_i+kc_j$。

\begin{property}
行列式与它的转置行列式相等：$D^T=D$。
\end{property}

因此行列式的行与列地位同等：对行成立的每一条性质，对列也同样成立，反之亦然。以下性质只对行叙述。

\begin{property}
对换行列式的两行（记作$r_i\leftrightarrow r_j$），行列式变号。
\end{property}

\begin{corollary}
若行列式有两行完全相同，则此行列式等于零。
\end{corollary}

（证明：对换这两行，由性质2得$D=-D$，故$D=0$。）

\begin{property}
行列式的某一行中所有元素都乘同一数$k$，等于用数$k$乘此行列式（记作$r_i(k)$），即
\[
\begin{vmatrix}
a_{11}&a_{12}&\cdots&a_{1n}\\
\vdots&\vdots&&\vdots\\
ka_{i1}&ka_{i2}&\cdots&ka_{in}\\
\vdots&\vdots&&\vdots\\
a_{n1}&a_{n2}&\cdots&a_{nn}
\end{vmatrix}
=k\begin{vmatrix}
a_{11}&a_{12}&\cdots&a_{1n}\\
\vdots&\vdots&&\vdots\\
a_{i1}&a_{i2}&\cdots&a_{in}\\
\vdots&\vdots&&\vdots\\
a_{n1}&a_{n2}&\cdots&a_{nn}
\end{vmatrix}.
\]
\end{property}

\begin{corollary}
行列式中某一行所有元素的公因子可以提到行列式记号外面。
\end{corollary}

\begin{corollary}
行列式中若有两行元素成比例，则此行列式等于零。
\end{corollary}

\begin{property}
若行列式的某一行元素都是两数之和，例如第$i$行的元素都是两数之和，则$D$等于下列两个行列式之和：这两个行列式除第$i$行外其余各行与$D$完全相同，而第$i$行分别是这两组数。
\end{property}

\begin{property}
把行列式的某一行各元素乘同一数然后加到另一行对应的元素上去（记作$r_i+kr_j$），行列式不变：
\[
\begin{vmatrix}
\vdots&\vdots&&\vdots\\
a_{i1}+ka_{j1}&a_{i2}+ka_{j2}&\cdots&a_{in}+ka_{jn}\\
\vdots&\vdots&&\vdots\\
a_{j1}&a_{j2}&\cdots&a_{jn}\\
\vdots&\vdots&&\vdots
\end{vmatrix}
=\begin{vmatrix}
\vdots&\vdots&&\vdots\\
a_{i1}&a_{i2}&\cdots&a_{in}\\
\vdots&\vdots&&\vdots\\
a_{j1}&a_{j2}&\cdots&a_{jn}\\
\vdots&\vdots&&\vdots
\end{vmatrix}.
\]
\end{property}

性质2、3、5给出了行列式关于行的三种运算（对换、倍乘、倍加），它们与矩阵的初等行变换一一对应。由性质4还可看出：若$n$阶行列式每个元素都表示成两数之和，则可分解成$2^n$个行列式之和。

\subsection{用性质化三角形计算行列式}

利用上述运算可简化行列式的计算，常用的一条路线是\textbf{化三角形}：

\begin{enumerate}
\item 观察结构：若某一行（列）零元多、或各行（列）元素之和相等，优先对该行（列）下手；
\item 用$r_i+kr_j$把主对角线以下的元素逐个化为$0$（化上三角形），必要时先用$r_i(k)$把主元化为$1$以减少分数；
\item 化成三角形后，行列式等于主对角线上各元素之积，注意累计每次$r_i\leftrightarrow r_j$带来的负号。
\end{enumerate}

若行列式的零元分布在副对角线一侧（“副对角线型”），可先把第1行依次与第$n,n-1,\cdots,2$行对换，共$n-1$次对换，化为上三角形；再仿此处理后$n-1$行，最终共需$\frac{n(n-1)}{2}$次对换，故
\[
\begin{vmatrix}
&&&a_{1n}\\&&a_{2,n-1}&\\&\vdots&&\\a_{n1}&&&
\end{vmatrix}
=(-1)^{\frac{n(n-1)}{2}}a_{1n}a_{2,n-1}\cdots a_{n1}.
\]

\corepoint{
三处最常见的错：
（1）$r_i+kr_j$与$r_j+kr_i$结果不同，前者是“$j$行的$k$倍加到$i$行”，不能套用加法交换律写成$kr_i+r_j$；
（2）把几个运算写在一起（如$r_2-r_1,\ r_4+5r_1$）时，后一次运算作用于前一次的结果，\textbf{次序一般不能颠倒}，颠倒会得出错误结果；
（3）用性质3提出公因子时只对\textbf{一行}提，若$n$阶行列式每行都提一次公因子$k$，外面应是$k^n$。
}

\subsection{分块三角行列式与递推法}

若$A$为$m$阶方阵、$B$为$n$阶方阵，$O$为零矩阵，则
\[
\begin{vmatrix}A&O\\C&B\end{vmatrix}
=\begin{vmatrix}A&C\\O&B\end{vmatrix}
=|A|\,|B|,
\]
其中$|A|$、$|B|$分别表示相应方阵的行列式。（证明思路：对前$m$行作$r_i+\lambda r_j$把$A$化为下三角形，再对后$n$列作$c_j+\lambda c_i$把$B$化为下三角形，即得下三角形行列式，其值为$|A||B|$。）

对含参数或阶数不定的$n$阶行列式，常用\textbf{递推法}：先按某行（列）展开或对特殊结构降阶，导出$D_n$与$D_{n-1}$（或$D_{n-2}$）的递推公式，再验证$n=1$（或$n=2$）时结论成立，最后由递推式求出通项。归纳法的主要步骤——导出递推公式、检验初始情形——不可省略。

\section{行列式按行（列）展开}

一般说来，低阶行列式比高阶行列式容易计算，因此考虑用低阶行列式表示高阶行列式。

在$n$阶行列式中，把$(i,j)$元$a_{ij}$所在的第$i$行和第$j$列划去后，留下来的$n-1$阶行列式叫做$(i,j)$元$a_{ij}$的\textbf{余子式}，记作$M_{ij}$；记
\[
A_{ij}=(-1)^{i+j}M_{ij},
\]
叫做$(i,j)$元$a_{ij}$的\textbf{代数余子式}。$A_{ij}$与$M_{ij}$只差一个符号，符号由$i+j$的奇偶决定（等价地由$(i,j)$位置“正负相间”的棋盘格决定，$(1,1)$元处为正）。

\begin{lemma}
一个$n$阶行列式，如果其中第$i$行所有元素除$(i,j)$元$a_{ij}$外都为零，那么这个行列式等于$a_{ij}$与它的代数余子式的乘积，即
\[
D=a_{ij}A_{ij}.
\]
\end{lemma}

证明思路：先证$i=j=1$的情形（此时$D=a_{11}M_{11}=a_{11}A_{11}$，是分块三角行列式的特例）；一般情形把第$i$行依次与第$i-1,i-2,\cdots,1$行对换、第$j$列依次与第$j-1,\cdots,1$列对换，经$(i-1)+(j-1)$次对换把$a_{ij}$换成$(1,1)$元，符号为$(-1)^{i+j-2}=(-1)^{i+j}$，再用$i=j=1$的结果。

\begin{theorem}
行列式等于它的任一行（列）的各元素与其对应的代数余子式乘积之和，即
\[
D=a_{i1}A_{i1}+a_{i2}A_{i2}+\cdots+a_{in}A_{in}\quad(i=1,2,\cdots,n),
\]
\[
D=a_{1j}A_{1j}+a_{2j}A_{2j}+\cdots+a_{nj}A_{nj}\quad(j=1,2,\cdots,n).
\]
\end{theorem}

这个定理叫做\textbf{行列式按行（列）展开法则}。证明思路：把第$i$行元素拆成$n$个“只有一个非零元”的行列式之和（性质4），再对每个用引理。

\begin{corollary}
行列式某一行（列）的元素与另一行（列）对应元素的代数余子式乘积之和等于零，即
\[
a_{i1}A_{j1}+a_{i2}A_{j2}+\cdots+a_{in}A_{jn}=0\quad(i\neq j),
\]
\[
a_{1i}A_{1j}+a_{2i}A_{2j}+\cdots+a_{ni}A_{nj}=0\quad(i\neq j).
\]
\end{corollary}

综合定理与推论，关于代数余子式有统一的重要性质：
\keyformula{\sum_{k=1}^{n}a_{ik}A_{jk}=\begin{cases}D,&i=j,\\0,&i\neq j,\end{cases}\qquad
\sum_{k=1}^{n}a_{ki}A_{kj}=\begin{cases}D,&i=j,\\0,&i\neq j.\end{cases}}

后一式也常写成$\delta_{ij}D$（$\delta_{ij}$为Kronecker符号）。记忆口诀：“\textbf{本行乘本行余子式得$D$，异行乘异行余子式得$0$}”。这一结论是第2章用伴随矩阵求逆矩阵的基础。

\subsection{按行（列）展开的计算步骤}

\begin{enumerate}
\item 选取零元最多的一行（列）作为展开行（列）；若没有明显的零元，用$r_i+kr_j$先把该行（列）化成只有一个非零元；
\item 按该行（列）展开，把$n$阶降为$n-1$阶，注意$(-1)^{i+j}$的符号；
\item 重复上述过程，直到化为二阶或三阶行列式，用对角线法则算到底。
\end{enumerate}

若需“造零”的那一行（列）中有某个元为$1$，就以它为主元，用$r_i-r_1$等运算把同列其余元素一次消为零；若主元不为$1$，可先用$r_i(k)$把它化为$1$，避免出现分数。

\subsection{范德蒙德行列式}

称形如下式的$n$阶行列式为\textbf{范德蒙德（Vandermonde）行列式}：
\[
D_n=\begin{vmatrix}
1&1&\cdots&1\\
x_1&x_2&\cdots&x_n\\
x_1^2&x_2^2&\cdots&x_n^2\\
\vdots&\vdots&&\vdots\\
x_1^{n-1}&x_2^{n-1}&\cdots&x_n^{n-1}
\end{vmatrix}
=\prod_{1\le j<i\le n}(x_i-x_j).
\]
证明用数学归纳法：从第$n$行开始，后行减去前行的$x_1$倍，再按第1列展开并把每列的公因子$(x_i-x_1)$提出，即化为$n-1$阶范德蒙德行列式。

\corepoint{
（1）范德蒙德行列式要求“列（行）成等比”，即每列是同一个数的$0,1,\cdots,n-1$次幂；若某行缺项或幂次跳跃，不可直接套公式。
（2）公式右端是\textbf{所有}$x_i-x_j$（$i>j$）的乘积，共$\frac{n(n-1)}{2}$个因子；符号规律是大下标减小下标，写反会整体差一个$(-1)^{\frac{n(n-1)}{2}}$。
（3）$D_n=0\iff$有两个$x_i$相等，这一判据常用于“证明某行列式为零”或“求方程$D_n=0$的根”。
}


% ---------------- 第 2 章 矩阵及其运算 ----------------
% 教材 §1 线性方程组和矩阵 / §2 矩阵的运算 / §3 逆矩阵
%      §4 克拉默法则 / §5 矩阵分块法
\chapter{矩阵及其运算}
\section{线性方程组和矩阵}

\subsection{数域}

如果数集$F$包含$0$和$1$，并且$F$中任意两个数的和、差、积、商（除数不为$0$）仍在$F$中，那么称$F$是一个\textbf{数域}。全体有理数之集$\mathbb{Q}$、全体实数之集$\mathbb{R}$、全体复数之集$\mathbb{C}$都是数域，分别称为\textbf{有理数域}、\textbf{实数域}、\textbf{复数域}，这是三个最常用的数域。本书的讨论一般在实数域上进行，但结论对一般数域同样成立。

\subsection{线性方程组与矩阵}

含$m$个方程、$n$个未知量的线性方程组（其通式见第3章§3），全部信息由两张“数表”决定：系数表与常数项表。把系数按原位置排成$m$行$n$列的数表，就得到\textbf{系数矩阵}$A$；把常数项排成一列得到\textbf{常数项列向量}$\beta$。于是方程组可简记为$AX=\beta$。

\begin{definition}
由数域$F$内$m\times n$个数$a_{ij}\;(i=1,\cdots,m;\ j=1,\cdots,n)$排成的$m$行$n$列的数表
\[
A=\begin{pmatrix}
a_{11}&a_{12}&\cdots&a_{1n}\\
a_{21}&a_{22}&\cdots&a_{2n}\\
\vdots&\vdots&&\vdots\\
a_{m1}&a_{m2}&\cdots&a_{mn}
\end{pmatrix}
\]
叫做$F$上的$m$行$n$列\textbf{矩阵}，简称$m\times n$矩阵，记作$A=(a_{ij})_{m\times n}$或$A=(a_{ij})$。数$a_{ij}$称为矩阵的\textbf{元素}（元），$i$为行标、$j$为列标。
\end{definition}

矩阵与行列式是两个不同的概念：行列式是一个\textbf{数}（只对$n$阶方阵才有意义），矩阵是一张\textbf{数表}（行数、列数可以任意），两者记号也不同（行列式用$|\ |$，矩阵用$(\ )$或$[\ ]$）。

\subsection{几类特殊矩阵}

\begin{itemize}
\item \textbf{零矩阵}：所有元素均为$0$的矩阵，记作$O$（$m\times n$零矩阵记作$O_{m\times n}$）；
\item \textbf{方阵}：行数与列数都等于$n$的矩阵，称为$n$阶方阵，$a_{11},a_{22},\cdots,a_{nn}$称为它的\textbf{主对角线元}；
\item \textbf{对角矩阵}：非主对角线元全为$0$的方阵，记作$\Lambda=\mathrm{diag}(\lambda_1,\lambda_2,\cdots,\lambda_n)$；主对角线元都相等的对角矩阵称为\textbf{标量矩阵}；主对角线元都是$1$的对角矩阵称为\textbf{单位矩阵}，记作$E$（或$E_n$）；
\item \textbf{同型矩阵}：行数相同、列数也相同的两个矩阵；同型矩阵对应元素全相等时称两矩阵\textbf{相等}，记作$A=B$；
\item \textbf{行矩阵与列矩阵}：$1\times n$矩阵称行矩阵，$m\times 1$矩阵称列矩阵（列向量）。
\end{itemize}

\section{矩阵的运算}

\subsection{矩阵的加法}

\begin{definition}
设$A=(a_{ij})$，$B=(b_{ij})$都是$m\times n$矩阵，规定$A$与$B$的\textbf{和}为矩阵$(a_{ij}+b_{ij})_{m\times n}$，记为$A+B$，即\textbf{对应元素相加}。只有\textbf{同型}矩阵才能相加。
\end{definition}

设$A=(a_{ij})_{m\times n}$，称$(-a_{ij})_{m\times n}$为$A$的\textbf{负矩阵}，记作$-A$。矩阵加法满足（$A,B,C$同为$m\times n$矩阵）：$A+B=B+A$，$(A+B)+C=A+(B+C)$，$A+O=A$，$A+(-A)=O$。规定$A-B=A+(-B)$。

\subsection{数与矩阵相乘}

\begin{definition}
数$k$与矩阵$A=(a_{ij})_{m\times n}$的乘积，规定为矩阵$(ka_{ij})_{m\times n}$，记作$kA$或$Ak$，即数$k$乘矩阵的每一个元素。
\end{definition}

数乘满足（$k,l$为数，$A,B$为同型矩阵）：$k(A+B)=kA+kB$，$(k+l)A=kA+lA$，$(kl)A=k(lA)$，$1\cdot A=A$，$0\cdot A=O$，$kO=O$。

\subsection{矩阵与矩阵相乘}

\begin{definition}
设$A=(a_{ij})_{m\times s}$，$B=(b_{ij})_{s\times n}$，规定$A$与$B$的\textbf{乘积}是一个$m\times n$矩阵$C=(c_{ij})$，其中
\[
c_{ij}=a_{i1}b_{1j}+a_{i2}b_{2j}+\cdots+a_{is}b_{sj}=\sum_{k=1}^{s}a_{ik}b_{kj},
\]
记作$C=AB$。即$c_{ij}$是$A$的第$i$行与$B$的第$j$列的对应元素乘积之和（“行乘列”）。
\end{definition}

三点必须记牢：(1)只有当左矩阵$A$的\textbf{列数}等于右矩阵$B$的\textbf{行数}时$AB$才有意义；(2)$AB$的行数等于$A$的行数、列数等于$B$的列数；(3)$AB$的$(i,j)$元等于$A$的第$i$行与$B$的第$j$列的对应元乘积之和，即“行乘列”。

矩阵乘法满足（假设运算都有意义）：\\
(1)$(AB)C=A(BC)$（结合律）；(2)$A(B+C)=AB+AC$，$(B+C)A=BA+CA$（分配律）；(3)$k(AB)=(kA)B=A(kB)$（$k$为数）；(4)$E_mA_{m\times n}=A_{m\times n}$，$A_{m\times n}E_n=A_{m\times n}$（单位矩阵是乘法的“幺元”）。

矩阵乘法有三条与数的乘法不同的特殊性：\\
(1)\textbf{不满足交换律}——虽然$AB$有意义但$BA$未必有意义；即使$AB,BA$都有意义也未必相等。一般$AB\neq BA$；若$AB=BA$，称$A$与$B$可交换。\\
(2)\textbf{不满足消去律}——一般由$AB=AC$（或$BA=CA$）不能断定$B=C$。\\
(3)\textbf{有零因子}——存在$A\neq O,B\neq O$使$AB=O$，故不能由$AB=O$断定$A=O$或$B=O$。

\subsection{方阵的幂与方阵的多项式}

设$A$为$n$阶方阵，规定$A^1=A$，$A^{k+1}=A^kA$（$k$为正整数），称为$A$的\textbf{幂}。有$A^kA^l=A^{k+l}$，$(A^k)^l=A^{kl}$（$k,l$为非负整数）。

\corepoint{
作为乘法不满足交换律的直接后果，一般$(AB)^k\neq A^kB^k$，且$A^k=B^k$推不出$A=B$。只有$A,B$可交换时才有二项式展开$(A+B)^2=A^2+2AB+B^2$，一般$(A+B)^2=A^2+AB+BA+B^2$。
}

利用方阵的幂可定义方阵的多项式：$f(A)=a_mA^m+a_{m-1}A^{m-1}+\cdots+a_1A+a_0E$，注意常数项要乘单位矩阵$E$。

\subsection{矩阵的转置}

\begin{definition}
把$m\times n$矩阵$A$的行换成同序数的列（第1行变成第1列，第2行变成第2列……），所得的$n\times m$矩阵称为$A$的\textbf{转置矩阵}，记作$A'$（也记作$A^{T}$）。
\end{definition}

转置满足（假设运算都可行）：
\begin{enumerate}
\item $(A')'=A$；
\item $(A+B)'=A'+B'$；
\item $(kA)'=kA'$（$k$为数）；
\item $(AB)'=B'A'$（注意顺序颠倒）；
\item $|A'|=|A|$（$A$为方阵时）。
\end{enumerate}

设$A=(a_{ij})$为$n$阶方阵。若$A'=A$，则称$A$为\textbf{对称矩阵}；若$A'=-A$，则称$A$为\textbf{反称矩阵}。对称矩阵的特点是元素以主对角线为对称轴对应相等（$a_{ij}=a_{ji}$）；反称矩阵的主对角线元全为$0$。

两个$n$阶对称矩阵之和仍是对称矩阵；若$A,B$是对称矩阵且可交换（$AB=BA$），则$AB$也是对称矩阵。

\subsection{方阵的行列式}

\begin{definition}
由$n$阶方阵$A$的元素所构成的行列式（各元素的位置不变）叫做方阵$A$的\textbf{行列式}，记作$|A|$或$\det A$。
\end{definition}

\begin{theorem}
设$A,B$为$n$阶方阵，$k$为数，则
\begin{enumerate}
\item $|kA|=k^n|A|$；
\item $|AB|=|A|\,|B|$，称此式为\textbf{行列式的乘法公式}；
\item $|A'|=|A|$。
\end{enumerate}
\end{theorem}

\keyformula{|A'|=|A|,\qquad |kA|=k^n|A|,\qquad |AB|=|A|\,|B|}

由乘法公式可推出$|AB|=|BA|$。注意$|A+B|\neq|A|+|B|$（一般不等），这是最容易与数的运算混淆的一条。

\subsection{共轭矩阵}

设$A=(a_{ij})$为复矩阵，$\bar a_{ij}$为$a_{ij}$的共轭复数，记$\bar A=(\bar a_{ij})$，称为$A$的\textbf{共轭矩阵}。共轭矩阵满足$\overline{(A+B)}=\bar A+\bar B$，$\overline{(kA)}=\bar k\bar A$，$\overline{(AB)}=\bar A\,\bar B$，$\overline{(A')}=(\bar A)'$。记$A^{H}=(\bar A)'$，称为$A$的\textbf{共轭转置矩阵}；若$A^{H}=A$，则称$A$为Hermite矩阵。

\section{逆矩阵}

\subsection{逆矩阵的定义与性质}

\begin{definition}
设$A$为$n$阶方阵，若存在$n$阶方阵$B$，使
\[
AB=BA=E,
\]
则称$A$是\textbf{可逆矩阵}（或$A$可逆），称$B$是$A$的\textbf{逆矩阵}。
\end{definition}

若$A$可逆，则其逆矩阵是唯一的：设$B_1,B_2$都满足定义，则$B_1=B_1E=B_1(AB_2)=(B_1A)B_2=EB_2=B_2$。因此把$A$的逆矩阵记作$A^{-1}$。对角矩阵只要主对角线元全不为零就可逆，且其逆仍为对角矩阵，$\mathrm{diag}(\lambda_1,\cdots,\lambda_n)^{-1}=\mathrm{diag}(\lambda_1^{-1},\cdots,\lambda_n^{-1})$。

可逆矩阵的性质：
\begin{enumerate}
\item 若$A$可逆，则$A^{-1}$可逆，且$(A^{-1})^{-1}=A$；
\item 若$A$可逆，数$k\neq0$，则$kA$可逆，且$(kA)^{-1}=\dfrac{1}{k}A^{-1}$；
\item 若$A,B$为同阶可逆方阵，则$AB$可逆，且$(AB)^{-1}=B^{-1}A^{-1}$（推广：$(A_1A_2\cdots A_s)^{-1}=A_s^{-1}\cdots A_2^{-1}A_1^{-1}$）；
\item 若$A$可逆，则$A'$可逆，且$(A')^{-1}=(A^{-1})'$；
\item 若$A$可逆，则$|A^{-1}|=|A|^{-1}$（由$|A||A^{-1}|=|AA^{-1}|=|E|=1$）。
\end{enumerate}

\corepoint{
$(AB)^{-1}=B^{-1}A^{-1}$必须\textbf{反序}，写成$A^{-1}B^{-1}$是错的（$A^{-1}B^{-1}$未必等于$(AB)^{-1}$，除非$A,B$可交换）。同理$(A+B)^{-1}\neq A^{-1}+B^{-1}$，一般$A+B$甚至未必可逆。
}

\subsection{伴随矩阵与可逆的充要条件}

\begin{definition}
设$A=(a_{ij})$为$n$阶方阵，$A_{ij}$是$|A|$中$(i,j)$元$a_{ij}$的代数余子式，则称$n$阶方阵
\[
A^{*}=\begin{pmatrix}
A_{11}&A_{21}&\cdots&A_{n1}\\
A_{12}&A_{22}&\cdots&A_{n2}\\
\vdots&\vdots&&\vdots\\
A_{1n}&A_{2n}&\cdots&A_{nn}
\end{pmatrix}
\]
为$A$的\textbf{伴随矩阵}（注意：$A^*$的$(i,j)$元是$A_{ji}$，即代数余子式要\textbf{转置}排列）。
\end{definition}

\begin{lemma}
设$A$为$n$阶方阵，则
\[
AA^{*}=A^{*}A=|A|E.
\]
\end{lemma}

证明思路：$AA^*$的$(i,j)$元为$\sum_{k=1}^{n}a_{ik}A_{jk}$，由第1章§5“本行乘本行代数余子式得$|A|$、异行乘异行代数余子式得$0$”，得$(i,j)$元为$|A|\delta_{ij}$。

\begin{theorem}
$n$阶方阵$A$可逆的充要条件是$|A|\neq0$；当$|A|\neq0$时
\[
A^{-1}=\frac{1}{|A|}A^{*}.
\]
\end{theorem}

当$|A|=0$时称$A$为\textbf{奇异矩阵}，当$|A|\neq0$时称$A$为\textbf{非奇异矩阵}。于是：$A$可逆$\iff A$非奇异$\iff|A|\neq0$。（另一等价说法是$R(A)=n$，见第3章§2。）

\keyformula{A^{-1}=\frac{1}{|A|}A^{*},\qquad AA^{*}=A^{*}A=|A|E}

若$A$可逆，规定$A^{-k}=(A^{-1})^{k}$（$k=1,2,\cdots$），此时$A^{k}A^{l}=A^{k+l}$对一切整数$k,l$成立。

\subsection{用伴随矩阵法求逆矩阵}

\begin{enumerate}
\item 计算$|A|$，若$|A|=0$则$A$不可逆，停止；
\item 逐个求出全部$n^2$个代数余子式$A_{ij}$（注意符号$(-1)^{i+j}$要算准）；
\item 把$A_{ij}$按\textbf{转置}位置排成$A^{*}$，即$A^*$的第$i$行第$j$列放$A_{ji}$；
\item 用$A^{-1}=\dfrac{1}{|A|}A^{*}$写出结果。
\end{enumerate}

对二阶方阵有便捷口诀（“主对调、副变号、除以行列式”）：
\keyformula{\begin{pmatrix}a&b\\c&d\end{pmatrix}^{-1}=\frac{1}{ad-bc}\begin{pmatrix}d&-b\\-c&a\end{pmatrix}\quad(ad-bc\neq0)}

\corepoint{
伴随矩阵法只适合二阶、三阶（$n\ge4$时要求$n^2$个余子式，计算量爆炸）。求具体数字矩阵的逆，用第3章§1的初等变换法：$(A\mid E)\xrightarrow{\text{行变换}}(E\mid A^{-1})$。
}

\subsection{用逆矩阵解矩阵方程}

设$A$为$n$阶可逆方阵：若$AX=B$则$X=A^{-1}B$（左乘$A^{-1}$）；若$XA=B$则$X=BA^{-1}$（右乘$A^{-1}$）；若$AXB=C$则$X=A^{-1}CB^{-1}$。
关键是$A^{-1}$要乘在$A$原来所在的一侧，因为矩阵乘法不可交换。若$A$不可逆，则转化为解线性方程组（第3章§3）。

\section{克拉默法则}

\begin{theorem}[克拉默法则]
设$n$个未知数、$n$个方程的线性方程组$AX=\beta$的系数行列式$D=|A|\neq0$，则方程组有唯一解
\[
x_1=\frac{D_1}{D},\quad x_2=\frac{D_2}{D},\quad\cdots,\quad x_n=\frac{D_n}{D},
\]
其中$D_j$是把$D$中第$j$列换成常数项$b_1,b_2,\cdots,b_n$所得到的$n$阶行列式。
\end{theorem}

\keyformula{x_j=\frac{D_j}{D}\quad(j=1,2,\cdots,n),\qquad D\neq0}

证明思路：把方程组写成$AX=\beta$，由$D=|A|\neq0$知$A$可逆，$X=A^{-1}\beta=\dfrac{1}{D}A^{*}\beta$，再逐分量算出分子即$D_j$（$A^*$的第$j$行与$\beta$的内积）。

\begin{corollary}[齐次方程组]
$n$个未知数、$n$个方程的齐次方程组$AX=0$，其系数行列式$D\neq0$时\textbf{只有零解}；反之，它\textbf{有非零解（即有无穷多个解）的充要条件是}$D=0$。
\end{corollary}

\keyformula{D\neq0\ \Longrightarrow\ \text{齐次方程组只有零解};\qquad
\text{齐次方程组有非零解}\iff D=0}

\subsection{使用要点与限制}

\begin{enumerate}
\item 克拉默法则\textbf{只适用于“方程个数＝未知数个数”且$D\neq0$}的方程组；这两个条件缺一不可；
\item 计算量大（$n$个未知数要算$n+1$个$n$阶行列式），实际求解应改用第3章的初等变换法；
\item 理论价值大于计算价值：它给出了“解是常数项与系数的显式表达式”这一结论，并直接推出“齐次方程组有非零解$\iff$系数行列式为零”，是第4、5章讨论线性相关性、特征值的出发点；
\item $D=0$时克拉默法则失效（无解或有无穷多解），此时必须用$R(A)$与$R(A\mid\beta)$的关系判定（第3章§3）。
\end{enumerate}

\corepoint{
常见错误：（1）把常数项替换到\textbf{行}上（应替换\textbf{列}）；（2）方程组中某个未知数在某个方程里“缺项”时，没有补$0$就直接写系数行列式；（3）$D=0$仍套公式得$x_j=0/0$；（4）用克拉默法则讨论“方程组有唯一解$\iff D\neq0$”——仅在$m=n$时成立，方程个数与未知数个数不等时必须用秩判定。
}

\section{矩阵分块法}

\subsection{分块矩阵的概念}

用若干条横线与纵线把矩阵$A$分成若干个小块，以子块为元素的矩阵称为\textbf{分块矩阵}。例如把$3\times4$矩阵按前两行、前两列切成四块，可记作
\[
A=\begin{pmatrix}A_{11}&A_{12}\\A_{21}&A_{22}\end{pmatrix},
\]
其中$A_{11}$是$2\times2$子块，$A_{12}$是$2\times2$子块，$A_{21},A_{22}$都是$1\times2$子块。分块方式不唯一，可依需要（按行分块、按列分块、按相同规格分块）选取，原则是\textbf{便于揭示结构}。

把$m\times n$矩阵$A$的每一行作为一个子块，得\textbf{按行分块}$A=(\alpha_1',\alpha_2',\cdots,\alpha_m')'$，$\alpha_i'$即$A$的第$i$行（称$A$的$m$个\textbf{行向量}）；把每一列作为一个子块，得\textbf{按列分块}$A=(\beta_1,\beta_2,\cdots,\beta_n)$，$\beta_j$为$A$的$n$个\textbf{列向量}。

\subsection{分块矩阵的运算}

分块矩阵的加法、数乘、乘法、转置与普通矩阵形式相同，只要\textbf{子块规格匹配}即可：

\begin{enumerate}
\item \textbf{加法}：$A,B$同型且采用\textbf{相同的分块法}，则$A+B$的对应子块相加；
\item \textbf{数乘}：$kA$的每个子块乘$k$；
\item \textbf{乘法}：设$A$为$m\times s$、$B$为$s\times n$矩阵，$A$的列分块法与$B$的行分块法一致（列组数等于行组数，对应子块规格相同），则按“子块相乘再相加”的规则得$C_{ij}=\sum_{k=1}^{t}A_{ik}B_{kj}$，即
\[
\begin{pmatrix}A_{11}&\cdots&A_{1t}\\ \vdots&&\vdots\\ A_{r1}&\cdots&A_{rt}\end{pmatrix}
\begin{pmatrix}B_{11}&\cdots&B_{1q}\\ \vdots&&\vdots\\ B_{t1}&\cdots&B_{tq}\end{pmatrix}
=\begin{pmatrix}C_{11}&\cdots&C_{1q}\\ \vdots&&\vdots\\ C_{r1}&\cdots&C_{rq}\end{pmatrix}.
\]
注意子块乘法\textbf{次序不能颠倒}（$A_{ik}B_{kj}\neq B_{kj}A_{ik}$）；
\item \textbf{转置}：分块矩阵转置时，不仅子块的位置转置，\textbf{每个子块本身也要转置}：
\[
\begin{pmatrix}A_{11}&A_{12}\\A_{21}&A_{22}\end{pmatrix}'=\begin{pmatrix}A_{11}'&A_{21}'\\A_{12}'&A_{22}'\end{pmatrix};
\]
\item \textbf{分块矩阵的幂}：对分块对角矩阵，幂等于各对角块分别求幂。
\end{enumerate}

\subsection{分块对角矩阵}

形如
\[
\Lambda=\begin{pmatrix}
A_1&&&\\&A_2&&\\&&\ddots&\\&&&A_s\end{pmatrix}
\]
（未写出处为零块，$A_i$为$n_i$阶方阵）的矩阵称为\textbf{分块对角矩阵}（准对角矩阵）。它有下列简便性质：
\begin{enumerate}
\item $|\Lambda|=|A_1|\,|A_2|\cdots|A_s|$；
\item $\Lambda$可逆$\iff$每个$A_i$都可逆，且$\Lambda^{-1}=\mathrm{diag}(A_1^{-1},A_2^{-1},\cdots,A_s^{-1})$；
\item $\Lambda^{k}=\mathrm{diag}(A_1^{k},A_2^{k},\cdots,A_s^{k})$；
\item 两个同规格分块对角矩阵的乘积仍为分块对角矩阵，对角块逐块相乘。
\end{enumerate}

\keyformula{\begin{vmatrix}A_1&&\\&\ddots&\\&&A_s\end{vmatrix}=|A_1|\,|A_2|\cdots|A_s|,\qquad \Lambda^{k}=\mathrm{diag}(A_1^{k},\cdots,A_s^{k})}

\subsection{分块三角矩阵}

形如$\begin{pmatrix}A&C\\O&B\end{pmatrix}$（$A$为$m$阶、$B$为$n$阶方阵）的矩阵称为\textbf{分块上三角矩阵}，$\begin{pmatrix}A&O\\C&B\end{pmatrix}$称为\textbf{分块下三角矩阵}。它们与三角形行列式一样“只取对角块”：
\[
\begin{vmatrix}A&C\\O&B\end{vmatrix}=\begin{vmatrix}A&O\\C&B\end{vmatrix}=|A|\,|B|.
\]
若$A,B$都可逆，则
\[
\begin{pmatrix}A&C\\O&B\end{pmatrix}^{-1}=\begin{pmatrix}A^{-1}&-A^{-1}CB^{-1}\\O&B^{-1}\end{pmatrix},
\]
下三角情形把$-A^{-1}CB^{-1}$换成$-B^{-1}CA^{-1}$放在左下方即可。验证方法是直接相乘看是否得单位矩阵；注意非对角块带负号，且$A^{-1}CB^{-1}$中三者的次序不能换。

\subsection{分块乘法的典型用途}

\begin{enumerate}
\item \textbf{把矩阵乘法写成“列的组合”}：$AB=A(\beta_1,\cdots,\beta_n)=(A\beta_1,\cdots,A\beta_n)$，即$AB$的第$j$列是$A$乘$B$的第$j$列；同理$AB$的第$i$行是$A$的第$i$行乘$B$。这是第3、4章讨论线性组合与秩的基本工具。
\item \textbf{把线性方程组写成列向量的线性组合}：$AX=x_1\beta_1+x_2\beta_2+\cdots+x_n\beta_n$，于是“$AX=\beta$有解”等价于“$\beta$可由$A$的列向量组线性表示”（第4章）。
\item \textbf{简化高阶运算}：按“同结构块”分块后，高阶矩阵的乘法、行列式、逆矩阵都可化为低阶问题（对角块、三角块尤其方便）。
\end{enumerate}

\corepoint{
分块运算的三个易错点：（1）分块乘法要求左矩阵的\textbf{列分组}与右矩阵的\textbf{行分组}完全一致，规格不匹配时不能乘；（2）转置时只转置子块位置、忘记子块内部也要转置；（3）分块三角矩阵求逆时负号与乘法次序$A^{-1}CB^{-1}$写错。矩阵的秩及秩的性质见第3章§2。
}


% ---------------- 第 3 章 矩阵的初等变换与线性方程组 ----------------
% 教材 §1 矩阵的初等变换 / §2 矩阵的秩 / §3 线性方程组的解
\chapter{矩阵的初等变换与线性方程组}
\section{矩阵的初等变换}

\subsection{初等变换的概念}

\begin{definition}
下面三种变换称为矩阵的\textbf{初等行变换}：
\begin{enumerate}
\item 对换矩阵的两行（对换$i,j$两行，记作$r_i\leftrightarrow r_j$）；
\item 以数$k\neq0$乘矩阵的某一行（第$i$行乘$k$，记作$r_i\times k$）；
\item 把矩阵的某一行乘数$k$加到另一行上去（第$j$行的$k$倍加到第$i$行，记作$r_i+kr_j$）。
\end{enumerate}
把上述定义中的“行”换成“列”，即得\textbf{初等列变换}的定义（记作$c_i\leftrightarrow c_j$、$c_i\times k$、$c_i+kc_j$）。初等行变换与初等列变换统称\textbf{初等变换}。
\end{definition}

三种初等变换都是可逆的，且其逆变换是同一类型的初等变换：
\[
r_i\leftrightarrow r_j\ \text{的逆变换是其本身};\qquad
r_i\times k\ \text{的逆变换是}\ r_i\times\frac{1}{k};\qquad
r_i+kr_j\ \text{的逆变换是}\ r_i+(-k)r_j.
\]

\begin{definition}
如果矩阵$A$经有限次初等变换变成矩阵$B$，就称矩阵$A$与$B$\textbf{等价}，记作$A\cong B$。
\end{definition}

等价是矩阵之间的一种关系，具有三条性质：反身性（$A\cong A$）、对称性（$A\cong B\Rightarrow B\cong A$）、传递性（$A\cong B,B\cong C\Rightarrow A\cong C$）。等价关系把全体$m\times n$矩阵分成若干等价类，同一等价类的矩阵“形状不同而本质相同”（秩相同，见§2）。

\subsection{行阶梯形矩阵与行最简形矩阵}

\begin{definition}
称矩阵$A$为\textbf{行阶梯形矩阵}，如果它满足：
\begin{enumerate}
\item 若有零行（元素全为零的行），则零行全部位于非零行的下方；
\item 各非零行左起第一个非零元的列标，由上至下严格递增。
\end{enumerate}
\end{definition}

行阶梯形矩阵的“台阶”逐行右移、逐行下降，形如阶梯。若再要求每个非零行的首非零元都是$1$、且这些$1$是它们所在列的唯一非零元，则称为\textbf{行最简形矩阵}。

\begin{theorem}
任何矩阵$A$总可以只经初等行变换化为行阶梯形矩阵，并可进一步化为行最简形矩阵。
\end{theorem}

\textbf{化行阶梯形的步骤}（自上而下，逐列推进）：
\begin{enumerate}
\item 在剩下的子矩阵中选一个非零元作为主元（一般选绝对值小、易计算者，或先用$r_i\leftrightarrow r_j$把非零元换到当前行）；
\item 用$r_j+kr_i$把主元\textbf{下方}同列的元素全部化为$0$（若主元不为$1$，可先用$r_i\times k$化为$1$以避免分数）；
\item 对下方的子矩阵重复上述过程，直至剩下的行全为零行。
\end{enumerate}

\textbf{化行最简形}：在行阶梯形的基础上，\textbf{自下而上}用$r_i+kr_j$把每个首非零元\textbf{上方}的元素也化为$0$，再用$r_i\times k$把首非零元化为$1$。

\corepoint{
（1）初等\textbf{行}变换与初等\textbf{列}变换一般不能混用：解线性方程组只允许行变换（列变换会改变未知量的对应关系）；求秩、求标准形时行、列变换都可用。
（2）化行阶梯形时，主元所在列下方必须全为$0$，否则不是阶梯形；主元位置选择不当会使计算量剧增。
（3）行最简形的首非零元是$1$且其所在列“上下都干净”，这是读出方程组的解、求逆矩阵的关键形态。
}

\subsection{初等矩阵}

\begin{definition}
由单位矩阵$E$经过\textbf{一次}初等变换得到的方阵，称为\textbf{初等矩阵}。初等矩阵共有三类：
\begin{enumerate}
\item 把单位矩阵的第$i,j$两行（列）对换，记作$E(i,j)$；
\item 以数$k\neq0$乘单位矩阵的第$i$行（列），记作$E(i(k))$；
\item 以数$k$乘单位矩阵的第$j$行加到第$i$行（或以$k$乘第$i$列加到第$j$列），记作$E(i,j(k))$。
\end{enumerate}
\end{definition}

\begin{theorem}
设$A$是一个$m\times n$矩阵，则
\begin{enumerate}
\item 对$A$施行一次初等\textbf{行}变换，相当于在$A$的\textbf{左边}乘一个相应的$m$阶初等矩阵；
\item 对$A$施行一次初等\textbf{列}变换，相当于在$A$的\textbf{右边}乘一个相应的$n$阶初等矩阵。
\end{enumerate}
\end{theorem}

初等矩阵的性质：
\begin{enumerate}
\item 初等矩阵都是可逆矩阵，且逆矩阵仍是同类型的初等矩阵：
\[
E(i,j)^{-1}=E(i,j),\qquad E(i(k))^{-1}=E\left(i\left(\tfrac{1}{k}\right)\right),\qquad E(i,j(k))^{-1}=E(i,j(-k));
\]
\item 初等矩阵的转置仍是初等矩阵：$E(i,j)'=E(i,j)$，$E(i(k))'=E(i(k))$，$E(i,j(k))'=E(j,i(k))$；
\item $|E(i,j)|=-1$，$|E(i(k))|=k$，$|E(i,j(k))|=1$，故初等矩阵的行列式都不为$0$。
\end{enumerate}

\begin{theorem}
方阵$A$可逆的充要条件是：存在有限个初等矩阵$P_1,P_2,\cdots,P_k$，使
\[
A=P_1P_2\cdots P_k.
\]
\end{theorem}

由此得\textbf{初等变换法求逆矩阵}：设$A$为$n$阶可逆矩阵，由$P_1P_2\cdots P_sA=E$得$P_1P_2\cdots P_sE=A^{-1}$。两式表明：若一系列初等行变换把$A$化成$E$，则同样的行变换把$E$化成$A^{-1}$。于是对分块矩阵$(A\mid E)$施行初等行变换，当$A$化成$E$时，右边的$E$就化成了$A^{-1}$：

\keyformula{(A\mid E)\xrightarrow{\ \text{初等行变换}\ }(E\mid A^{-1})}

类似地，若$A$可逆，则
\[
(A\mid B)\xrightarrow{\text{行变换}}(E\mid A^{-1}B),\qquad
\begin{pmatrix}A\\B\end{pmatrix}\xrightarrow{\text{列变换}}\begin{pmatrix}E\\BA^{-1}\end{pmatrix}.
\]
这两条一次性解出了矩阵方程$AX=B$与$XA=B$，不必先求$A^{-1}$。

\corepoint{
$(A\mid E)\to(E\mid A^{-1})$过程中\textbf{只能做行变换}，一旦用了列变换结果必错。若中途出现全零行，说明$A$不可逆（$|A|=0$），应立即停止并判定“不可逆”。
}

\section{矩阵的秩}

\subsection{秩的概念}

在$m\times n$矩阵$A$中，任取$k$行与$k$列（$k\le\min\{m,n\}$），位于这些行列交叉处的$k^2$个元素按原位置构成的$k$阶行列式，称为$A$的一个\textbf{$k$阶子式}。$A$共有$C_m^kC_n^k$个$k$阶子式。

\begin{definition}
矩阵$A$中非零子式的最高阶数称为矩阵$A$的\textbf{秩}，记作$R(A)$（或$r(A)$、$\mathrm{rank}(A)$）。
\end{definition}

规定零矩阵的秩为$0$。由定义直接得：
\begin{enumerate}
\item $0\le R(A)\le\min\{m,n\}$；
\item $R(A)=0\iff A=O$；
\item $R(A)\ge r\iff A$中至少有一个$r$阶子式不为零；
\item $R(A)<r\iff A$中所有$r$阶子式全为零；
\item $R(A)=r\iff A$有$r$阶非零子式，而所有$r+1$阶子式（若存在）全为零。
\end{enumerate}

\begin{theorem}
矩阵经初等变换后，其秩不变。
\end{theorem}

证明思路：只需证一次初等变换不改变秩。对“倍乘”“倍加”变换，非零子式的对应关系可直接搬运；对“对换”变换，子式至多变号。反过来，因初等变换可逆，秩也不会增大，故秩不变。

\begin{corollary}
$R(A)=r$的充要条件是$A$经初等变换可化为
\[
\begin{pmatrix}E_r&O\\O&O\end{pmatrix}_{m\times n},
\]
称之为$A$的\textbf{等价标准形}。即存在$m$阶可逆矩阵$P$与$n$阶可逆矩阵$Q$，使$PAQ=\begin{pmatrix}E_r&O\\O&O\end{pmatrix}$。
\end{corollary}

\subsection{求矩阵的秩}

由定义求秩要检查大量子式，实际计算用初等变换：

\begin{enumerate}
\item 用初等行变换（也可兼用列变换）把$A$化为\textbf{行阶梯形矩阵}；
\item 行阶梯形矩阵中\textbf{非零行的行数}就是$R(A)$。
\end{enumerate}

道理是：行阶梯形矩阵的非零行数等于其非零子式的最高阶数，而初等变换不改变秩。

\corepoint{
（1）$R(A)$是一个\textbf{数}（非负整数），不是矩阵；写“$R(A)=0$”表示零矩阵，不是“$A=0$”之外的意思。
（2）$R(A)=n$（$n$为阶数）$\iff|A|\neq0\iff A$可逆$\iff A\cong E_n$，这四条等价说法必须能互相转换。
（3）求秩时若用了列变换，最后读“非零行数”仍然正确（秩不变），但若目的是解方程组，列变换会打乱未知量次序。
}

\subsection{秩的性质与不等式}

设$A$为$m\times n$矩阵，$B$为$n\times p$矩阵，$P,Q$为可逆矩阵，则
\begin{enumerate}
\item $R(A')=R(A)$（转置不改变秩）；
\item $R(PA)=R(AQ)=R(PAQ)=R(A)$（乘可逆矩阵不改变秩）；
\item $R(A+B)\le R(A)+R(B)$（$A,B$同型）；
\item $R(AB)\le\min\{R(A),R(B)\}$；
\item $R(AB)\ge R(A)+R(B)-n$（Sylvester不等式）；
\item 特别地，若$AB=O$，则$R(A)+R(B)\le n$；
\item $\max\{R(A),R(B)\}\le R(A,B)\le R(A)+R(B)$，其中$(A,B)$是把$A,B$按列拼成的分块矩阵；
\item $R\begin{pmatrix}A&O\\O&B\end{pmatrix}=R(A)+R(B)$；
\item 若$A$为$n$阶方阵且$A^2=A$，则$R(A)+R(A-E)=n$。
\end{enumerate}

\keyformula{R(AB)\le\min\{R(A),R(B)\},\qquad R(AB)\ge R(A)+R(B)-n,\qquad AB=O\Rightarrow R(A)+R(B)\le n}

\corepoint{
$R(A+B)\le R(A)+R(B)$一般\textbf{不能取等}，$R(AB)\le\min\{R(A),R(B)\}$同样不能取等；而$AB=O$时$R(A)+R(B)\le n$是“取等”的常见情形（当且仅当$A$的零空间维数恰好等于$R(B)$）。含未知参数的秩讨论，方法是：先化行阶梯形，令某个主元恰为$0$，分别讨论参数取值。
}

\subsection{满秩矩阵}

设$A$为$m\times n$矩阵：
\begin{itemize}
\item 若$R(A)=m$，称$A$为\textbf{行满秩矩阵}；
\item 若$R(A)=n$，称$A$为\textbf{列满秩矩阵}；
\item 若$R(A)=\min\{m,n\}$，简称$A$为\textbf{满秩矩阵}；
\item $n$阶方阵$A$满秩$\iff R(A)=n\iff|A|\neq0\iff A$可逆。
\end{itemize}

满秩矩阵的等价标准形形态：
\[
A\ \text{列满秩}\iff A\cong\begin{pmatrix}E_n\\O\end{pmatrix};\qquad
A\ \text{行满秩}\iff A\cong\begin{pmatrix}E_m&O\end{pmatrix}.
\]

\section{线性方程组的解}

\subsection{解的判定定理}

设有$m$个方程、$n$个未知量的线性方程组
\[
\begin{cases}
a_{11}x_1+a_{12}x_2+\cdots+a_{1n}x_n=b_1,\\
a_{21}x_1+a_{22}x_2+\cdots+a_{2n}x_n=b_2,\\
\quad\cdots\cdots\cdots\cdots\\
a_{m1}x_1+a_{m2}x_2+\cdots+a_{mn}x_n=b_m,
\end{cases}
\]
记系数矩阵$A=(a_{ij})_{m\times n}$，未知量列向量$X=(x_1,\cdots,x_n)'$，常数项列向量$\beta=(b_1,\cdots,b_m)'$，则方程组可写成矩阵形式$AX=\beta$；记\textbf{增广矩阵}
\[
B=(A\mid\beta)=\left(\begin{array}{cccc|c}
a_{11}&a_{12}&\cdots&a_{1n}&b_1\\
a_{21}&a_{22}&\cdots&a_{2n}&b_2\\
\vdots&\vdots&&\vdots&\vdots\\
a_{m1}&a_{m2}&\cdots&a_{mn}&b_m
\end{array}\right).
\]
若把$A$按列分块为$A=(\alpha_1,\alpha_2,\cdots,\alpha_n)$，则方程组又可写成向量形式
\[
x_1\alpha_1+x_2\alpha_2+\cdots+x_n\alpha_n=\beta.
\]
于是“方程组有解”$\iff$“$\beta$可由$A$的列向量组线性表示”。若方程组有解，称它是\textbf{相容的}，否则称\textbf{不相容}。若$\beta\neq0$，称方程组为\textbf{非齐次}的；若$\beta=0$，即$AX=0$，称为\textbf{齐次}线性方程组（它总有零解，故总是相容的）。

\begin{theorem}[有解判定]
$n$元线性方程组$AX=\beta$有解的充要条件是
\[
R(A)=R(A\mid\beta).
\]
\end{theorem}

\begin{theorem}[解的个数]
设$R(A)=R(A\mid\beta)=r$，$n$为未知量个数，则
\begin{enumerate}
\item 若$r=n$，方程组有\textbf{唯一解}；
\item 若$r<n$，方程组有\textbf{无穷多解}，其中有$n-r$个自由未知量；
\item 若$R(A)<R(A\mid\beta)$，方程组\textbf{无解}。
\end{enumerate}
\end{theorem}

\keyformula{R(A)=R(A\mid\beta)=r:\quad r=n\Rightarrow\text{唯一解},\qquad r<n\Rightarrow\text{无穷多解};\qquad
R(A)<R(A\mid\beta)\Rightarrow\text{无解}}

判定定理的直观解释：把增广矩阵化为行阶梯形后，若出现形如$(0\ \cdots\ 0\mid d)$（$d\neq0$）的行，即“$0=d$”这样的矛盾方程，则无解——这正是$R(A)<R(A\mid\beta)$；若不出现矛盾行，则剩下的等式都能解，自由未知量个数就是$n-r$。

\subsection{齐次线性方程组}

对$AX=0$（$A$为$m\times n$矩阵），因$R(A)\le R(A\mid 0)=R(A)$恒成立，故\textbf{齐次方程组总是相容}（至少有零解）。

\begin{theorem}
设$A$为$m\times n$矩阵，则
\begin{enumerate}
\item $AX=0$只有零解$\iff R(A)=n$（$A$列满秩）；
\item $AX=0$有非零解$\iff R(A)<n$；
\item 特别地，当$m<n$（方程个数少于未知量个数）时，$AX=0$必有非零解；
\item 当$m=n$时，$AX=0$有非零解$\iff|A|=0$（与克拉默法则的推论一致）。
\end{enumerate}
\end{theorem}

当$R(A)=r<n$时，$AX=0$的全部解可表示为
\[
X=k_1\xi_1+k_2\xi_2+\cdots+k_{n-r}\xi_{n-r}\quad(k_1,\cdots,k_{n-r}\text{为任意常数}),
\]
其中$\xi_1,\cdots,\xi_{n-r}$是$AX=0$的$n-r$个线性无关的解向量（\textbf{基础解系}）。求法即下面的“自由未知量法”；基础解系与解空间的完整理论见第4章§5。

\subsection{非齐次线性方程组}

设$AX=\beta$（$\beta\neq0$），称$AX=0$为它的\textbf{导出组}（对应的齐次方程组）。

\begin{theorem}
设$\eta_1,\eta_2$都是$AX=\beta$的解，$\xi$是$AX=0$的解，则
\begin{enumerate}
\item $\eta_1-\eta_2$是$AX=0$的解；
\item $\eta_1+\xi$仍是$AX=\beta$的解。
\end{enumerate}
\end{theorem}

\begin{corollary}[通解结构]
若$AX=\beta$有解，$R(A)=r<n$，$\eta^{*}$是$AX=\beta$的一个\textbf{特解}，$\xi_1,\cdots,\xi_{n-r}$是导出组$AX=0$的基础解系，则$AX=\beta$的\textbf{通解}为
\[
X=\eta^{*}+k_1\xi_1+k_2\xi_2+\cdots+k_{n-r}\xi_{n-r}\quad(k_1,\cdots,k_{n-r}\text{为任意常数}).
\]
即“非齐次通解＝一个特解＋导出组通解”。
\end{corollary}

\subsection{用初等变换求通解（步骤）}

\begin{enumerate}
\item 写出增广矩阵$B=(A\mid\beta)$，只用\textbf{初等行变换}化为行阶梯形，判断$R(A)$与$R(A\mid\beta)$：若$R(A)<R(A\mid\beta)$，无解，停止；
\item 继续化为\textbf{行最简形}（首非零元为$1$，所在列其余元为$0$）；
\item 由行最简形写出同解的方程组：首非零元对应的未知量用其余未知量表示；
\item 把\textbf{自由未知量}（非首非零元对应的未知量）依次取$(1,0,\cdots,0),(0,1,\cdots,0),\cdots$，并令常数项为$0$，得导出组的基础解系$\xi_1,\cdots,\xi_{n-r}$；再令自由未知量全取$0$，得特解$\eta^*$；
\item 写出通解$X=\eta^*+k_1\xi_1+\cdots+k_{n-r}\xi_{n-r}$。
\end{enumerate}

\corepoint{
（1）解方程组\textbf{只能做初等行变换}，绝不能做列变换（列变换相当于把未知量改名，通解会对不上）。
（2）判断有无解必须用$R(A)$与$R(A\mid\beta)$，只看$R(A)$会误判。
（3）$R(A)=R(A\mid\beta)<n$时“有无穷多解”，答案必须写出\textbf{通解}含$n-r$个任意常数，只写一个特解算错。
（4）齐次方程组$AX=0$的基础解系必须取$n-R(A)$个\textbf{线性无关}的解；若自由未知量取值组选取不当，可能取到线性相关的向量组。
（5）含参数的方程组，先化阶梯形，再对参数使主元为零的取值逐个讨论（此时秩可能下降，解的个数随之改变）。
}

\subsection{与矩阵方程、矩阵秩的联系}

\begin{enumerate}
\item $AX=B$（$B$为$m\times p$矩阵）可视为$p$个“系数矩阵相同、常数项不同”的方程组同时求解，用$(A\mid B)\xrightarrow{\text{行变换}}$行最简形处理；若$A$可逆，则$X=A^{-1}B$。
\item $AX=0$有非零解$\iff R(A)<n$，这条判据在第4章用于判定向量组线性相关，在第5章用于求特征向量（$(\lambda E-A)X=0$有非零解$\iff|\lambda E-A|=0$）。
\item $R(AB)\le\min\{R(A),R(B)\}$与$AB=O\Rightarrow R(A)+R(B)\le n$可由“$B$的每一列都是$AX=0$的解”直接看出。
\end{enumerate}


% ---------------- 第 4 章 向量组的线性相关性 ----------------
% 教材 §1 向量组及其线性组合 / §2 向量组的线性相关性 / §3 向量组的秩
%      §4 向量空间 / §5 线性方程组的解的结构
\chapter{向量组的线性相关性}
% 第4章 向量组的线性相关性（同济《线性代数》第七版 第4章）
% 节次与教材一致：§1 向量组及其线性组合；§2 向量组的线性相关性；§3 向量组的秩；
% §4 向量空间（第七版把向量空间移到第4章第4节）；§5 线性方程组的解的结构。
% 只用 \section / \subsection，正文不写例题与习题。

\section{向量组及其线性组合}

\subsection{$n$维向量的概念}
\begin{definition}[$n$维向量]
数域$F$内的$n$个数$a_1,a_2,\dots,a_n$组成的有序数组$\boldsymbol{\alpha}$叫做数域$F$上的\textbf{$n$维向量}（简称向量），记作
\[
\boldsymbol{\alpha}=(a_1,a_2,\dots,a_n)
\quad\text{或}\quad
\boldsymbol{\alpha}=
\begin{pmatrix}
a_1\\ a_2\\ \vdots\\ a_n
\end{pmatrix}.
\]
前者称为\textbf{行向量}，后者称为\textbf{列向量}；数$a_i$称为该向量的第$i$个\textbf{分量}。分量全为实数的向量称为\textbf{实向量}，分量中含有复数的向量称为\textbf{复向量}。本书如无特别声明，只讨论实向量。
\end{definition}

按第2章的规定，行向量就是行矩阵，列向量就是列矩阵，行向量与列向量都按矩阵的运算规则进行运算。$n$维列向量与$n$维行向量是两个不同的向量；未指明行、列时，一律当作\textbf{列向量}，行向量用$\boldsymbol{\alpha}^{T}$等表示。

当$n=3$时，$n$维向量可看作以原点为起点的有向线段，即解析几何中的向量；当$n>3$时它不再有直观的几何形象，只是沿用几何术语。
\begin{definition}[向量组]
若干个同维数的列向量（或同维数的行向量）所组成的集合叫做\textbf{向量组}。一个$m\times n$矩阵的全体列向量是一个含$n$个$m$维列向量的向量组，其全体行向量是一个含$m$个$n$维行向量的向量组。
\end{definition}

含有限个向量的有序向量组与矩阵一一对应：$m$个$n$维列向量$\boldsymbol{\alpha}_1,\boldsymbol{\alpha}_2,\dots,\boldsymbol{\alpha}_m$构成$n\times m$矩阵
\[
A=(\boldsymbol{\alpha}_1,\boldsymbol{\alpha}_2,\dots,\boldsymbol{\alpha}_m),
\]
称$A$为该向量组所构成的矩阵。这是把向量组的讨论转化为矩阵讨论的桥梁。

\subsection{线性运算}
\begin{definition}[向量加法]
设$\boldsymbol{\alpha}=(a_1,a_2,\dots,a_n),\boldsymbol{\beta}=(b_1,b_2,\dots,b_n)$都是$n$维向量，则向量$(a_1+b_1,a_2+b_2,\dots,a_n+b_n)$叫做$\boldsymbol{\alpha}$与$\boldsymbol{\beta}$的\textbf{和}，记作$\boldsymbol{\alpha}+\boldsymbol{\beta}$。
\end{definition}

\begin{definition}[数乘]
设$\boldsymbol{\alpha}=(a_1,a_2,\dots,a_n)$为数域$F$上的$n$维向量，$k\in F$，则向量$(ka_1,ka_2,\dots,ka_n)$叫做数$k$与向量$\boldsymbol{\alpha}$的\textbf{乘积}，记作$k\boldsymbol{\alpha}$。
\end{definition}

向量的加法与数乘统称为向量的\textbf{线性运算}。对任意$\boldsymbol{\alpha},\boldsymbol{\beta},\boldsymbol{\gamma}$及$k,l\in F$，线性运算满足八条性质：
\begin{enumerate}
    \item $\boldsymbol{\alpha}+\boldsymbol{\beta}=\boldsymbol{\beta}+\boldsymbol{\alpha}$；
    \item $(\boldsymbol{\alpha}+\boldsymbol{\beta})+\boldsymbol{\gamma}=\boldsymbol{\alpha}+(\boldsymbol{\beta}+\boldsymbol{\gamma})$；
    \item $\boldsymbol{\alpha}+\boldsymbol{0}=\boldsymbol{\alpha}$（存在零向量）；
    \item $\boldsymbol{\alpha}+(-\boldsymbol{\alpha})=\boldsymbol{0}$（存在负向量）；
    \item $1\boldsymbol{\alpha}=\boldsymbol{\alpha}$；
    \item $k(l\boldsymbol{\alpha})=(kl)\boldsymbol{\alpha}$；
    \item $k(\boldsymbol{\alpha}+\boldsymbol{\beta})=k\boldsymbol{\alpha}+k\boldsymbol{\beta}$；
    \item $(k+l)\boldsymbol{\alpha}=k\boldsymbol{\alpha}+l\boldsymbol{\alpha}$。
\end{enumerate}
由这八条还可推出$0\boldsymbol{\alpha}=\boldsymbol{0}$，$(-1)\boldsymbol{\alpha}=-\boldsymbol{\alpha}$，$k\boldsymbol{0}=\boldsymbol{0}$。这八条正是第6章"线性空间"公理的原型。

\subsection{线性组合与线性表示}
\begin{definition}[线性组合、线性表示]
给定向量组$A:\boldsymbol{\alpha}_1,\boldsymbol{\alpha}_2,\dots,\boldsymbol{\alpha}_m$，对于任何一组实数$k_1,k_2,\dots,k_m$，表达式
\[
k_1\boldsymbol{\alpha}_1+k_2\boldsymbol{\alpha}_2+\dots+k_m\boldsymbol{\alpha}_m
\]
称为向量组$A$的一个\textbf{线性组合}，$k_1,k_2,\dots,k_m$称为这个线性组合的\textbf{系数}。

给定向量组$A:\boldsymbol{\alpha}_1,\dots,\boldsymbol{\alpha}_m$和向量$\boldsymbol{b}$，若存在一组数$\lambda_1,\dots,\lambda_m$使
\[
\boldsymbol{b}=\lambda_1\boldsymbol{\alpha}_1+\lambda_2\boldsymbol{\alpha}_2+\dots+\lambda_m\boldsymbol{\alpha}_m,
\]
则称向量$\boldsymbol{b}$是向量组$A$的线性组合，或称$\boldsymbol{b}$\textbf{能由向量组$A$线性表示}。
\end{definition}

"$\boldsymbol{b}$能由$A$线性表示"等价于线性方程组
\[
(\boldsymbol{\alpha}_1,\boldsymbol{\alpha}_2,\dots,\boldsymbol{\alpha}_m)
\begin{pmatrix}\lambda_1\\ \lambda_2\\ \vdots\\ \lambda_m\end{pmatrix}
=\boldsymbol{b},\qquad\text{即}\quad A\boldsymbol{x}=\boldsymbol{b}
\]
有解，解向量就是表示系数。

\begin{definition}[向量组的线性表示与等价]
设有两个向量组$A:\boldsymbol{\alpha}_1,\dots,\boldsymbol{\alpha}_r$及$B:\boldsymbol{\beta}_1,\dots,\boldsymbol{\beta}_s$。若$B$中每个向量都能由向量组$A$线性表示，则称\textbf{向量组$B$能由向量组$A$线性表示}；若$A$与$B$能相互线性表示，则称这两个\textbf{向量组等价}。
\end{definition}

将$A,B$所构成的矩阵分别记作$A=(\boldsymbol{\alpha}_1,\dots,\boldsymbol{\alpha}_r)$，$B=(\boldsymbol{\beta}_1,\dots,\boldsymbol{\beta}_s)$。$B$能由$A$线性表示，即存在矩阵$K=(k_{ij})_{r\times s}$使
\[
(\boldsymbol{\beta}_1,\dots,\boldsymbol{\beta}_s)=(\boldsymbol{\alpha}_1,\dots,\boldsymbol{\alpha}_r)K,
\quad\text{即}\quad B=AK,
\]
$K$称为这一线性表示的\textbf{系数矩阵}。于是：$C=AB$时，$C$的列向量组能由$A$的列向量组线性表示（$B$为系数矩阵），$C$的行向量组能由$B$的行向量组线性表示（$A$为系数矩阵）。

矩阵的行（列）等价与向量组的等价相对应：矩阵$A$经初等行变换化为$B$，则$A$的行向量组与$B$的行向量组等价；$A$与$B$列等价，则它们的列向量组等价。

\subsection{线性表示的秩判据}
\begin{theorem}[线性表示的秩判据]
向量$\boldsymbol{b}$能由向量组$A:\boldsymbol{\alpha}_1,\dots,\boldsymbol{\alpha}_m$线性表示的充分必要条件是矩阵$A=(\boldsymbol{\alpha}_1,\dots,\boldsymbol{\alpha}_m)$的秩等于矩阵$B=(A,\boldsymbol{b})$的秩，即
\[
R(A)=R(A,\boldsymbol{b})\quad\Longleftrightarrow\quad R(\boldsymbol{\alpha}_1,\dots,\boldsymbol{\alpha}_m)
=R(\boldsymbol{\alpha}_1,\dots,\boldsymbol{\alpha}_m,\boldsymbol{b}).
\]
\end{theorem}

\begin{theorem}[向量组线性表示的秩判据]
向量组$B:\boldsymbol{\beta}_1,\dots,\boldsymbol{\beta}_s$能由向量组$A:\boldsymbol{\alpha}_1,\dots,\boldsymbol{\alpha}_r$线性表示的充分必要条件是矩阵$A=(\boldsymbol{\alpha}_1,\dots,\boldsymbol{\alpha}_r)$的秩等于矩阵$(A,B)=(\boldsymbol{\alpha}_1,\dots,\boldsymbol{\alpha}_r,\boldsymbol{\beta}_1,\dots,\boldsymbol{\beta}_s)$的秩，即
\[
R(A)=R(A,B).
\]
\end{theorem}

\begin{corollary}[向量组等价的秩判据]
向量组$A$与向量组$B$等价的充分必要条件是
\[
R(A)=R(B)=R(A,B).
\]
\end{corollary}

\begin{theorem}[表示关系的秩不等式]
若向量组$B$能由向量组$A$线性表示，则
\[
R(B)\le R(A).
\]
\end{theorem}

由$B=AK$及矩阵秩的性质$R(AK)\le R(A)$即得。此定理说明：向量组经线性表示后秩不会增大；等价的向量组必有相同的秩（但秩相同未必等价）。

\keyformula{\boldsymbol{b}\ \text{可由}\ A\ \text{表示}\iff R(A)=R(A,\boldsymbol{b}),\qquad B\ \text{可由}\ A\ \text{表示}\iff R(A)=R(A,B)}

\keyformula{A\ \text{与}\ B\ \text{等价}\iff R(A)=R(B)=R(A,B),\qquad B\ \text{可由}\ A\ \text{表示}\ \Rightarrow\ R(B)\le R(A)}

\subsection{判定与计算步骤}
\begin{enumerate}
    \item 写增广矩阵：把$\boldsymbol{b}$（或向量组$B$的每个向量）排在$A$的右边，得$(A,\boldsymbol{b})$或$(A,B)$。
    \item 用初等行变换把$(A,\boldsymbol{b})$化为行阶梯形，比较$R(A)$与$R(A,\boldsymbol{b})$：相等则能表示，否则不能。
    \item 若需要求出表示式，继续把$(A,\boldsymbol{b})$化为\textbf{行最简形}，写出$A\boldsymbol{x}=\boldsymbol{b}$的通解；通解中的自由未知数任意取值，就给出了表示式（表示式一般不唯一，除非$R(A)=m$，此时唯一）。
\end{enumerate}

\corepoint{
三个易错点：\\
（1）比较秩时一定比较$R(A)$与$R(A,\boldsymbol{b})$，不能只算$R(\boldsymbol{b})$；\\
（2）"不能表示"的判据是$R(A)\ne R(A,\boldsymbol{b})$（此时必是$R(A,\boldsymbol{b})=R(A)+1$）；\\
（3）表示系数不唯一时，不同自由未知数取值对应不同的表示式，但都有效；若题目要求"写出表示式"，须写出带任意常数的通式。
}

\section{向量组的线性相关性}

\subsection{线性相关与线性无关}
\begin{definition}[线性相关、线性无关]
给定向量组$A:\boldsymbol{\alpha}_1,\boldsymbol{\alpha}_2,\dots,\boldsymbol{\alpha}_m$，若存在\textbf{不全为零}的数$k_1,k_2,\dots,k_m$使
\[
k_1\boldsymbol{\alpha}_1+k_2\boldsymbol{\alpha}_2+\dots+k_m\boldsymbol{\alpha}_m=\boldsymbol{0},
\]
则称向量组$A$是\textbf{线性相关}的；否则称它\textbf{线性无关}，即
\[
k_1\boldsymbol{\alpha}_1+\dots+k_m\boldsymbol{\alpha}_m=\boldsymbol{0}
\ \Longrightarrow\ k_1=k_2=\dots=k_m=0.
\]
\end{definition}

几点直接结论：
\begin{enumerate}
    \item 只含一个向量的向量组：$\boldsymbol{\alpha}=\boldsymbol{0}$时线性相关，$\boldsymbol{\alpha}\ne\boldsymbol{0}$时线性无关；
    \item 含两个向量$\boldsymbol{\alpha}_1,\boldsymbol{\alpha}_2$的向量组线性相关$\iff$两向量的分量对应成比例，几何意义是两向量共线；
    \item 三个向量线性相关的几何意义是三向量共面；
    \item 含零向量的向量组一定线性相关。
\end{enumerate}

线性相关与线性表示的联系：向量组$A:\boldsymbol{\alpha}_1,\dots,\boldsymbol{\alpha}_m$线性相关，等价于齐次线性方程组
\[
x_1\boldsymbol{\alpha}_1+x_2\boldsymbol{\alpha}_2+\dots+x_m\boldsymbol{\alpha}_m=\boldsymbol{0},
\quad\text{即}\quad A\boldsymbol{x}=\boldsymbol{0}
\]
有非零解。

\subsection{相关性的秩判据}
\begin{theorem}[线性相关性的秩判据]
向量组$\boldsymbol{\alpha}_1,\boldsymbol{\alpha}_2,\dots,\boldsymbol{\alpha}_m$线性相关的充分必要条件是它所构成的矩阵$A=(\boldsymbol{\alpha}_1,\dots,\boldsymbol{\alpha}_m)$的秩小于向量个数$m$；向量组线性无关的充分必要条件是$R(A)=m$。即
\[
\boldsymbol{\alpha}_1,\dots,\boldsymbol{\alpha}_m\ \text{线性相关}\iff R(A)<m,
\qquad
\boldsymbol{\alpha}_1,\dots,\boldsymbol{\alpha}_m\ \text{线性无关}\iff R(A)=m.
\]
\end{theorem}

\textbf{证明}\quad 必要性：设向量组线性相关，则存在不全为零的$x_1,\dots,x_m$使$A\boldsymbol{x}=\boldsymbol{0}$，即齐次方程组$A\boldsymbol{x}=\boldsymbol{0}$有非零解，由第3章定理知$R(A)<m$。充分性：若$R(A)<m$，则$A\boldsymbol{x}=\boldsymbol{0}$有非零解$\boldsymbol{x}=(k_1,\dots,k_m)^{T}\ne\boldsymbol{0}$，即$k_1\boldsymbol{\alpha}_1+\dots+k_m\boldsymbol{\alpha}_m=\boldsymbol{0}$且系数不全为零，故向量组线性相关。\hfill$\square$

这一判据把"线性相关性"完全转化为"矩阵的秩"，是本章全部计算的基础。

\begin{corollary}
当$A$为$n$阶方阵时，其列向量组$\boldsymbol{\alpha}_1,\dots,\boldsymbol{\alpha}_n$线性无关的充分必要条件是$|A|\ne 0$；线性相关的充分必要条件是$|A|=0$。
\end{corollary}

\begin{corollary}[加长向量组]
设$\boldsymbol{\alpha}_i=(a_{1i},a_{2i},\dots,a_{ri})^{T}$，$\boldsymbol{\beta}_i=(a_{1i},\dots,a_{ri},a_{r+1,i},\dots,a_{r+s,i})^{T}$，$i=1,2,\dots,m$，即$\boldsymbol{\beta}_i$是$\boldsymbol{\alpha}_i$的"加长"。若$r$维向量组$\boldsymbol{\alpha}_1,\dots,\boldsymbol{\alpha}_m$线性无关，则$r+s$维向量组$\boldsymbol{\beta}_1,\dots,\boldsymbol{\beta}_m$也线性无关。
\end{corollary}

等价说法：\textbf{线性无关的向量组，截短（去掉分量）后未必无关；而线性相关的向量组，加长后仍相关。}（"短无关$\Rightarrow$长无关"，"长相关$\Rightarrow$短相关"。）

\begin{theorem}[向量个数多于维数必相关]
设$\boldsymbol{\alpha}_1,\boldsymbol{\alpha}_2,\dots,\boldsymbol{\alpha}_m$是$m$个$n$维向量，若$m>n$，则向量组必定线性相关。
\end{theorem}

\textbf{证明}\quad $A=(\boldsymbol{\alpha}_1,\dots,\boldsymbol{\alpha}_m)$是$n\times m$矩阵，$R(A)\le\min\{n,m\}=n<m$，由秩判据知线性相关。\hfill$\square$

\subsection{相关性的基本结论}
\begin{theorem}[相关性的三条基本结论]
\begin{enumerate}
    \item 若向量组$A:\boldsymbol{\alpha}_1,\dots,\boldsymbol{\alpha}_r$线性相关，则加上若干向量所得的向量组$B$也线性相关；反之，若$B$线性无关，则其任何部分组（特别地$A$）都线性无关。
    \item $m$个$n$维向量组成的向量组，当$n<m$时一定线性相关；特别地，$n+1$个$n$维向量一定线性相关。
    \item 设向量组$A:\boldsymbol{\alpha}_1,\dots,\boldsymbol{\alpha}_m$线性无关，而向量组$B:\boldsymbol{\alpha}_1,\dots,\boldsymbol{\alpha}_m,\boldsymbol{b}$线性相关，则$\boldsymbol{b}$必能由$A$线性表示，且表示式唯一。
\end{enumerate}
\end{theorem}

\textbf{证明}\quad (1) 若有不全为零的$k_1,\dots,k_r$使$\sum k_i\boldsymbol{\alpha}_i=\boldsymbol{0}$，则把它们与$B$中新增向量的系数取$0$，便得$B$的非平凡零组合，故$B$相关。\\
(3) 由$B$相关，存在不全为零的$\lambda_1,\dots,\lambda_m,\lambda$使$\lambda_1\boldsymbol{\alpha}_1+\dots+\lambda_m\boldsymbol{\alpha}_m+\lambda\boldsymbol{b}=\boldsymbol{0}$。若$\lambda=0$，则$\boldsymbol{\alpha}_1,\dots,\boldsymbol{\alpha}_m$线性相关，矛盾，故$\lambda\ne0$，于是
\[
\boldsymbol{b}=-\frac{\lambda_1}{\lambda}\boldsymbol{\alpha}_1-\dots-\frac{\lambda_m}{\lambda}\boldsymbol{\alpha}_m .
\]
唯一性：若有两种表示，相减得$\boldsymbol{\alpha}_1,\dots,\boldsymbol{\alpha}_m$的一个零组合，由无关性知系数全为零，两种表示相同。\hfill$\square$

\begin{theorem}["至少一个可由其余表示"]
向量组$\boldsymbol{\alpha}_1,\boldsymbol{\alpha}_2,\dots,\boldsymbol{\alpha}_m$（$m\ge2$）线性相关的充分必要条件是其中至少有一个向量能由其余$m-1$个向量线性表示。
\end{theorem}

注意：线性相关时\textbf{并非每个}向量都能由其余向量表示。例如$\boldsymbol{\alpha}_1=(1,0),\boldsymbol{\alpha}_2=(2,0),\boldsymbol{\alpha}_3=(0,1)$线性相关，但$\boldsymbol{\alpha}_3$不能由$\boldsymbol{\alpha}_1,\boldsymbol{\alpha}_2$表示。这是最容易出错的一处。

\subsection{判定步骤与常用手法}
\begin{enumerate}
    \item \textbf{定义法}：设$k_1\boldsymbol{\alpha}_1+\dots+k_m\boldsymbol{\alpha}_m=\boldsymbol{0}$，解出$k_i$是否必须全为零。
    \item \textbf{秩法（首选）}：写出矩阵$A=(\boldsymbol{\alpha}_1,\dots,\boldsymbol{\alpha}_m)$，用初等行变换求$R(A)$，与$m$比较。$R(A)<m$相关，$R(A)=m$无关。
    \item \textbf{行列式法}：当向量个数等于维数（方阵）时，算$|A|$，$|A|=0$相关，$|A|\ne0$无关。
    \item \textbf{矩阵等式法}：若$(\boldsymbol{\beta}_1,\dots,\boldsymbol{\beta}_s)=(\boldsymbol{\alpha}_1,\dots,\boldsymbol{\alpha}_r)K$且$\boldsymbol{\alpha}$组无关，则$\boldsymbol{\beta}$组无关$\iff R(K)=s$。
\end{enumerate}

\keyformula{k_1\boldsymbol{\alpha}_1+\dots+k_m\boldsymbol{\alpha}_m=\boldsymbol{0}\ \text{有非零解}\iff A\boldsymbol{x}=\boldsymbol{0}\ \text{有非零解}\iff R(A)<m}

\keyformula{A=(\boldsymbol{\alpha}_1,\dots,\boldsymbol{\alpha}_m)\ \text{为方阵时：}\ \text{无关}\iff |A|\ne0,\qquad m>n\ \Rightarrow\ \text{必相关}}

\corepoint{
判断相关性的三条"铁律"与三个陷阱：\\
（1）秩与"向量个数"比，不是与"维数"比：$R(A)=m$才无关；\\
（2）"$m>n$必相关"只说个数超过维数，反过来不成立——$m\le n$也可能相关；\\
（3）相关$\Rightarrow$存在某个向量可由其余表示，但\textbf{不能指定是哪一个}；判断"$\boldsymbol{\beta}$能否由$\boldsymbol{\alpha}$组表示"仍要用"线性表示的秩判据"。
}

\section{向量组的秩}

\subsection{最大无关组与秩}
\begin{definition}[最大线性无关向量组]
设有向量组$A$，如果在$A$中能选出$r$个向量$\boldsymbol{\alpha}_1,\dots,\boldsymbol{\alpha}_r$，满足
\begin{enumerate}
    \item[(i)] 向量组$A_0:\boldsymbol{\alpha}_1,\dots,\boldsymbol{\alpha}_r$线性无关；
    \item[(ii)] 向量组$A$中任意$r+1$个向量（如果$A$中有$r+1$个向量的话）都线性相关，
\end{enumerate}
那么称$A_0$是向量组$A$的一个\textbf{最大线性无关向量组}（简称\textbf{最大无关组}）。
\end{definition}

\begin{definition}[向量组的秩]
最大无关组所含向量的个数$r$称为\textbf{向量组$A$的秩}，记作$R(A)$或$R(\boldsymbol{\alpha}_1,\dots,\boldsymbol{\alpha}_m)$。只含零向量的向量组没有最大无关组，规定它的秩为$0$。
\end{definition}

由定义直接得到：
\begin{enumerate}
    \item 若向量组$A$线性无关，则$A$本身就是它的最大无关组，且$R(A)=m$；反之$R(A)=m$时$A$线性无关。
    \item 向量组$A$与它的任一最大无关组等价（最大无关组是$A$的部分组，故能被$A$表示；而$A$中任一向量与$A_0$组成的$r+1$个向量相关，由"无关组添加向量后相关，则新向量可由原组唯一表示"知它可由$A_0$表示）。
    \item 一个向量组的最大无关组一般不唯一，但所含向量个数相同（见下节定理）。
    \item $n$维向量组至多含$n$个向量的最大无关组，故$R(A)\le n$。
\end{enumerate}

\begin{corollary}[最大无关组的等价定义]
设向量组$A_0:\boldsymbol{\alpha}_1,\dots,\boldsymbol{\alpha}_r$是向量组$A$的一个部分组，且满足
\begin{enumerate}
    \item[(i)] $A_0$线性无关；
    \item[(ii)] $A$中任一向量都能由$A_0$线性表示，
\end{enumerate}
则$A_0$便是$A$的最大无关组。
\end{corollary}

\textbf{证明}\quad 只需证$A$中任意$r+1$个向量线性相关。设$\boldsymbol{b}_1,\dots,\boldsymbol{b}_{r+1}$是$A$中任取的$r+1$个向量，由条件(ii)它们都能由$A_0$线性表示，故其构成的矩阵的秩$R\le R(A_0)=r<r+1$，由秩判据知它们线性相关。\hfill$\square$

用这个等价定义求秩最方便：只需找一组"既无关、又能表示全体"的向量。

\subsection{秩与矩阵的秩}
\begin{theorem}
设有两个$n$维向量组$A:\boldsymbol{\alpha}_1,\dots,\boldsymbol{\alpha}_r$，$B:\boldsymbol{\beta}_1,\dots,\boldsymbol{\beta}_s$。若$A$线性无关，且$A$能由$B$线性表示，则$r\le s$。
\end{theorem}

\begin{corollary}
若向量组$A,B$都线性无关且等价，则$r=s$；进而，同一向量组的任意两个最大无关组等价，且所含向量个数相同。
\end{corollary}

\begin{theorem}[矩阵的秩等于行秩、列秩]
设$A$是$n\times m$矩阵，则$A$的列向量组的秩等于$A$的秩，也等于$A$的行向量组的秩，即
\[
R(A)=A\ \text{的列向量组的秩}=A\ \text{的行向量组的秩}.
\]
\end{theorem}

\textbf{证明}\quad 设$R(A)=r$，则$A$至少含一个$r$阶子式$D_r\ne0$。若$D_r$位于$A$的第$k_1<k_2<\dots<k_r$列上，则由这$r$个列向量构成的矩阵$A_1=(\boldsymbol{\alpha}_{k_1},\dots,\boldsymbol{\alpha}_{k_r})$含非零$r$阶子式，故$R(A_1)=r$，即这$r$列线性无关。又任取$A$的$r+1$个列向量，它们构成的矩阵是$A$的子矩阵，其秩$\le R(A)=r<r+1$，故这$r+1$列线性相关。因此$\boldsymbol{\alpha}_{k_1},\dots,\boldsymbol{\alpha}_{k_r}$是列向量组的一个最大无关组，列向量组的秩为$r$。行向量组同理（因$R(A)=R(A^{T})$）。\hfill$\square$

由证明可见：\textbf{$D_r$所在的$r$个列（行）向量就是列（行）向量组的一个最大无关组。}

\begin{theorem}
若向量组$B$能由向量组$A$线性表示，则$R(B)\le R(A)$；若$A,B$等价，则$R(A)=R(B)$。
\end{theorem}

增加了$r$个向量，秩最多增加$r$：$R(\boldsymbol{\alpha}_1,\dots,\boldsymbol{\alpha}_m,\boldsymbol{\beta}_1,\dots,\boldsymbol{\beta}_s)\le R(\boldsymbol{\alpha}_1,\dots,\boldsymbol{\alpha}_m)+R(\boldsymbol{\beta}_1,\dots,\boldsymbol{\beta}_s)$。

\subsection{求秩与最大无关组的步骤}
\begin{enumerate}
    \item 以$\boldsymbol{\alpha}_1,\dots,\boldsymbol{\alpha}_m$为列构成矩阵$A$。
    \item 只用\textbf{初等行变换}把$A$化为行阶梯形（行变换不改变列向量组的线性关系，这是关键）。
    \item 行阶梯形中\textbf{非零行的行数}$=R(A)=$向量组的秩。
    \item 行阶梯形中\textbf{每个非零行首元所在的列号}$k_1,\dots,k_r$，对应的原向量$\boldsymbol{\alpha}_{k_1},\dots,\boldsymbol{\alpha}_{k_r}$就是一个最大无关组。
    \item 若要把其余向量用最大无关组表示，把$A$继续化为\textbf{行最简形}$B=(\boldsymbol{b}_1,\dots,\boldsymbol{b}_m)$，则$A\boldsymbol{x}=\boldsymbol{0}$与$B\boldsymbol{x}=\boldsymbol{0}$同解，读出$B$中非首元列的关系式，即得原向量间的同一关系式。
\end{enumerate}

\corepoint{
（1）求最大无关组\textbf{只能用行变换}。若用了列变换，列向量组的对应关系被破坏，"首元所在列"不再是原向量的编号。\\
（2）秩相同不等于等价：$(\boldsymbol{\alpha}_1)=(1,0)^{T}$与$(\boldsymbol{\alpha}_2)=(2,0)^{T}$秩都是1、互相等价；但$(1,0)^{T}$与$(0,1)^{T}$秩都为1却互不能表示。判等价必须验$R(A)=R(B)=R(A,B)$。\\
（3）零向量组（只含零向量）的秩规定为$0$，没有最大无关组。
}

\keyformula{R(\boldsymbol{\alpha}_1,\dots,\boldsymbol{\alpha}_m)=R\bigl((\boldsymbol{\alpha}_1,\dots,\boldsymbol{\alpha}_m)\bigr)=\text{行阶梯形非零行数}}

\keyformula{\boldsymbol{\alpha}\ \text{组无关}\iff R(\boldsymbol{\alpha}_1,\dots,\boldsymbol{\alpha}_m)=m,\qquad B\ \text{可由}\ A\ \text{表示}\Rightarrow R(B)\le R(A)}

\section{向量空间}

\subsection{向量空间与子空间}
\begin{definition}[向量空间]
设$V$为$n$维向量的集合，如果集合$V$非空，且集合$V$对于向量的加法及数乘两种运算封闭，那么就称集合$V$为\textbf{向量空间}。所谓封闭，是指
\[
\boldsymbol{a}\in V,\ \boldsymbol{b}\in V\ \Rightarrow\ \boldsymbol{a}+\boldsymbol{b}\in V;
\qquad
\boldsymbol{a}\in V,\ \lambda\in\mathbb{R}\ \Rightarrow\ \lambda\boldsymbol{a}\in V.
\]
\end{definition}

$n$维向量的全体$\mathbb{R}^{n}$是向量空间。当$n>3$时它没有直观几何意义，但$n=3$时它就是取定坐标原点后的点空间。

判断一个集合是否为向量空间，实质是检验两条封闭性（八条运算律在$\mathbb{R}^{n}$中自动成立）。典型结论：
\begin{enumerate}
    \item 齐次线性方程组$A\boldsymbol{x}=\boldsymbol{0}$的解集$S=\{\boldsymbol{x}\mid A\boldsymbol{x}=\boldsymbol{0}\}$是向量空间，称为\textbf{解空间}；
    \item 非齐次线性方程组$A\boldsymbol{x}=\boldsymbol{b}$（$\boldsymbol{b}\ne\boldsymbol{0}$）的解集\textbf{不是}向量空间（对加法、数乘都不封闭）；
    \item 形如$(0,a_2,\dots,a_n)^{T}$的向量全体是向量空间（第一个分量恒为0，对运算封闭）；
    \item 由$\boldsymbol{\alpha},\boldsymbol{b}$生成的集合$V=\{\boldsymbol{x}\mid \boldsymbol{x}=\lambda\boldsymbol{\alpha}+\mu\boldsymbol{b},\ \lambda,\mu\in\mathbb{R}\}$是向量空间；一般地，由向量组$\boldsymbol{\alpha}_1,\dots,\boldsymbol{\alpha}_m$生成的
    \[
    L=\{\boldsymbol{x}\mid \boldsymbol{x}=\lambda_1\boldsymbol{\alpha}_1+\dots+\lambda_m\boldsymbol{\alpha}_m,\ \lambda_i\in\mathbb{R}\}
    \]
    是向量空间，记作$L=\operatorname{Span}\{\boldsymbol{\alpha}_1,\dots,\boldsymbol{\alpha}_m\}$，它与向量组$\boldsymbol{\alpha}_1,\dots,\boldsymbol{\alpha}_m$等价。若两向量组等价，则它们生成的向量空间相等。
\end{enumerate}

\begin{definition}[子空间]
设有向量空间$V$及$V_0$，若$V_0\subseteq V$，就称$V_0$是$V$的\textbf{子空间}。任何由$n$维向量组成的向量空间都是$\mathbb{R}^{n}$的子空间。
\end{definition}

\subsection{维数与坐标}
\begin{definition}[基与维数]
设$V$为向量空间，如果$r$个向量$\boldsymbol{\alpha}_1,\dots,\boldsymbol{\alpha}_r\in V$，且满足
\begin{enumerate}
    \item[(i)] $\boldsymbol{\alpha}_1,\dots,\boldsymbol{\alpha}_r$线性无关；
    \item[(ii)] $V$中任一向量都可由$\boldsymbol{\alpha}_1,\dots,\boldsymbol{\alpha}_r$线性表示，
\end{enumerate}
那么称$\boldsymbol{\alpha}_1,\dots,\boldsymbol{\alpha}_r$为向量空间$V$的一个\textbf{基}，$r$称为向量空间$V$的\textbf{维数}，记作$\dim V=r$，并称$V$为$r$维向量空间。只含零向量的向量空间没有基，规定其维数为$0$。
\end{definition}

把向量空间看作向量组，则$V$的基就是向量组的最大无关组，$V$的维数就是向量组的秩；$V$可表示为$V=\operatorname{Span}\{\boldsymbol{\alpha}_1,\dots,\boldsymbol{\alpha}_r\}$。

关于基的三点注意：
\begin{enumerate}
    \item 同一向量空间的任意两组基所含向量个数相同（都等于$\dim V$）；
    \item $r$维向量空间$V$中任意$r$个线性无关的向量都是$V$的基；
    \item 若$\boldsymbol{\alpha}_1,\dots,\boldsymbol{\alpha}_r$能线性表示$V$中任一向量，且$r=\dim V$，则它们是$V$的一个基。
\end{enumerate}

\begin{theorem}[基扩充定理]
有限维向量空间$V$中任意一组线性无关的向量都可以扩充成$V$的一个基。
\end{theorem}

\textbf{证明}\quad 设$V$是有限维的，存在有限向量组$S$使$V=\operatorname{Span}S$。设$\boldsymbol{\alpha}_1,\dots,\boldsymbol{\alpha}_s$线性无关。从$S$中依次取向量，凡能保持线性无关的就添入，得到含$\boldsymbol{\alpha}_1,\dots,\boldsymbol{\alpha}_s$的向量组$A_0$，它线性无关且在$S$中极大。于是$S$中每个向量都可由$A_0$表示，从而$V=\operatorname{Span}S\subseteq\operatorname{Span}A_0$，故$A_0$是$V$的基，且包含原来的$\boldsymbol{\alpha}_1,\dots,\boldsymbol{\alpha}_s$。\hfill$\square$

\begin{definition}[坐标]
设$\boldsymbol{\alpha}_1,\dots,\boldsymbol{\alpha}_r$是$r$维向量空间$V$的一个基，对任意$\boldsymbol{\alpha}\in V$，若有$x_1,\dots,x_r\in\mathbb{R}$使
\[
\boldsymbol{\alpha}=x_1\boldsymbol{\alpha}_1+x_2\boldsymbol{\alpha}_2+\dots+x_r\boldsymbol{\alpha}_r,
\]
则有序数组$(x_1,x_2,\dots,x_r)$称为向量$\boldsymbol{\alpha}$在基$\boldsymbol{\alpha}_1,\dots,\boldsymbol{\alpha}_r$下的\textbf{坐标}。基给定后，坐标是唯一的。
\end{definition}

在$\mathbb{R}^{n}$中取单位坐标向量组$\boldsymbol{e}_1,\dots,\boldsymbol{e}_n$为基，则以$(x_1,\dots,x_n)^{T}$为分量的向量$\boldsymbol{x}$在该基下的坐标就是$(x_1,\dots,x_n)^{T}$本身，故$\boldsymbol{e}_1,\dots,\boldsymbol{e}_n$称为$\mathbb{R}^{n}$的\textbf{自然基}。

\subsection{基变换与坐标变换}
\begin{definition}[基变换公式与过渡矩阵]
设$\boldsymbol{\alpha}_1,\dots,\boldsymbol{\alpha}_n$及$\boldsymbol{\beta}_1,\dots,\boldsymbol{\beta}_n$是向量空间$V$的两个基，且
\[
\begin{cases}
\boldsymbol{\beta}_1=p_{11}\boldsymbol{\alpha}_1+p_{21}\boldsymbol{\alpha}_2+\dots+p_{n1}\boldsymbol{\alpha}_n,\\
\boldsymbol{\beta}_2=p_{12}\boldsymbol{\alpha}_1+p_{22}\boldsymbol{\alpha}_2+\dots+p_{n2}\boldsymbol{\alpha}_n,\\
\qquad\qquad\cdots\cdots\cdots\cdots\\
\boldsymbol{\beta}_n=p_{1n}\boldsymbol{\alpha}_1+p_{2n}\boldsymbol{\alpha}_2+\dots+p_{nn}\boldsymbol{\alpha}_n,
\end{cases}
\]
即
\[
(\boldsymbol{\beta}_1,\boldsymbol{\beta}_2,\dots,\boldsymbol{\beta}_n)
=(\boldsymbol{\alpha}_1,\boldsymbol{\alpha}_2,\dots,\boldsymbol{\alpha}_n)
\begin{pmatrix}
p_{11}&p_{12}&\dots&p_{1n}\\
p_{21}&p_{22}&\dots&p_{2n}\\
\vdots&\vdots&\ddots&\vdots\\
p_{n1}&p_{n2}&\dots&p_{nn}
\end{pmatrix}
=(\boldsymbol{\alpha}_1,\dots,\boldsymbol{\alpha}_n)P .
\]
该式称为\textbf{基变换公式}，矩阵$P$称为由基$\boldsymbol{\alpha}_1,\dots,\boldsymbol{\alpha}_n$到基$\boldsymbol{\beta}_1,\dots,\boldsymbol{\beta}_n$的\textbf{过渡矩阵}。由$\boldsymbol{\beta}_1,\dots,\boldsymbol{\beta}_n$线性无关知$P$可逆。
\end{definition}

注意$P$的各列恰是新基向量在旧基下的坐标——即"新基$=$旧基$\times P$"，$P$的列是$\boldsymbol{\beta}_j$的坐标，方向不能反。

\begin{theorem}[坐标变换公式]
设向量空间$V$的基$\boldsymbol{\alpha}_1,\dots,\boldsymbol{\alpha}_n$到基$\boldsymbol{\beta}_1,\dots,\boldsymbol{\beta}_n$的过渡矩阵是$P$，$V$中向量$\boldsymbol{\alpha}$在基$\boldsymbol{\alpha}_1,\dots,\boldsymbol{\alpha}_n$下的坐标是$(x_1,\dots,x_n)^{T}$，在基$\boldsymbol{\beta}_1,\dots,\boldsymbol{\beta}_n$下的坐标是$(x_1',\dots,x_n')^{T}$，则
\[
\begin{pmatrix}x_1\\x_2\\\vdots\\x_n\end{pmatrix}
=P
\begin{pmatrix}x_1'\\x_2'\\\vdots\\x_n'\end{pmatrix},
\qquad\text{即}\quad
\begin{pmatrix}x_1'\\x_2'\\\vdots\\x_n'\end{pmatrix}
=P^{-1}
\begin{pmatrix}x_1\\x_2\\\vdots\\x_n\end{pmatrix}.
\]
\end{theorem}

\textbf{证明}\quad 由$\boldsymbol{\alpha}=(\boldsymbol{\alpha}_1,\dots,\boldsymbol{\alpha}_n)\boldsymbol{x}=(\boldsymbol{\beta}_1,\dots,\boldsymbol{\beta}_n)\boldsymbol{x}'$及$(\boldsymbol{\beta}_1,\dots,\boldsymbol{\beta}_n)=(\boldsymbol{\alpha}_1,\dots,\boldsymbol{\alpha}_n)P$，得$(\boldsymbol{\alpha}_1,\dots,\boldsymbol{\alpha}_n)\boldsymbol{x}=(\boldsymbol{\alpha}_1,\dots,\boldsymbol{\alpha}_n)P\boldsymbol{x}'$，由$\boldsymbol{\alpha}_1,\dots,\boldsymbol{\alpha}_n$线性无关得$\boldsymbol{x}=P\boldsymbol{x}'$。\hfill$\square$

\keyformula{(\boldsymbol{\beta}_1,\dots,\boldsymbol{\beta}_n)=(\boldsymbol{\alpha}_1,\dots,\boldsymbol{\alpha}_n)P\quad(\text{基变换}),\qquad \boldsymbol{x}=P\boldsymbol{x}',\quad \boldsymbol{x}'=P^{-1}\boldsymbol{x}\quad(\text{坐标变换})}

\keyformula{\dim\operatorname{Span}\{\boldsymbol{\alpha}_1,\dots,\boldsymbol{\alpha}_m\}=R(\boldsymbol{\alpha}_1,\dots,\boldsymbol{\alpha}_m)}

\corepoint{
（1）方向："由旧基到新基"的过渡矩阵$P$满足$(\text{新基})=(\text{旧基})P$，$P$的\textbf{第$j$列}是新基第$j$个向量在旧基下的坐标。\\
（2）坐标变换是"反着走"的：$\boldsymbol{x}=P\boldsymbol{x}'$（旧坐标$=$过渡矩阵$\times$新坐标）。记混方向是本节最常见的错误。\\
（3）求$P$的标准做法：把$(\text{旧基}\mid\text{新基})$用初等行变换化为$(E\mid P)$，即$P=A^{-1}B$，其中$A,B$分别以旧、新基向量为列。
}

\section{线性方程组的解的结构}

本节从"解空间"的角度重述第3章用初等变换得到的结论，是第4章向量理论的应用。

\subsection{齐次线性方程组的解集}
设齐次线性方程组
\[
\begin{cases}
a_{11}x_1+a_{12}x_2+\dots+a_{1n}x_n=0,\\
a_{21}x_1+a_{22}x_2+\dots+a_{2n}x_n=0,\\
\qquad\qquad\cdots\cdots\cdots\cdots\\
a_{m1}x_1+a_{m2}x_2+\dots+a_{mn}x_n=0,
\end{cases}
\]
记系数矩阵$A=(a_{ij})_{m\times n}$，则方程组可写作$A\boldsymbol{x}=\boldsymbol{0}$。其解集
\[
S=\{\boldsymbol{x}\mid A\boldsymbol{x}=\boldsymbol{0}\}\subseteq\mathbb{R}^{n}
\]
对加法与数乘封闭，是一个向量空间，称为该齐次方程组的\textbf{解空间}。若$\boldsymbol{\xi}_1,\dots,\boldsymbol{\xi}_{t}$是解空间的一个基，则方程组的通解为
\[
\boldsymbol{x}=c_1\boldsymbol{\xi}_1+c_2\boldsymbol{\xi}_2+\dots+c_t\boldsymbol{\xi}_{t}\qquad(c_1,\dots,c_t\in\mathbb{R}).
\]

\begin{definition}[基础解系]
齐次线性方程组$A\boldsymbol{x}=\boldsymbol{0}$的解空间的一个基，称为该齐次线性方程组的\textbf{基础解系}。
\end{definition}

求通解只需求出基础解系。

\begin{theorem}[解空间的维数]
设$m\times n$矩阵$A$的秩$R(A)=r$，则$n$元齐次线性方程组$A\boldsymbol{x}=\boldsymbol{0}$的解空间$S$的维数为
\[
\dim S=n-r,
\]
即基础解系恰含$n-r$个向量（$n-r$也就是自由未知数的个数）。
\end{theorem}

当$R(A)=n$时方程组只有零解，没有基础解系（$S$只含零向量）；当$R(A)=r<n$时，基础解系含$n-r$个向量，且任意$n-r$个线性无关的解都可构成基础解系，故基础解系不唯一，通解形式也不唯一。

\subsection{求基础解系的步骤}
\begin{enumerate}
    \item 对$A$作初等行变换，化为\textbf{行最简形}，定出首元所在列（对应"主变量"）与其余列（对应"自由未知数"）。
    \item 把$n-r$个自由未知数依次取$(1,0,\dots,0),(0,1,\dots,0),\dots,(0,\dots,0,1)$，回代求出主变量。
    \item 这样得到的$n-r$个解向量即为基础解系$\boldsymbol{\xi}_1,\dots,\boldsymbol{\xi}_{n-r}$，它们线性无关（因为自由未知数部分的坐标向量是单位坐标向量组）。
    \item 写出通解$\boldsymbol{x}=c_1\boldsymbol{\xi}_1+\dots+c_{n-r}\boldsymbol{\xi}_{n-r}$。
\end{enumerate}

\subsection{非齐次线性方程组解的结构}
设非齐次线性方程组$A\boldsymbol{x}=\boldsymbol{b}$（$\boldsymbol{b}\ne\boldsymbol{0}$），其对应的齐次方程组$A\boldsymbol{x}=\boldsymbol{0}$称为它的\textbf{导出组}。解向量具有两条性质：

\textbf{性质1}\quad 设$\boldsymbol{\eta}_1,\boldsymbol{\eta}_2$都是$A\boldsymbol{x}=\boldsymbol{b}$的解，则$\boldsymbol{\eta}_1-\boldsymbol{\eta}_2$是导出组$A\boldsymbol{x}=\boldsymbol{0}$的解。

\textbf{性质2}\quad 设$\boldsymbol{\eta}$是$A\boldsymbol{x}=\boldsymbol{b}$的解，$\boldsymbol{\xi}$是导出组$A\boldsymbol{x}=\boldsymbol{0}$的解，则$\boldsymbol{\eta}+\boldsymbol{\xi}$仍是$A\boldsymbol{x}=\boldsymbol{b}$的解。

由此得到解的\textbf{结构定理}：若$\boldsymbol{\eta}^{*}$是$A\boldsymbol{x}=\boldsymbol{b}$的一个\textbf{特解}，$\boldsymbol{\xi}_1,\dots,\boldsymbol{\xi}_{n-r}$是导出组的基础解系，则$A\boldsymbol{x}=\boldsymbol{b}$的通解为
\[
\boldsymbol{x}=\boldsymbol{\eta}^{*}+c_1\boldsymbol{\xi}_1+c_2\boldsymbol{\xi}_2+\dots+c_{n-r}\boldsymbol{\xi}_{n-r}
\qquad(c_1,\dots,c_{n-r}\in\mathbb{R}).
\]
即：\textbf{非齐次方程组的通解$=$它的一个特解$+$导出组的通解}。反过来，它的任一解都必可写成这种形式（由性质1，$\boldsymbol{\eta}-\boldsymbol{\eta}^{*}$是导出组的解）。

\keyformula{A\boldsymbol{x}=\boldsymbol{0}:\quad \dim S=n-R(A),\qquad \boldsymbol{x}=c_1\boldsymbol{\xi}_1+\dots+c_{n-r}\boldsymbol{\xi}_{n-r}}

\keyformula{A\boldsymbol{x}=\boldsymbol{b}:\quad \boldsymbol{x}=\boldsymbol{\eta}^{*}+c_1\boldsymbol{\xi}_1+\dots+c_{n-r}\boldsymbol{\xi}_{n-r}\quad(\boldsymbol{\eta}^{*}\ \text{为任一特解})}

\begin{table}[htbp]
\centering
\caption{$n$元线性方程组$A\boldsymbol{x}=\boldsymbol{b}$解的判定}
\begin{tabular}{@{}lll@{}}
\toprule
条件 & 结论 & 解的结构\\
\midrule
$R(A)\ne R(A,\boldsymbol{b})$ & 无解 & ——\\
$R(A)=R(A,\boldsymbol{b})=n$ & 有唯一解 & 只有特解$\boldsymbol{\eta}^{*}$\\
$R(A)=R(A,\boldsymbol{b})=r<n$ & 有无穷多解 & $\boldsymbol{\eta}^{*}+c_1\boldsymbol{\xi}_1+\dots+c_{n-r}\boldsymbol{\xi}_{n-r}$\\
\bottomrule
\end{tabular}
\end{table}

\corepoint{
（1）"无穷多解"时通解中\textbf{任意常数的个数}恒为$n-R(A)$，与$m$无关；写答案时漏掉常数或个数写错都是硬伤。\\
（2）特解只需\textbf{一个}，通常取自由未知数全为$0$时得到的那个解；特解不同，通解的形式不同但解集相同。\\
（3）验证答案：把通解代回$A\boldsymbol{x}=\boldsymbol{b}$；再检查基础解系是否满足$A\boldsymbol{\xi}_i=\boldsymbol{0}$且彼此线性无关。\\
（4）非齐次方程组的解集\textbf{不是}向量空间，"两个解之和"一般不再是解，故不能对非齐次组套用"解空间""基"的说法。
}


% ---------------- 第 5 章 相似矩阵及二次型 ----------------
% 教材 §1 向量的内积、长度及正交性 / §2 方阵的特征值与特征向量 / §3 相似矩阵
%      §4 对称矩阵的对角化 / §5 二次型及其标准形 / §6 用配方法化二次型成标准形 / §7 正定二次型
\chapter{相似矩阵及二次型}
% 第5章 相似矩阵及二次型（同济《线性代数》第七版 第5章）§1–§4
% §1 向量的内积、长度及正交性；§2 方阵的特征值与特征向量；§3 相似矩阵；§4 对称矩阵的对角化
% §5–§7（二次型）见 10.二次型与二次曲面.tex。正文不写例题与习题。

\section{向量的内积、长度及正交性}

\subsection{内积及其性质}
\begin{definition}[内积]
设有$n$维向量
\[
\boldsymbol{x}=\begin{pmatrix}x_1\\x_2\\\vdots\\x_n\end{pmatrix},
\qquad
\boldsymbol{y}=\begin{pmatrix}y_1\\y_2\\\vdots\\y_n\end{pmatrix},
\]
令
\[
(\boldsymbol{x},\boldsymbol{y})=x_1y_1+x_2y_2+\dots+x_ny_n=\sum_{i=1}^{n}x_iy_i,
\]
称$(\boldsymbol{x},\boldsymbol{y})$为向量$\boldsymbol{x}$与$\boldsymbol{y}$的\textbf{内积}。
\end{definition}

内积是两向量之间的运算，其结果是一个实数。用矩阵记号表示，当$\boldsymbol{x},\boldsymbol{y}$都是列向量时，
\[
(\boldsymbol{x},\boldsymbol{y})=\boldsymbol{x}^{T}\boldsymbol{y}=\boldsymbol{y}^{T}\boldsymbol{x}.
\]
$n$维向量的内积是解析几何中数量积（点积）的推广，但$n>3$时没有直观的"长度""夹角"，因此反过来用内积去定义长度和夹角。

内积满足下列性质（$\boldsymbol{x},\boldsymbol{y},\boldsymbol{z}$为$n$维向量，$\lambda\in\mathbb{R}$）：
\begin{enumerate}
    \item $(\boldsymbol{x},\boldsymbol{y})=(\boldsymbol{y},\boldsymbol{x})$（对称性）；
    \item $(\lambda\boldsymbol{x},\boldsymbol{y})=\lambda(\boldsymbol{x},\boldsymbol{y})$（关于第一个变量的齐次性）；
    \item $(\boldsymbol{x}+\boldsymbol{y},\boldsymbol{z})=(\boldsymbol{x},\boldsymbol{z})+(\boldsymbol{y},\boldsymbol{z})$（关于第一个变量的可加性）；
    \item $(\boldsymbol{x},\boldsymbol{x})\ge0$，且$(\boldsymbol{x},\boldsymbol{x})=0\iff\boldsymbol{x}=\boldsymbol{0}$（正定性）；
    \item $(\boldsymbol{x},\lambda\boldsymbol{y}+\mu\boldsymbol{z})=\lambda(\boldsymbol{x},\boldsymbol{y})+\mu(\boldsymbol{x},\boldsymbol{z})$（关于第二个变量的线性性）。
\end{enumerate}

\subsection{夹角与正交}
\begin{definition}[长度、单位向量]
令
\[
\|\boldsymbol{x}\|=\sqrt{(\boldsymbol{x},\boldsymbol{x})}=\sqrt{x_1^2+x_2^2+\dots+x_n^2},
\]
称$\|\boldsymbol{x}\|$为$n$维向量$\boldsymbol{x}$的\textbf{长度}（或范数）。长度为$1$的向量称为\textbf{单位向量}；若$\boldsymbol{x}\ne\boldsymbol{0}$，取$\boldsymbol{x}/\|\boldsymbol{x}\|$即得单位向量，这一过程称为把向量\textbf{单位化}。
\end{definition}

向量长度具有与解析几何中向量长度相同的三条基本性质：
\begin{enumerate}
    \item 非负性：$\|\boldsymbol{x}\|\ge0$，且$\|\boldsymbol{x}\|=0\iff\boldsymbol{x}=\boldsymbol{0}$；
    \item 齐次性：$\|\lambda\boldsymbol{x}\|=|\lambda|\,\|\boldsymbol{x}\|$；
    \item 三角不等式：$\|\boldsymbol{x}+\boldsymbol{y}\|\le\|\boldsymbol{x}\|+\|\boldsymbol{y}\|$。
\end{enumerate}

\begin{theorem}[施瓦茨不等式]
对任意$n$维向量$\boldsymbol{x},\boldsymbol{y}$，有
\[
(\boldsymbol{x},\boldsymbol{y})^{2}\le(\boldsymbol{x},\boldsymbol{x})(\boldsymbol{y},\boldsymbol{y}),
\qquad\text{即}\qquad
|(\boldsymbol{x},\boldsymbol{y})|\le\|\boldsymbol{x}\|\,\|\boldsymbol{y}\|.
\]
\end{theorem}

由施瓦茨不等式，当$\boldsymbol{x}\ne\boldsymbol{0},\boldsymbol{y}\ne\boldsymbol{0}$时$\left|\dfrac{(\boldsymbol{x},\boldsymbol{y})}{\|\boldsymbol{x}\|\|\boldsymbol{y}\|}\right|\le1$，故可定义夹角：

\begin{definition}[夹角、正交]
当$\boldsymbol{x}\ne\boldsymbol{0},\boldsymbol{y}\ne\boldsymbol{0}$时，令
\[
\theta=\arccos\frac{(\boldsymbol{x},\boldsymbol{y})}{\|\boldsymbol{x}\|\,\|\boldsymbol{y}\|}
\qquad(0\le\theta\le\pi),
\]
称$\theta$为$n$维向量$\boldsymbol{x}$与$\boldsymbol{y}$的\textbf{夹角}。当$(\boldsymbol{x},\boldsymbol{y})=0$时，称向量$\boldsymbol{x}$与$\boldsymbol{y}$\textbf{正交}，记作$\boldsymbol{x}\perp\boldsymbol{y}$。
\end{definition}

零向量与任何向量都正交。由两两正交的非零向量构成的向量组称为\textbf{正交向量组}；由单位向量构成的正交向量组称为\textbf{规范正交向量组}（标准正交向量组）。

\begin{theorem}
若$n$维向量$\boldsymbol{\alpha}_1,\boldsymbol{\alpha}_2,\dots,\boldsymbol{\alpha}_r$是一组两两正交的非零向量，则它们线性无关。
\end{theorem}

\textbf{证明}\quad 设有$\lambda_1\boldsymbol{\alpha}_1+\dots+\lambda_r\boldsymbol{\alpha}_r=\boldsymbol{0}$。用$\boldsymbol{\alpha}_i$与两端作内积，由两两正交得$\lambda_i(\boldsymbol{\alpha}_i,\boldsymbol{\alpha}_i)=0$，而$(\boldsymbol{\alpha}_i,\boldsymbol{\alpha}_i)=\|\boldsymbol{\alpha}_i\|^{2}>0$，故$\lambda_i=0$（$i=1,\dots,r$），所以线性无关。\hfill$\square$

推论：$n$维向量空间中两两正交的非零向量不多于$n$个。

\subsection{规范正交基}
\begin{definition}[规范正交基]
设$\boldsymbol{\alpha}_1,\dots,\boldsymbol{\alpha}_r$是向量空间$V$（$V\subseteq\mathbb{R}^{n}$）的一个基，若它们两两正交且都是单位向量，则称$\boldsymbol{\alpha}_1,\dots,\boldsymbol{\alpha}_r$是$V$的一个\textbf{规范正交基}（标准正交基）。
\end{definition}

$\boldsymbol{e}_1,\dots,\boldsymbol{e}_n$就是$\mathbb{R}^{n}$的一个规范正交基。若$\boldsymbol{e}_1,\dots,\boldsymbol{e}_n$是$V$的规范正交基，则$V$中任一向量$\boldsymbol{x}$可表示为$\boldsymbol{x}=\lambda_1\boldsymbol{e}_1+\dots+\lambda_n\boldsymbol{e}_n$，两端与$\boldsymbol{e}_i$作内积得

\keyformula{
\lambda_i=(\boldsymbol{x},\boldsymbol{e}_i),\qquad
\boldsymbol{x}=\sum_{i=1}^{n}(\boldsymbol{x},\boldsymbol{e}_i)\,\boldsymbol{e}_i .
}

即在规范正交基下，坐标就是内积，不必解方程组——这是取规范正交基的最大好处。

\subsubsection{施密特（Schmidt）正交化方法}
设$\boldsymbol{\alpha}_1,\dots,\boldsymbol{\alpha}_r$线性无关，要由此构造与它等价的规范正交组，分两步：

\textbf{(1) 正交化}\quad 取
\[
\begin{aligned}
\boldsymbol{\beta}_1&=\boldsymbol{\alpha}_1,\\
\boldsymbol{\beta}_2&=\boldsymbol{\alpha}_2-\frac{(\boldsymbol{\alpha}_2,\boldsymbol{\beta}_1)}{(\boldsymbol{\beta}_1,\boldsymbol{\beta}_1)}\boldsymbol{\beta}_1,\\
\boldsymbol{\beta}_3&=\boldsymbol{\alpha}_3-\frac{(\boldsymbol{\alpha}_3,\boldsymbol{\beta}_1)}{(\boldsymbol{\beta}_1,\boldsymbol{\beta}_1)}\boldsymbol{\beta}_1
-\frac{(\boldsymbol{\alpha}_3,\boldsymbol{\beta}_2)}{(\boldsymbol{\beta}_2,\boldsymbol{\beta}_2)}\boldsymbol{\beta}_2,\\
&\ \ \vdots\\
\boldsymbol{\beta}_r&=\boldsymbol{\alpha}_r-\sum_{j=1}^{r-1}\frac{(\boldsymbol{\alpha}_r,\boldsymbol{\beta}_j)}{(\boldsymbol{\beta}_j,\boldsymbol{\beta}_j)}\boldsymbol{\beta}_j .
\end{aligned}
\]

\textbf{(2) 单位化}\quad 令
\[
\boldsymbol{e}_i=\frac{\boldsymbol{\beta}_i}{\|\boldsymbol{\beta}_i\|}\qquad(i=1,2,\dots,r).
\]
则$\boldsymbol{e}_1,\dots,\boldsymbol{e}_r$就是与$\boldsymbol{\alpha}_1,\dots,\boldsymbol{\alpha}_r$等价的规范正交组。

该方法的优点是：不仅整体等价，而且对每个$i$都有
\[
\operatorname{Span}\{\boldsymbol{\beta}_1,\dots,\boldsymbol{\beta}_i\}=\operatorname{Span}\{\boldsymbol{\alpha}_1,\dots,\boldsymbol{\alpha}_i\},
\qquad i=1,2,\dots,r .
\]

\subsection{正交矩阵与正交变换}
\begin{definition}[正交矩阵]
如果$n$阶实矩阵$A$满足
\[
A^{T}A=E\quad(\text{即}\ A^{-1}=A^{T}),
\]
则称$A$为\textbf{正交矩阵}。
\end{definition}

用$A$的列向量表示，$A^{T}A=E$即$a_{ij}$满足$(\boldsymbol{\alpha}_i,\boldsymbol{\alpha}_j)=\delta_{ij}$。因此：\textbf{$A$为正交矩阵$\iff$$A$的列向量都是单位向量且两两正交}，即列向量组是$\mathbb{R}^{n}$的规范正交基；由于$A^{T}A=E$与$AA^{T}=E$等价，结论对行向量同样成立。

$n$阶正交矩阵$A$有以下性质：
\begin{enumerate}
    \item $A$可逆，且$A^{-1}=A^{T}$；
    \item $A^{T}$（从而$A^{-1}$）也是正交矩阵；
    \item 若$A,B$都是$n$阶正交矩阵，则$AB$也是正交矩阵；
    \item $|A|=1$或$|A|=-1$；
    \item 对任意$n$维列向量$\boldsymbol{x},\boldsymbol{y}$，有$\|A\boldsymbol{x}\|=\|\boldsymbol{x}\|$且$(A\boldsymbol{x},A\boldsymbol{y})=(\boldsymbol{x},\boldsymbol{y})$（保长度、保内积）。
\end{enumerate}

\begin{definition}[正交变换]
若$P$为正交矩阵，则线性变换$\boldsymbol{x}=P\boldsymbol{y}$称为\textbf{正交变换}。
\end{definition}

由性质5，正交变换保持向量的长度（从而保持线段长度、三角形形状）和内积不变。$\mathbb{R}^{n}$中向量在两组规范正交基之间的坐标变换就是正交变换。

\keyformula{(\boldsymbol{x},\boldsymbol{y})=\boldsymbol{x}^{T}\boldsymbol{y},\qquad \|\boldsymbol{x}\|=\sqrt{(\boldsymbol{x},\boldsymbol{x})},\qquad \boldsymbol{x}\perp\boldsymbol{y}\iff\boldsymbol{x}^{T}\boldsymbol{y}=0}

\keyformula{A^{T}A=E\iff A^{-1}=A^{T}\iff A\ \text{的列（行）向量组为}\ \mathbb{R}^{n}\ \text{的规范正交基}}

\corepoint{
（1）正交化公式的分母是$(\boldsymbol{\beta}_j,\boldsymbol{\beta}_j)=\|\boldsymbol{\beta}_j\|^{2}$，分子是$(\boldsymbol{\alpha}_i,\boldsymbol{\beta}_j)$，减的是$\boldsymbol{\beta}_j$（已正交化的向量），不能写成$\boldsymbol{\alpha}_j$。\\
（2）必须"先正交化、后单位化"；顺序颠倒会破坏正交性。\\
（3）施密特正交化只在同一线性无关组内部使用；若两组向量本来就要合并正交化，须按给定顺序整体处理。\\
（4）内积对第二个变量是线性的，对第一个变量也是线性的（实向量情形），但写成$(k\boldsymbol{x},\boldsymbol{y})=k(\boldsymbol{x},\boldsymbol{y})$时注意别把$k$漏在括号外。
}

\section{方阵的特征值与特征向量}

\subsection{特征值与特征向量的概念}
\begin{definition}[特征值、特征向量]
设$A$是$n$阶矩阵，如果数$\lambda$和$n$维\textbf{非零}列向量$\boldsymbol{x}$使关系式
\[
A\boldsymbol{x}=\lambda\boldsymbol{x}
\]
成立，那么这样的数$\lambda$称为矩阵$A$的\textbf{特征值}，非零向量$\boldsymbol{x}$称为$A$的对应于特征值$\lambda$的\textbf{特征向量}。
\end{definition}

$A\boldsymbol{x}=\lambda\boldsymbol{x}$可写成
\[
(\lambda E-A)\boldsymbol{x}=\boldsymbol{0},
\]
这是$n$个未知数$n$个方程的齐次线性方程组，它有非零解的充分必要条件是系数行列式
\[
|\lambda E-A|=0,
\]
即
\[
\begin{vmatrix}
\lambda-a_{11}&-a_{12}&\dots&-a_{1n}\\
-a_{21}&\lambda-a_{22}&\dots&-a_{2n}\\
\vdots&\vdots&&\vdots\\
-a_{n1}&-a_{n2}&\dots&\lambda-a_{nn}
\end{vmatrix}=0 .
\]
这是以$\lambda$为未知数的一元$n$次方程，称为矩阵$A$的\textbf{特征方程}；其左端$|\lambda E-A|$是$\lambda$的$n$次多项式，称为$A$的\textbf{特征多项式}。

于是：
\begin{enumerate}
    \item $A$的特征值就是特征方程的根；$n$阶矩阵在复数范围内有$n$个特征值（重根按重数计）。
    \item $A$的属于特征值$\lambda_i$的特征向量，就是齐次方程组$(\lambda_iE-A)\boldsymbol{x}=\boldsymbol{0}$的\textbf{非零}解向量。
    \item 该方程组的解空间$N(\lambda_iE-A)$称为$A$的关于$\lambda_i$的\textbf{特征子空间}；其中除零向量外，其余向量全是属于$\lambda_i$的特征向量。
\end{enumerate}

\subsection{特征值的性质}
设$n$阶矩阵$A=(a_{ij})$的特征值为$\lambda_1,\lambda_2,\dots,\lambda_n$，则
\begin{enumerate}
    \item $\lambda_1+\lambda_2+\dots+\lambda_n=a_{11}+a_{22}+\dots+a_{nn}=\operatorname{tr}(A)$（$A$的迹）；
    \item $\lambda_1\lambda_2\cdots\lambda_n=|A|$。
\end{enumerate}
由此可得：$A$可逆$\iff$$A$的$n$个特征值全不为零；$|A|=0$$\iff$$0$是$A$的特征值。

特征值还有下列常用性质（$\lambda$为$A$的特征值，$\boldsymbol{x}$为对应特征向量）：
\begin{enumerate}
    \item $k\lambda$是$kA$的特征值；$\lambda^{m}$是$A^{m}$的特征值；
    \item 若$f(x)=a_mx^{m}+\dots+a_1x+a_0$是多项式，则$f(\lambda)$是$f(A)=a_mA^{m}+\dots+a_1A+a_0E$的特征值，且$\boldsymbol{x}$仍是对应特征向量；
    \item 当$A$可逆时，$\lambda\ne0$，$\lambda^{-1}$是$A^{-1}$的特征值，$\dfrac{|A|}{\lambda}$是$A^{*}$（伴随矩阵）的特征值；
    \item $A^{T}$与$A$有相同的特征值（但特征向量一般不同）。
\end{enumerate}

\begin{theorem}[不同特征值的特征向量线性无关]
设$\lambda_1,\lambda_2,\dots,\lambda_m$是$n$阶矩阵$A$的$m$个\textbf{互不相等}的特征值，$\boldsymbol{p}_1,\dots,\boldsymbol{p}_m$依次是与它们对应的特征向量，则$\boldsymbol{p}_1,\dots,\boldsymbol{p}_m$线性无关。
\end{theorem}

\begin{corollary}
设$\lambda_1\ne\lambda_2$是$A$的两个不同特征值，$\boldsymbol{p}_{1},\dots,\boldsymbol{p}_{r}$与$\boldsymbol{q}_{1},\dots,\boldsymbol{q}_{s}$分别是属于$\lambda_1$与$\lambda_2$的线性无关特征向量，则$\boldsymbol{p}_1,\dots,\boldsymbol{p}_r,\boldsymbol{q}_1,\dots,\boldsymbol{q}_s$线性无关。
\end{corollary}

注意：属于同一特征值的特征向量的任意非零线性组合仍是该特征值的特征向量（因为特征子空间是向量空间）；但属于两个不同特征值的特征向量之和\textbf{不再是}特征向量。

\subsection{求特征值与特征向量的步骤}
\begin{enumerate}
    \item 计算特征多项式$|\lambda E-A|$，解特征方程$|\lambda E-A|=0$，求出全部特征值$\lambda_1,\dots,\lambda_n$（重根按重数记）。
    \item 对每个特征值$\lambda_i$，解齐次方程组$(\lambda_iE-A)\boldsymbol{x}=\boldsymbol{0}$，求其基础解系$\boldsymbol{\xi}_1,\dots,\boldsymbol{\xi}_{t}$。
    \item 则属于$\lambda_i$的全部特征向量为$k_1\boldsymbol{\xi}_1+\dots+k_t\boldsymbol{\xi}_t$（$k_1,\dots,k_t$不全为零）。
\end{enumerate}

\corepoint{
（1）特征向量必须\textbf{非零}，答"全部特征向量"时要写"不全为零"的线性组合，不能只写基础解系。\\
（2）特征方程用$|\lambda E-A|=0$或$|A-\lambda E|=0$都可以（两者只差$(-1)^{n}$），但\textbf{代入求特征向量时必须用同一个形式}，否则符号出错。\\
（3）解$(\lambda_iE-A)\boldsymbol{x}=\boldsymbol{0}$时，若$\lambda_i$是重根，得到的线性无关特征向量个数可能少于重数（这正是"不可对角化"的原因）。\\
（4）验算：$\sum\lambda_i=\operatorname{tr}(A)$、$\prod\lambda_i=|A|$，两式能查出大部分计算错误；再把$\boldsymbol{x}$代回$A\boldsymbol{x}=\lambda\boldsymbol{x}$验证。
}

\keyformula{A\boldsymbol{x}=\lambda\boldsymbol{x}\ (\boldsymbol{x}\ne\boldsymbol{0})\quad\Longleftrightarrow\quad (\lambda E-A)\boldsymbol{x}=\boldsymbol{0}\ \text{有非零解}\quad\Longleftrightarrow\quad |\lambda E-A|=0}

\keyformula{\sum_{i=1}^{n}\lambda_i=\operatorname{tr}(A),\qquad \prod_{i=1}^{n}\lambda_i=|A|,\qquad f(\lambda)\ \text{是}\ f(A)\ \text{的特征值},\qquad \lambda^{-1}\ \text{是}\ A^{-1}\ \text{的特征值}}

\section{相似矩阵}

\subsection{相似矩阵的概念与性质}
\begin{definition}[相似]
设$A,B$都是$n$阶矩阵，若有可逆矩阵$P$，使
\[
P^{-1}AP=B,
\]
则称$B$是$A$的\textbf{相似矩阵}，或说矩阵$A$与$B$\textbf{相似}；对$A$进行运算$P^{-1}AP$称为对$A$进行\textbf{相似变换}，可逆矩阵$P$称为把$A$变成$B$的\textbf{相似变换矩阵}。
\end{definition}

矩阵的相似关系具有：
\begin{enumerate}
    \item 自反性：$A$与$A$相似；
    \item 对称性：$A$与$B$相似，则$B$与$A$相似；
    \item 传递性：$A$与$B$相似、$B$与$C$相似，则$A$与$C$相似。
\end{enumerate}

\begin{theorem}[相似矩阵的特征多项式相同]
若$n$阶矩阵$A$与$B$相似，则$A$与$B$的特征多项式相同，从而$A$与$B$的特征值亦相同。
\end{theorem}

\textbf{证明}\quad 由$P^{-1}AP=B$得
\[
|B-\lambda E|=|P^{-1}AP-P^{-1}(\lambda E)P|=|P^{-1}(A-\lambda E)P|
=|P|^{-1}|A-\lambda E||P|=|A-\lambda E| .
\]
\hfill$\square$

由此还可得到：相似矩阵有相同的迹、相同的行列式、相同的秩。反之不成立——特征值相同（甚至$|A|=|B|$、$\operatorname{tr}A=\operatorname{tr}B$）的两个矩阵未必相似。

\begin{corollary}
若$n$阶矩阵$A$与对角矩阵
\[
\Lambda=\begin{pmatrix}\lambda_1&&&\\&\lambda_2&&\\&&\ddots&\\&&&\lambda_n\end{pmatrix}
\]
相似，则$\lambda_1,\lambda_2,\dots,\lambda_n$即是$A$的$n$个特征值。
\end{corollary}

若$A=P\Lambda P^{-1}$，则$A^{k}=P\Lambda^{k}P^{-1}$，且对多项式$f$有$f(A)=Pf(\Lambda)P^{-1}$，而对角矩阵的幂与多项式极易计算。这正是要求"把矩阵化为对角矩阵"的原因。

\subsection{方阵可对角化的条件}
\begin{theorem}[可对角化的充要条件]
$n$阶矩阵$A$与对角矩阵相似（即$A$能对角化）的充分必要条件是$A$有$n$个线性无关的特征向量。

进一步，$P^{-1}AP=\Lambda$为对角矩阵$\iff$$P$的$n$个列向量是$A$的$n$个线性无关的特征向量，且$\Lambda$的主对角线元素依次是这些特征向量所对应的特征值。
\end{theorem}

\textbf{证明要点}\quad 设$P=(\boldsymbol{p}_1,\dots,\boldsymbol{p}_n)$可逆，则$P^{-1}AP=\Lambda\iff AP=P\Lambda\iff(A\boldsymbol{p}_1,\dots,A\boldsymbol{p}_n)=(\lambda_1\boldsymbol{p}_1,\dots,\lambda_n\boldsymbol{p}_n)\iff A\boldsymbol{p}_i=\lambda_i\boldsymbol{p}_i$。由于$P$可逆$\iff$$\boldsymbol{p}_1,\dots,\boldsymbol{p}_n$线性无关，即得结论。

\begin{corollary}
若$n$阶矩阵$A$的$n$个特征值互不相等，则$A$与对角矩阵相似。
\end{corollary}

反之不成立：有重特征值时也可能可对角化，只是\textbf{不一定}。

\subsection{几何重数与代数重数}
\begin{definition}[代数重数、几何重数]
设$\lambda_0$是$n$阶矩阵$A$的特征值，若$\lambda_0$是特征方程的$m$重根，称$m$为$\lambda_0$的\textbf{代数重数}；称$\lambda_0$对应的特征子空间（即$(\lambda_0E-A)\boldsymbol{x}=\boldsymbol{0}$的解空间）的维数为$\lambda_0$的\textbf{几何重数}，它等于属于$\lambda_0$的线性无关特征向量的最大个数，即$n-R(\lambda_0E-A)$。
\end{definition}

\begin{theorem}
矩阵的任一特征值的几何重数不超过它的代数重数，即
\[
n-R(\lambda_0E-A)\le m .
\]
\end{theorem}

\begin{theorem}[可对角化的重数判据]
$n$阶矩阵$A$与对角矩阵相似的充分必要条件是：$A$的每一个特征值的几何重数都等于它的代数重数，即对每个特征值$\lambda_i$有
\[
n-R(\lambda_iE-A)=m_i .
\]
等价地：$\sum_i\left[n-R(\lambda_iE-A)\right]=n$，即$A$共有$n$个线性无关的特征向量。
\end{theorem}

于是判断能否对角化有两条路：(a) 数一数线性无关特征向量是否够$n$个；(b) 比较每个特征值的两个重数。

\subsection{对角化的步骤}
\begin{enumerate}
    \item 求$A$的全部特征值$\lambda_1,\dots,\lambda_s$（互异）及其代数重数$m_1,\dots,m_s$。
    \item 对每个$\lambda_i$解$(\lambda_iE-A)\boldsymbol{x}=\boldsymbol{0}$，求基础解系，得$t_i=n-R(\lambda_iE-A)$个线性无关特征向量。
    \item 若$\sum t_i<n$，则$A$不可对角化（重特征值处向量不够）；若$\sum t_i=n$，令
    \[
    P=(\boldsymbol{p}_1,\boldsymbol{p}_2,\dots,\boldsymbol{p}_n),\qquad
    \Lambda=\operatorname{diag}(\lambda_1',\lambda_2',\dots,\lambda_n'),
    \]
    则$P^{-1}AP=\Lambda$。
\end{enumerate}

\corepoint{
（1）$P$中列向量的\textbf{排列顺序}必须与$\Lambda$中对角元的顺序一一对应；顺序错则等式不成立（这是阅卷最常见的扣分点）。\\
（2）$P$不唯一：同一特征值的特征向量可换成任意非零倍数或其他等价组，$\Lambda$中同特征值的对角元也可交换。\\
（3）"特征值互不相同$\Rightarrow$可对角化"成立，但"有重根$\Rightarrow$不可对角化"不成立。\\
（4）注意相似与合同、等价是三种不同的关系：相似是$P^{-1}AP$，合同是$C^{T}AC$，等价是$PAQ$。
}

\keyformula{B=P^{-1}AP\ \Rightarrow\ |\lambda E-B|=|\lambda E-A|,\quad \operatorname{tr}B=\operatorname{tr}A,\quad |B|=|A|}

\keyformula{A\ \text{可对角化}\iff A\ \text{有}\ n\ \text{个线性无关的特征向量}\iff n-R(\lambda_iE-A)=m_i\ (\forall i)}

\keyformula{P^{-1}AP=\Lambda,\quad P=(\boldsymbol{p}_1,\dots,\boldsymbol{p}_n),\quad A\boldsymbol{p}_i=\lambda_i\boldsymbol{p}_i}

\section{对称矩阵的对角化}

实对称矩阵是一类一定能对角化、而且可以用\textbf{正交}矩阵对角化的矩阵，这是二次型理论的基础。

\subsection{实对称矩阵特征值与特征向量的性质}
\begin{theorem}[性质1：特征值为实数]
实对称矩阵$A$的特征值都是实数。
\end{theorem}

证明思路：设$\lambda$为特征值、$\boldsymbol{x}$为对应复特征向量，取共轭转置比较$A\bar{\boldsymbol{x}}=\bar\lambda\bar{\boldsymbol{x}}$与$\boldsymbol{x}^{T}A=\lambda\boldsymbol{x}^{T}$，得$(\lambda-\bar\lambda)\bar{\boldsymbol{x}}^{T}\boldsymbol{x}=0$，而$\bar{\boldsymbol{x}}^{T}\boldsymbol{x}>0$，故$\lambda=\bar\lambda$。此时$(\lambda E-A)\boldsymbol{x}=\boldsymbol{0}$是实系数方程组，特征向量可取实向量。

\begin{theorem}[性质2：不同特征值的特征向量正交]
设$\lambda_1,\lambda_2$是实对称矩阵$A$的两个特征值，$\boldsymbol{p}_1,\boldsymbol{p}_2$是对应的特征向量。若$\lambda_1\ne\lambda_2$，则$\boldsymbol{p}_1$与$\boldsymbol{p}_2$正交。
\end{theorem}

\textbf{证明}\quad $\lambda_1\boldsymbol{p}_1^{T}\boldsymbol{p}_2=(A\boldsymbol{p}_1)^{T}\boldsymbol{p}_2=\boldsymbol{p}_1^{T}A^{T}\boldsymbol{p}_2=\boldsymbol{p}_1^{T}A\boldsymbol{p}_2=\lambda_2\boldsymbol{p}_1^{T}\boldsymbol{p}_2$，故$(\lambda_1-\lambda_2)\boldsymbol{p}_1^{T}\boldsymbol{p}_2=0$，而$\lambda_1\ne\lambda_2$，得$\boldsymbol{p}_1^{T}\boldsymbol{p}_2=0$。\hfill$\square$

\begin{theorem}[性质3：重根的几何重数等于代数重数]
设$A$为$n$阶实对称矩阵，$\lambda_0$是$A$的特征方程的$m$重根，则矩阵$\lambda_0E-A$的秩$R(\lambda_0E-A)=n-m$，从而对应于$\lambda_0$恰有$m$个线性无关的特征向量。
\end{theorem}

\begin{theorem}[实对称矩阵的正交对角化]
设$A$为$n$阶实对称矩阵，则必有正交矩阵$P$，使
\[
P^{-1}AP=P^{T}AP=\Lambda=
\begin{pmatrix}\lambda_1&&&\\&\lambda_2&&\\&&\ddots&\\&&&\lambda_n\end{pmatrix},
\]
其中$\lambda_1,\dots,\lambda_n$是$A$的$n$个特征值。
\end{theorem}

即：实对称矩阵一定可以正交相似对角化，且相似变换矩阵可取为正交矩阵（此时$P^{-1}=P^{T}$，求逆变成转置）。

\subsection{用正交矩阵对角化的步骤}
\begin{enumerate}
    \item 求出$A$的全部互不相等的特征值$\lambda_1,\dots,\lambda_s$，重数依次为$m_1,\dots,m_s$（$\sum m_i=n$）。
    \item 对每个重特征值$\lambda_i$，解$(\lambda_iE-A)\boldsymbol{x}=\boldsymbol{0}$，求基础解系（恰含$m_i$个向量），再把这$m_i$个向量\textbf{施密特正交化并单位化}，得$m_i$个两两正交的单位特征向量。
    \item 因不同特征值的特征向量已自动正交，把这$n$个两两正交的单位特征向量构成正交矩阵$P$，便有
    \[
    P^{-1}AP=P^{T}AP=\Lambda,
    \]
    其中$\Lambda$对角元的排列次序与$P$中列向量的排列次序相对应。
\end{enumerate}

\corepoint{
（1）\textbf{同一重特征值}对应的特征向量只保证线性无关，不保证正交，必须自己做施密特正交化；\textbf{不同特征值}的特征向量已正交，不能再混在一起正交化（会破坏原有的方向）。\\
（2）最后必须\textbf{单位化}，否则$P$只是可逆矩阵而不是正交矩阵，"$P^{-1}AP$"就不再等于"$P^{T}AP$"。\\
（3）$A$是实对称矩阵时一定成功；若$A$不是对称矩阵，即使特征值全是实数，也不能保证存在正交$P$。\\
（4）验算：$P^{T}P=E$、$P^{T}AP=\Lambda$两条都要检查。
}

\keyformula{A^{T}=A\ \Longrightarrow\ \exists\,P\ (P^{T}P=E):\ P^{-1}AP=P^{T}AP=\operatorname{diag}(\lambda_1,\dots,\lambda_n)}

\keyformula{\lambda_i\in\mathbb{R};\qquad \lambda_i\ne\lambda_j\ \Rightarrow\ \boldsymbol{p}_i\perp\boldsymbol{p}_j;\qquad R(\lambda_0E-A)=n-m_0}

% 第5章 相似矩阵及二次型（同济《线性代数》第七版 第5章）§5–§7 + 拓展（二次曲面）
% §5 二次型及其标准形；§6 用配方法化二次型成标准形；§7 正定二次型。
% 二次曲面为教材之外的拓展节（对应第8章空间解析几何相关内容），保留作几何背景。
% §1–§4 见 8.特征值&特征向量&相似矩阵.tex。正文不写例题与习题。

\section{二次型及其标准形}

\subsection{二次型及其矩阵}
\begin{definition}[二次型]
含有$n$个变量$x_1,x_2,\dots,x_n$的二次齐次函数
\[
\begin{aligned}
f(x_1,x_2,\dots,x_n)={}&a_{11}x_1^{2}+a_{22}x_2^{2}+\dots+a_{nn}x_n^{2}\\
&+2a_{12}x_1x_2+2a_{13}x_1x_3+\dots+2a_{n-1,n}x_{n-1}x_n
\end{aligned}
\]
称为\textbf{二次型}。当系数$a_{ij}$全为实数时称为\textbf{实二次型}，含有复数系数时称为复二次型。本书只讨论实二次型。
\end{definition}

取$a_{ij}=a_{ji}$（$i<j$），则$2a_{ij}x_ix_j=a_{ij}x_ix_j+a_{ji}x_jx_i$，于是
\[
f=\sum_{i=1}^{n}\sum_{j=1}^{n}a_{ij}x_ix_j .
\]
记$\boldsymbol{x}=(x_1,x_2,\dots,x_n)^{T}$，
\[
A=\begin{pmatrix}
a_{11}&a_{12}&\dots&a_{1n}\\
a_{21}&a_{22}&\dots&a_{2n}\\
\vdots&\vdots&\ddots&\vdots\\
a_{n1}&a_{n2}&\dots&a_{nn}
\end{pmatrix}
\quad(a_{ij}=a_{ji}),
\]
则二次型可用矩阵记号写成
\[
f=\boldsymbol{x}^{T}A\boldsymbol{x},
\]
其中$A$为\textbf{对称矩阵}，称为二次型$f$的\textbf{矩阵}；$A$的秩称为二次型$f$的\textbf{秩}。

任给一个二次型，就唯一确定一个对称矩阵；任给一个对称矩阵，也可唯一确定一个二次型。因此二次型与对称矩阵之间存在\textbf{一一对应}关系，对称矩阵$A$也叫做二次型$f$的矩阵，$f$叫做对称矩阵$A$的二次型。

只含平方项的二次型称为\textbf{标准形}；形如
\[
f=y_1^{2}+\dots+y_p^{2}-y_{p+1}^{2}-\dots-y_{p+q}^{2}
\]
的标准形称为\textbf{规范形}。

\subsection{合同矩阵}
我们讨论的主要问题是：是否存在可逆线性变换$\boldsymbol{x}=C\boldsymbol{y}$（若存在，如何寻找），使二次型化为标准形。把$\boldsymbol{x}=C\boldsymbol{y}$代入$f=\boldsymbol{x}^{T}A\boldsymbol{x}$，得
\[
f=\boldsymbol{x}^{T}A\boldsymbol{x}=(C\boldsymbol{y})^{T}A(C\boldsymbol{y})
=\boldsymbol{y}^{T}(C^{T}AC)\boldsymbol{y}.
\]
记$B=C^{T}AC$，则$B$仍为对称矩阵，且$R(B)=R(A)$（$C$可逆），$B$就是变换后二次型的矩阵。

\begin{definition}[合同]
设$A$和$B$是$n$阶矩阵，若有可逆矩阵$C$，使
\[
B=C^{T}AC,
\]
则称矩阵$A$与$B$\textbf{合同}。
\end{definition}

合同关系具有自反性、对称性、传递性。若$A$为对称矩阵，则$B=C^{T}AC$也是对称矩阵，且$R(B)=R(A)$——即\textbf{经可逆线性变换，二次型的矩阵变为与之合同的矩阵，二次型的秩不变}。

于是"用可逆线性变换化二次型为标准形"等价于"求可逆矩阵$C$使$C^{T}AC$为对角矩阵"，这个问题称为把对称矩阵$A$\textbf{合同对角化}。

\subsection{用正交变换化二次型为标准形}
\begin{theorem}[正交变换化标准形]
任给二次型$f=\sum_{i,j=1}^{n}a_{ij}x_ix_j=\boldsymbol{x}^{T}A\boldsymbol{x}$（$A^{T}=A$），总有\textbf{正交变换}$\boldsymbol{x}=P\boldsymbol{y}$，使$f$化为标准形
\[
f=\lambda_1y_1^{2}+\lambda_2y_2^{2}+\dots+\lambda_ny_n^{2},
\]
其中$\lambda_1,\lambda_2,\dots,\lambda_n$是$A$的$n$个特征值。
\end{theorem}

这是因为实对称矩阵$A$必有正交矩阵$P$使$P^{-1}AP=P^{T}AP=\Lambda$（第4节定理），令$\boldsymbol{x}=P\boldsymbol{y}$即得。此时$P$的列向量是$A$的$n$个两两正交的单位特征向量。

\begin{corollary}
任给$n$元二次型$f=\boldsymbol{x}^{T}A\boldsymbol{x}$，总有可逆变换$\boldsymbol{x}=C\boldsymbol{z}$，使$f(C\boldsymbol{z})$为规范形。
\end{corollary}

证明思路：先由定理取正交变换化为$\lambda_1y_1^{2}+\dots+\lambda_ny_n^{2}$，设$f$的秩为$r$，则恰有$r$个$\lambda_i\ne0$；再令$y_i=\dfrac{z_i}{\sqrt{|\lambda_i|}}$（$i\le r$），$y_i=z_i$（$i>r$），即得系数为$\pm1$或$0$的规范形。

\subsubsection{正交变换法的步骤}
\begin{enumerate}
    \item 写出二次型的矩阵$A$（注意：交叉项$x_ix_j$的系数要\textbf{平分}到$a_{ij}$与$a_{ji}$两处）。
    \item 求$A$的特征值$\lambda_1,\dots,\lambda_n$及对应的两两正交的单位特征向量$\boldsymbol{p}_1,\dots,\boldsymbol{p}_n$（同一重特征值处需施密特正交化）。
    \item 令$P=(\boldsymbol{p}_1,\dots,\boldsymbol{p}_n)$，作正交变换$\boldsymbol{x}=P\boldsymbol{y}$，则
    \[
    f=\lambda_1y_1^{2}+\lambda_2y_2^{2}+\dots+\lambda_ny_n^{2}.
    \]
\end{enumerate}

正交变换的优点是保持几何形状不变（保长度、保内积），因此常用于判断二次曲面/二次曲线的类型。

\subsection*{附：用初等变换（合同变换）化标准形}
配方法用矩阵语言表达，就是对$A$作一列"协调一致"的初等行、列变换（$C^{T}AC$）。把单位矩阵放在$A$下面构成$2n\times n$矩阵$\begin{pmatrix}A\\E\end{pmatrix}$，每对$A$作一次列变换，就同时对$A$作一次相应的行变换（左乘、右乘同一个初等矩阵的转置），即
\[
\begin{pmatrix}A\\E\end{pmatrix}\longrightarrow\begin{pmatrix}C^{T}AC\\C\end{pmatrix}.
\]
当$C^{T}AC$化为对角矩阵时，下面的$E$就化成了所求的合同变换矩阵$C$。这就是\textbf{初等变换法}（也叫合同变换法），与配方法等价。

\keyformula{f=\boldsymbol{x}^{T}A\boldsymbol{x}\ (A^{T}=A),\qquad \boldsymbol{x}=C\boldsymbol{y}\ \Rightarrow\ f=\boldsymbol{y}^{T}(C^{T}AC)\boldsymbol{y},\qquad B=C^{T}AC\iff A\ \text{与}\ B\ \text{合同}}

\keyformula{\boldsymbol{x}=P\boldsymbol{y}\ (P^{T}P=E)\ \Rightarrow\ f=\lambda_1y_1^{2}+\dots+\lambda_ny_n^{2},\quad \lambda_i\ \text{为}\ A\ \text{的特征值}}

\corepoint{
（1）写矩阵时交叉项系数必须\textbf{除以2}：$f=2x_1x_2$对应$a_{12}=a_{21}=1$。这是二次型部分最高频的错误。\\
（2）正交变换得到的标准形系数\textbf{恰是特征值}，且系数顺序与$P$中列向量顺序一致；配方法得到的标准形系数则与所用变换有关，不唯一。\\
（3）正交变换化标准形\textbf{不改变}二次型的秩与正、负惯性指数；配方法同样不改变。\\
（4）正交变换下$C^{T}AC=C^{-1}AC$（因为$C^{T}=C^{-1}$），所以"合同"与"相似"在正交变换这一特殊情形下重合；一般情形两者不同。
}

\section{用配方法化二次型成标准形}

用正交变换化标准形需要求特征值、解特征向量，计算量大；若只要求"用可逆线性变换化标准形"，用拉格朗日配方法更简便，只是所得标准形不唯一。

\subsection{配方法的两种情形}
\textbf{情形一：二次型中含有某个变量的平方项。}\quad 不妨设含$x_1^{2}$，把含有$x_1$的所有项集中起来配成完全平方，剩下的是不含$x_1$的$n-1$元二次型，再对它重复上述步骤。具体地，若
\[
f=a_{11}x_1^{2}+2a_{12}x_1x_2+\dots+2a_{1n}x_1x_n+g(x_2,\dots,x_n),
\]
则
\[
f=\frac{1}{a_{11}}\bigl(a_{11}x_1+a_{12}x_2+\dots+a_{1n}x_n\bigr)^{2}
+g_1(x_2,\dots,x_n),
\]
其中$g_1$仍是二次型。令
\[
y_1=a_{11}x_1+a_{12}x_2+\dots+a_{1n}x_n,\quad y_i=x_i\ (i\ge2),
\]
这是可逆线性变换（系数矩阵为上三角且对角元非零），在此变换下$f$化为只含$y_2,\dots,y_n$的二次型，继续配方即可。

\textbf{情形二：二次型中不含任何平方项。}\quad 此时必有某个交叉项，例如$2a_{12}x_1x_2$（$a_{12}\ne0$）。先作可逆变换
\[
\begin{cases}
x_1=y_1+y_2,\\
x_2=y_1-y_2,\\
x_i=y_i\quad(i\ge3),
\end{cases}
\]
则$2a_{12}x_1x_2=2a_{12}(y_1^{2}-y_2^{2})$，出现了平方项，就化为情形一，继续配方。

\subsection{配方法的步骤}
\begin{enumerate}
    \item 若含平方项，就按上面的公式对第一个平方项变量配方，把含它的项全部吸收进完全平方，得到一个新的二次型（变量少一个）。
    \item 若不含平方项，先用$x_1=y_1+y_2,\ x_2=y_1-y_2$造出平方项。
    \item 重复以上两步，直到只剩平方项，得标准形$f=k_1z_1^{2}+\dots+k_rz_r^{2}$。
    \item 把各步的变换\textbf{复合}（依次代入），写出总的可逆线性变换$\boldsymbol{x}=C\boldsymbol{z}$及变换矩阵$C$。复合时注意"先作的变换在右边"。
\end{enumerate}

\corepoint{
（1）配方时必须把\textbf{所有}含该变量的项（包括交叉项）一次吸收干净，否则后面无法继续。\\
（2）情形二中若同时有多个交叉项，只需挑一个造平方项即可，不必全处理。\\
（3）最后求变换矩阵时，要把分步变换依次代入合成一个矩阵，不能只写最后一步的矩阵。\\
（4）标准形中系数为$0$的项可以省略，但写变换矩阵时对应的变量仍要保留（否则矩阵不满秩，不是可逆变换）。\\
（5）配方法不改变二次型的秩与惯性指数，可用这一点验算。
}

\keyformula{a_{11}x_1^{2}+2\sum_{j\ge2}a_{1j}x_1x_j=\frac{1}{a_{11}}\Bigl(\sum_{j\ge1}a_{1j}x_j\Bigr)^{2}-\frac{1}{a_{11}}\Bigl(\sum_{j\ge2}a_{1j}x_j\Bigr)^{2}}

\keyformula{x_1=y_1+y_2,\ x_2=y_1-y_2\ \Rightarrow\ 2a_{12}x_1x_2=2a_{12}y_1^{2}-2a_{12}y_2^{2}}

\keyformula{f\ \xrightarrow{\ \boldsymbol{x}=C\boldsymbol{y}\ }\ k_1y_1^{2}+\dots+k_ny_n^{2}\qquad(C\ \text{可逆})}

\section{正定二次型}

\subsection{惯性定理}
二次型的标准形显然不唯一，但标准形中所含\textbf{项数}（即二次型的秩）是确定的。不仅如此，在限定变换为实变换时，标准形中\textbf{正系数的个数}也是不变的：

\begin{theorem}[惯性定理]
设二次型$f=\boldsymbol{x}^{T}A\boldsymbol{x}$的秩为$r$，且有两个可逆变换$\boldsymbol{x}=C_1\boldsymbol{y}$及$\boldsymbol{x}=C_2\boldsymbol{z}$使
\[
f=k_1y_1^{2}+\dots+k_ry_r^{2}\quad(k_i\ne0),
\qquad
f=l_1z_1^{2}+\dots+l_rz_r^{2}\quad(l_i\ne0),
\]
则$k_1,\dots,k_r$中正数的个数与$l_1,\dots,l_r$中正数的个数相等（从而负数的个数也相等）。
\end{theorem}

二次型的标准形中正系数的个数称为二次型的\textbf{正惯性指数}，负系数的个数称为\textbf{负惯性指数}。若正惯性指数为$p$、秩为$r$，则负惯性指数为$r-p$，且$f$的\textbf{规范形}唯一确定：
\[
f=z_1^{2}+\dots+z_p^{2}-z_{p+1}^{2}-\dots-z_r^{2}.
\]

因此，两个$n$元实二次型合同$\iff$它们有相同的秩与相同的正惯性指数。

\subsection{正定二次型与正定矩阵}
\begin{definition}[正定、负定]
设二次型$f=\boldsymbol{x}^{T}A\boldsymbol{x}$，若对任何$\boldsymbol{x}\ne\boldsymbol{0}$都有$f(\boldsymbol{x})>0$（显然$f(\boldsymbol{0})=0$），则称$f$为\textbf{正定二次型}，并称对称矩阵$A$是\textbf{正定}的；若对任何$\boldsymbol{x}\ne\boldsymbol{0}$都有$f(\boldsymbol{x})<0$，则称$f$为\textbf{负定二次型}，并称$A$是\textbf{负定}的。
\end{definition}

\begin{theorem}[正定的标准形判据]
$n$元二次型$f=\boldsymbol{x}^{T}A\boldsymbol{x}$为正定的充分必要条件是：它的标准形的$n$个系数全为正，即它的规范形的$n$个系数全为$1$，亦即它的\textbf{正惯性指数等于$n$}。
\end{theorem}

\textbf{证明}\quad 设可逆变换$\boldsymbol{x}=C\boldsymbol{y}$使$f=\sum k_iy_i^{2}$。充分性：若$k_i>0$，则对$\boldsymbol{x}\ne\boldsymbol{0}$有$\boldsymbol{y}=C^{-1}\boldsymbol{x}\ne\boldsymbol{0}$，故$f=\sum k_iy_i^{2}>0$。必要性：反证，若某个$k_i\le0$，取$\boldsymbol{y}=\boldsymbol{e}_i$（单位坐标向量），则$\boldsymbol{x}=C\boldsymbol{e}_i\ne\boldsymbol{0}$而$f(\boldsymbol{e}_i)=k_i\le0$，与正定矛盾。\hfill$\square$

\begin{corollary}[特征值判据]
对称矩阵$A$为正定的充分必要条件是$A$的特征值全为正。
\end{corollary}

\begin{corollary}[$A=Q^{T}Q$判据]
实二次型$f=\boldsymbol{x}^{T}A\boldsymbol{x}$为正定的充分必要条件是：存在可逆矩阵$Q$，使
\[
A=Q^{T}Q .
\]
\end{corollary}

\begin{theorem}[顺序主子式判据，Hurwitz]
对称矩阵$A$为正定的充分必要条件是：$A$的\textbf{各阶顺序主子式全为正}，即
\[
a_{11}>0,\qquad
\begin{vmatrix}a_{11}&a_{12}\\a_{21}&a_{22}\end{vmatrix}>0,\qquad\dots,\qquad
\begin{vmatrix}
a_{11}&a_{12}&\dots&a_{1n}\\
a_{21}&a_{22}&\dots&a_{2n}\\
\vdots&\vdots&&\vdots\\
a_{n1}&a_{n2}&\dots&a_{nn}
\end{vmatrix}>0 .
\]
对称矩阵$A$为负定的充分必要条件是：\textbf{奇数阶顺序主子式为负，偶数阶顺序主子式为正}。
\end{theorem}

这里"$k$阶顺序主子式"指$A$的前$k$行、前$k$列构成的$k$阶子式。

\subsection{正定性的判定步骤}
\begin{enumerate}
    \item 写出二次型的对称矩阵$A$。
    \item 首选\textbf{顺序主子式}：计算$\Delta_1,\Delta_2,\dots,\Delta_n$。全正$\Rightarrow$正定；奇数阶为负、偶数阶为正$\Rightarrow$负定。
    \item 若只需判断"是否正定"而矩阵阶数较大，可先用特征值或配方法求出正惯性指数，看是否为$n$。
    \item 含参数的问题：令$\Delta_k>0$（或负定条件）解出参数范围。
\end{enumerate}

\corepoint{
（1）顺序主子式必须\textbf{从左上角依次取}，不能取任意主子式；只验$|A|>0$不能推出正定（如$\operatorname{diag}(-1,-1)$）。\\
（2）正定$\Rightarrow$$|A|>0$；但$|A|>0$不一定正定，也不一定是负定判据。\\
（3）负定的判据符号是"奇负偶正"，与正定"全正"别混。\\
（4）只有\textbf{实对称矩阵}才谈正定；非对称矩阵即使$x^{T}Ax>0$恒成立，通常也先对称化再讨论。\\
（5）几何意义：二元正定二次型$f(x,y)=c\ (c>0)$是一族以原点为中心的椭圆；三元正定二次型$f(x,y,z)=c\ (c>0)$是一族椭球面——这是正定性最直观的图景。
}

\keyformula{f\ \text{正定}\iff p=n\iff \lambda_i>0\ (\forall i)\iff \Delta_k>0\ (k=1,\dots,n)\iff \exists\,Q\ \text{可逆},\ A=Q^{T}Q}

\keyformula{f\ \text{负定}\iff (-1)^{k}\Delta_k>0\ (k=1,\dots,n),\qquad \text{规范形：}\ z_1^{2}+\dots+z_p^{2}-z_{p+1}^{2}-\dots-z_r^{2}}

\section{拓展：空间中的曲面与二次曲面}

本节是教材之外的几何拓展，把第5节的二次型理论落到三维空间，用于识别三元二次方程所表示的曲面。

\subsection{曲面方程与几类特殊曲面}
给定空间直角坐标系，若曲面$S$上点的坐标都满足方程$F(x,y,z)=0$，且坐标满足该方程的点都在$S$上，则称$F(x,y,z)=0$为曲面$S$的\textbf{方程}，$S$为该方程的\textbf{图形}。

\textbf{球面}\quad 球心$M_0(x_0,y_0,z_0)$、半径$r$的球面方程为
\[
(x-x_0)^{2}+(y-y_0)^{2}+(z-z_0)^{2}=r^{2},
\]
展开后是三元二次方程，其特点是：$x^{2},y^{2},z^{2}$的系数相同且没有$xy,yz,zx$项。具有这三个条件的方程，配方后若右端$k>0$为球面，$k=0$为一点（点球面），$k<0$无图形（虚球面）。

\textbf{柱面}\quad 平行于定直线并沿定曲线$C$移动的直线$L$形成的轨迹叫做\textbf{柱面}，$C$叫\textbf{准线}，$L$叫\textbf{母线}。只含两个变量的方程在空间中表示柱面，其母线平行于方程中不出现的那个变量所对应的坐标轴。例如$\dfrac{x^{2}}{a^{2}}+\dfrac{y^{2}}{b^{2}}=1$、$\dfrac{x^{2}}{a^{2}}-\dfrac{y^{2}}{b^{2}}=1$、$x^{2}=2py$分别表示母线平行于$z$轴的椭圆柱面、双曲柱面、抛物柱面。

\textbf{旋转曲面}\quad 由一条平面曲线$C$绕该平面上的一条定直线$L$旋转一周所成的曲面叫做\textbf{旋转曲面}，$C$叫母线，$L$叫旋转轴。位于$yOz$面上的曲线$f(y,z)=0$绕$z$轴旋转所得曲面方程为
\[
f\bigl(\pm\sqrt{x^{2}+y^{2}},\ z\bigr)=0 ;
\]
绕$y$轴旋转则为$f\bigl(y,\pm\sqrt{x^{2}+z^{2}}\bigr)=0$。口诀：\textbf{与旋转轴同名的变量不动，另一变量换成另两个变量平方和的平方根}。

\textbf{空间曲线}\quad 空间曲线可看作两个曲面的交线，其一般方程为
\[
\begin{cases}F(x,y,z)=0,\\ G(x,y,z)=0,\end{cases}
\]
也可写成参数方程$x=x(t),\ y=y(t),\ z=z(t)$。以$C$为准线、母线垂直于平面$\pi$的柱面与$\pi$的交线叫做$C$在$\pi$上的\textbf{投影}；从$C$的方程中消去$z$得柱面方程$F(x,y)=0$，则$\begin{cases}F(x,y)=0\\ z=0\end{cases}$就是$C$在$xOy$面上的投影曲线。

\subsection{二次曲面的类型}
一般的三元二次方程可写成
\[
\boldsymbol{x}^{T}A\boldsymbol{x}+\boldsymbol{v}^{T}\boldsymbol{x}+a_{44}=0,
\qquad
A=(a_{ij})_{3\times3}\ (a_{ij}=a_{ji}),\ \boldsymbol{x}=(x,y,z)^{T},\ \boldsymbol{v}=(a_{14},a_{24},a_{34})^{T}.
\]
由第5节的正交变换定理，存在正交变换$\boldsymbol{x}=P\boldsymbol{y}$把二次型$\boldsymbol{x}^{T}A\boldsymbol{x}$化为$\lambda_1x'^{2}+\lambda_2y'^{2}+\lambda_3z'^{2}$，即存在适当的直角坐标系使原方程化为
\[
\lambda_1x'^{2}+\lambda_2y'^{2}+\lambda_3z'^{2}+a'_{14}x'+a'_{24}y'+a'_{34}z'+a'_{44}=0 .
\]
对系数的各种情况分类讨论，即得下列九类二次曲面：

\begin{table}[htbp]
\centering
\caption{九类二次曲面的标准方程}
\begin{adjustbox}{max width=\linewidth}
\begin{tabular}{@{}lll@{}}
\toprule
名称 & 标准方程 & 方程特点\\
\midrule
椭圆柱面 & $\dfrac{x^{2}}{a^{2}}+\dfrac{y^{2}}{b^{2}}=1$ & 缺一个变量，母线平行于该轴\\
双曲柱面 & $\dfrac{x^{2}}{a^{2}}-\dfrac{y^{2}}{b^{2}}=1$ & 缺一个变量，两平方项异号\\
抛物柱面 & $x^{2}=2py\ (p>0)$ & 缺一个变量，含一次项\\
椭球面 & $\dfrac{x^{2}}{a^{2}}+\dfrac{y^{2}}{b^{2}}+\dfrac{z^{2}}{c^{2}}=1$ & 三个平方项系数同号\\
单叶双曲面 & $\dfrac{x^{2}}{a^{2}}+\dfrac{y^{2}}{b^{2}}-\dfrac{z^{2}}{c^{2}}=1$ & 两正一负，常数项为$1$\\
双叶双曲面 & $\dfrac{x^{2}}{a^{2}}-\dfrac{y^{2}}{b^{2}}-\dfrac{z^{2}}{c^{2}}=1$ & 一正两负，常数项为$1$\\
椭圆抛物面 & $\dfrac{x^{2}}{2p}+\dfrac{y^{2}}{2q}=z\ (p,q\ \text{同号})$ & 两平方项同号，一次项\\
双曲抛物面 & $\dfrac{x^{2}}{2p}-\dfrac{y^{2}}{2q}=z\ (p,q\ \text{同号})$ & 两平方项异号，一次项\\
二次锥面 & $\dfrac{x^{2}}{a^{2}}+\dfrac{y^{2}}{b^{2}}-\dfrac{z^{2}}{c^{2}}=0$ & 二次齐次，常数项为$0$\\
\bottomrule
\end{tabular}
\end{adjustbox}
\end{table}

\textbf{椭球面}\quad $\dfrac{x^{2}}{a^{2}}+\dfrac{y^{2}}{b^{2}}+\dfrac{z^{2}}{c^{2}}=1$（$a,b,c>0$），$a,b,c$称为三个半轴。$a=b=c$时为球面；有两个相等时为旋转椭球面。图形关于三个坐标面、三个坐标轴及原点都对称，且被限制在长方体$|x|\le a,\ |y|\le b,\ |z|\le c$内。用$z=h$（$|h|<c$）截割，截线是椭圆
\[
\Bigl(\frac{x}{a\sqrt{1-h^{2}/c^{2}}}\Bigr)^{2}+\Bigl(\frac{y}{b\sqrt{1-h^{2}/c^{2}}}\Bigr)^{2}=1,
\]
随$|h|\to c$椭圆缩成点$(0,0,\pm c)$。

\textbf{单叶双曲面}\quad $\dfrac{x^{2}}{a^{2}}+\dfrac{y^{2}}{b^{2}}-\dfrac{z^{2}}{c^{2}}=1$。用$z=h$截得椭圆，半轴随$|h|$增大而增大；用$y=k$截得双曲线，$|k|<b$时实轴平行于$x$轴，$|k|>b$时实轴平行于$z$轴，$|k|=b$时截线退化为两条相交直线。

\textbf{双叶双曲面}\quad $\dfrac{x^{2}}{a^{2}}-\dfrac{y^{2}}{b^{2}}-\dfrac{z^{2}}{c^{2}}=1$。用$z=h$或$y=h$截得双曲线；用$x=h$（$|h|\ge a$）截得椭圆。

\textbf{椭圆抛物面}\quad $\dfrac{x^{2}}{2p}+\dfrac{y^{2}}{2q}=z$（$p,q$同号）。用$z=h$（$h$与$p,q$同号）截得椭圆；用$x=h$或$y=h$截得抛物线。

\textbf{双曲抛物面（马鞍面）}\quad $\dfrac{x^{2}}{2p}-\dfrac{y^{2}}{2q}=z$（$p,q$同号）。用$z=h\ne0$截得双曲线（$h$与$p,q$同号时实轴平行于$x$轴，异号时平行于$y$轴）；用$z=0$截得两条相交直线；用$x=h$或$y=h$截得抛物线。

\textbf{二次锥面}\quad $\dfrac{x^{2}}{a^{2}}+\dfrac{y^{2}}{b^{2}}-\dfrac{z^{2}}{c^{2}}=0$。是二次齐次方程，若$M_0(x_0,y_0,z_0)$在锥面上，则整条过原点与$M_0$的直线都在锥面上（锥面的"直纹"性质）；用$z=h\ne0$截得椭圆，与$z=0$只交于原点。$a=b$时为圆锥面。

\begin{figure}[htbp]
\centering
\begin{tikzpicture}
\begin{axis}[width=0.46\textwidth,height=6cm,view={25}{18},axis lines=center,ticks=none,
  xlabel=$x$,ylabel=$y$,zlabel=$z$,label style={font=\small}]
\addplot3[mesh,draw=black!75,line width=0.3pt,domain=0:360,y domain=-1.6:1.6,samples=30,samples y=8]
({1.0*cos(x)},{0.75*sin(x)},{y});
\end{axis}
\end{tikzpicture}\hfill
\begin{tikzpicture}
\begin{axis}[width=0.46\textwidth,height=6cm,view={25}{18},axis lines=center,ticks=none,
  xlabel=$x$,ylabel=$y$,zlabel=$z$,label style={font=\small}]
\addplot3[mesh,draw=black!75,line width=0.3pt,domain=0:360,y domain=-1.5:1.5,samples=32,samples y=12]
({(0.55+0.20*y^2)*cos(x)},{(0.55+0.20*y^2)*sin(x)},{y});
\end{axis}
\end{tikzpicture}
\caption{左：柱面（母线平行于$z$轴，准线为$xOy$面上的曲线）\quad 右：旋转曲面（$yOz$面上曲线绕$z$轴旋转，$f(\pm\sqrt{x^2+y^2},z)=0$）}
\end{figure}

\begin{figure}[htbp]
\centering
\begin{tikzpicture}
\begin{axis}[width=0.55\textwidth,height=6.2cm,view={28}{18},axis lines=center,ticks=none,
  xlabel=$x$,ylabel=$y$,zlabel=$z$,label style={font=\small}]
\addplot3[mesh,draw=black!75,line width=0.3pt,domain=0:360,y domain=0:180,samples=30,samples y=12]
({1.6*sin(y)*cos(x)},{1.1*sin(y)*sin(x)},{1.2*cos(y)});
\end{axis}
\end{tikzpicture}
\caption{椭球面$\dfrac{x^{2}}{a^{2}}+\dfrac{y^{2}}{b^{2}}+\dfrac{z^{2}}{c^{2}}=1$}
\end{figure}

\begin{figure}[htbp]
\centering
\begin{tikzpicture}
\begin{axis}[width=0.46\textwidth,height=6cm,view={28}{18},axis lines=center,ticks=none,
  xlabel=$x$,ylabel=$y$,zlabel=$z$,label style={font=\small}]
\addplot3[mesh,draw=black!75,line width=0.3pt,domain=0:360,y domain=-1.1:1.1,samples=30,samples y=12]
({1.3*cosh(y)*cos(x)},{1.0*cosh(y)*sin(x)},{0.9*sinh(y)});
\end{axis}
\end{tikzpicture}\hfill
\begin{tikzpicture}
\begin{axis}[width=0.46\textwidth,height=6cm,view={28}{18},axis lines=center,ticks=none,
  xlabel=$x$,ylabel=$y$,zlabel=$z$,label style={font=\small}]
\addplot3[mesh,draw=black!75,line width=0.3pt,domain=0:360,y domain=0:1.35,samples=28,samples y=10]
({1.3*cosh(y)},{0.9*sinh(y)*cos(x)},{1.0*sinh(y)*sin(x)});
\addplot3[mesh,draw=black!75,line width=0.3pt,domain=0:360,y domain=0:1.35,samples=28,samples y=10]
({-1.3*cosh(y)},{0.9*sinh(y)*cos(x)},{1.0*sinh(y)*sin(x)});
\end{axis}
\end{tikzpicture}
\caption{左：单叶双曲面$\dfrac{x^{2}}{a^{2}}+\dfrac{y^{2}}{b^{2}}-\dfrac{z^{2}}{c^{2}}=1$\quad 右：双叶双曲面$\dfrac{x^{2}}{a^{2}}-\dfrac{y^{2}}{b^{2}}-\dfrac{z^{2}}{c^{2}}=1$}
\end{figure}

\begin{figure}[htbp]
\centering
\begin{tikzpicture}
\begin{axis}[width=0.46\textwidth,height=6cm,view={28}{18},axis lines=center,ticks=none,
  xlabel=$x$,ylabel=$y$,zlabel=$z$,label style={font=\small}]
\addplot3[mesh,draw=black!75,line width=0.3pt,samples=24,domain=-1.2:1.2,y domain=-1.2:1.2,samples y=16]
{x^2+y^2};
\end{axis}
\end{tikzpicture}\hfill
\begin{tikzpicture}
\begin{axis}[width=0.46\textwidth,height=6cm,view={28}{18},axis lines=center,ticks=none,
  xlabel=$x$,ylabel=$y$,zlabel=$z$,label style={font=\small}]
\addplot3[mesh,draw=black!75,line width=0.3pt,samples=24,domain=-1.2:1.2,y domain=-1.2:1.2,samples y=16]
{x^2-y^2};
\end{axis}
\end{tikzpicture}
\caption{左：椭圆抛物面$\dfrac{x^{2}}{2p}+\dfrac{y^{2}}{2q}=z$\quad 右：双曲抛物面（马鞍面）$\dfrac{x^{2}}{2p}-\dfrac{y^{2}}{2q}=z$}
\end{figure}

\begin{figure}[htbp]
\centering
\begin{tikzpicture}
\begin{axis}[width=0.42\textwidth,height=6.4cm,view={28}{18},axis lines=center,ticks=none,
  xlabel=$x$,ylabel=$y$,zlabel=$z$,label style={font=\small}]
\addplot3[mesh,draw=black!75,line width=0.3pt,domain=0:360,y domain=-1.3:1.3,samples=30,samples y=10]
({y*cos(x)},{0.85*y*sin(x)},{y});
\end{axis}
\end{tikzpicture}
\caption{二次锥面$\dfrac{x^{2}}{a^{2}}+\dfrac{y^{2}}{b^{2}}-\dfrac{z^{2}}{c^{2}}=0$：$z=h$截得椭圆，$z=0$时仅交于原点$O$}
\end{figure}

\corepoint{
（1）判别二次曲面类型的顺序：先看方程是否为二次齐次（无一次项、常数项）$\to$锥面；再看是否有一次项$\to$抛物面（椭圆/双曲看平方项同号与否）；再看常数项与平方项符号正负个数$\to$椭球面（三同号）、单叶（两正一负）、双叶（一正两负）。\\
（2）"单叶""双叶"的判据是常数为$1$时负号平方项的个数：一个为单叶，两个为双叶；常数项为$0$则是锥面。\\
（3）截痕法（用坐标平面或平行平面去截）是识别曲面的通用手段：一组截线是椭圆、另一组是抛物线$\to$椭圆抛物面；一组椭圆、一组双曲线$\to$单叶/双叶双曲面（看是否连通）。\\
（4）含$x,y,z$交叉项的方程，先用正交变换把二次型化为标准形（求$A$的特征值），再按上表判别，这就是二次型理论在几何中的直接应用。
}


% ---------------- 第 6 章 线性空间与线性变换（选学） ----------------
% 教材 §1 线性空间的定义与性质 / §2 维数、基与坐标 / §3 基变换与坐标变换
%      §4 线性变换 / §5 线性变换的矩阵表示式
\chapter{线性空间与线性变换}
% 第6章 线性空间与线性变换（同济《线性代数》第七版 第6章，选学但保留）
% §1 线性空间的定义与性质；§2 维数、基与坐标；§3 基变换与坐标变换；
% §4 线性变换；§5 线性变换的矩阵表示式。正文不写例题与习题。

\section{线性空间的定义与性质}

第4章的向量是"有序数组"，本章把向量与向量空间的概念公理化：元素是什么并不重要，重要的是元素之间定义了满足八条规律的加法与数乘。

\subsection{线性空间的定义}
\begin{definition}[线性空间]
设$V$是一个非空集合，$F$是一个数域（本书取实数域$\mathbb{R}$）。如果在$V$中定义了一个\textbf{加法}，即对任意$\boldsymbol{\alpha},\boldsymbol{\beta}\in V$，总有唯一的$\boldsymbol{\gamma}\in V$与之对应，记作$\boldsymbol{\gamma}=\boldsymbol{\alpha}+\boldsymbol{\beta}$；又在$V$中定义了一个\textbf{数与元素的乘法}（简称数乘），即对任一$\lambda\in F$与任一$\boldsymbol{\alpha}\in V$，总有唯一的$\boldsymbol{\delta}\in V$与之对应，记作$\boldsymbol{\delta}=\lambda\boldsymbol{\alpha}$；并且这两种运算满足以下八条运算规律（设$\boldsymbol{\alpha},\boldsymbol{\beta},\boldsymbol{\gamma}\in V$，$\lambda,\mu\in F$）：
\begin{enumerate}
    \item $\boldsymbol{\alpha}+\boldsymbol{\beta}=\boldsymbol{\beta}+\boldsymbol{\alpha}$；
    \item $(\boldsymbol{\alpha}+\boldsymbol{\beta})+\boldsymbol{\gamma}=\boldsymbol{\alpha}+(\boldsymbol{\beta}+\boldsymbol{\gamma})$；
    \item 在$V$中存在零元素$\boldsymbol{0}$，对任何$\boldsymbol{\alpha}\in V$都有$\boldsymbol{\alpha}+\boldsymbol{0}=\boldsymbol{\alpha}$；
    \item 对任何$\boldsymbol{\alpha}\in V$，存在$\boldsymbol{\alpha}$的负元素$\boldsymbol{\beta}\in V$，使$\boldsymbol{\alpha}+\boldsymbol{\beta}=\boldsymbol{0}$；
    \item $1\boldsymbol{\alpha}=\boldsymbol{\alpha}$；
    \item $\lambda(\mu\boldsymbol{\alpha})=(\lambda\mu)\boldsymbol{\alpha}$；
    \item $(\lambda+\mu)\boldsymbol{\alpha}=\lambda\boldsymbol{\alpha}+\mu\boldsymbol{\alpha}$；
    \item $\lambda(\boldsymbol{\alpha}+\boldsymbol{\beta})=\lambda\boldsymbol{\alpha}+\lambda\boldsymbol{\beta}$；
\end{enumerate}
那么就称$V$为（数域$F$上的）\textbf{线性空间}（或称向量空间），$V$中的元素统称为\textbf{向量}。
\end{definition}

八条规律的含义：$(i)(ii)$是加法的交换律、结合律；$(iii)(iv)$保证加法有逆运算；$(v)(vi)$与$(vii)(viii)$是数乘的结合律与分配律，$(v)$保证非零数乘有逆运算。满足这八条规律的加法与数乘统称为\textbf{线性运算}。

\corepoint{
线性空间$=$集合$+$线性运算。同一个集合若定义两种不同的线性运算，就构成不同的线性空间；若所定义的运算不满足八条规律，就不是线性空间。所以"运算"才是线性空间的本质，元素本身是什么并不重要。
}

常见与常见反例：
\begin{enumerate}
    \item 次数不超过$n$的多项式全体$P[x]_n=\{a_nx^{n}+\dots+a_1x+a_0\}$，按通常的多项式加法与数乘，构成线性空间（对运算封闭，八条规律自动满足）。
    \item 次数\textbf{恰等于}$n$的多项式全体，按通常运算\textbf{不}构成线性空间：两个$n$次多项式相加，最高次项可能抵消而次数降低，对加法不封闭。
    \item 形如$A\sin(x+B)$的函数全体，按函数加法与数乘构成线性空间。
    \item 正实数全体$\mathbb{R}^{+}$，若定义$a\oplus b=ab$，$\lambda\odot a=a^{\lambda}$，则构成线性空间（零元素是数$1$，$a$的负元素是$a^{-1}$）——这说明"向量"完全可以是数。
    \item 检验一个集合是否为线性空间时，若运算不是通常的实数加、乘，必须逐条验证八条规律，只验封闭性不够。
\end{enumerate}

\subsection{线性空间的性质}
设$V$是线性空间，则
\begin{enumerate}
    \item $V$中的零元素是唯一的；
    \item $V$中每个元素的负元素是唯一的，$\boldsymbol{\alpha}$的负元素记作$-\boldsymbol{\alpha}$；
    \item $0\boldsymbol{\alpha}=\boldsymbol{0}$，$(-1)\boldsymbol{\alpha}=-\boldsymbol{\alpha}$，$\lambda\boldsymbol{0}=\boldsymbol{0}$；
    \item 若$\lambda\boldsymbol{\alpha}=\boldsymbol{0}$，则$\lambda=0$或$\boldsymbol{\alpha}=\boldsymbol{0}$。
\end{enumerate}
（第4条在一般线性空间中不再"显然"，须由八条规律推出：若$\lambda\ne0$，两端乘$\lambda^{-1}$得$\boldsymbol{\alpha}=\lambda^{-1}\boldsymbol{0}=\boldsymbol{0}$。）

\subsection{子空间}
\begin{definition}[子空间]
设$V$是一个线性空间，$L$是$V$的一个非空子集，如果$L$对于$V$中所定义的加法与数乘两种运算也构成一个线性空间，则称$L$为$V$的\textbf{线性子空间}，简称\textbf{子空间}。
\end{definition}

\begin{theorem}[子空间的判定]
线性空间$V$的非空子集$L$构成$V$的子空间的充分必要条件是：$L$对于$V$中的线性运算封闭，即
\[
\boldsymbol{\alpha},\boldsymbol{\beta}\in L\ \Rightarrow\ \boldsymbol{\alpha}+\boldsymbol{\beta}\in L;
\qquad
\boldsymbol{\alpha}\in L,\ \lambda\in F\ \Rightarrow\ \lambda\boldsymbol{\alpha}\in L .
\]
\end{theorem}

因为$L\subseteq V$，八条规律中的$(i)(ii)(v)(vi)(vii)(viii)$自动成立，只需封闭性即可保证$(iii)(iv)$（由$0\boldsymbol{\alpha}=\boldsymbol{0}$、$(-1)\boldsymbol{\alpha}=-\boldsymbol{\alpha}$，零元素与负元素都在$L$中）。

由$V$中向量组$\boldsymbol{\alpha}_1,\dots,\boldsymbol{\alpha}_m$所生成的子空间为
\[
L=\{\boldsymbol{\gamma}\mid \boldsymbol{\gamma}=\lambda_1\boldsymbol{\alpha}_1+\dots+\lambda_m\boldsymbol{\alpha}_m,\ \lambda_i\in F\}
=\operatorname{Span}\{\boldsymbol{\alpha}_1,\dots,\boldsymbol{\alpha}_m\},
\]
它是包含$\boldsymbol{\alpha}_1,\dots,\boldsymbol{\alpha}_m$的最小子空间。

\keyformula{L\le V\ \text{（子空间）}\iff L\ne\varnothing\ \text{且}\ L\ \text{对加法、数乘封闭}}

\keyformula{\operatorname{Span}\{\boldsymbol{\alpha}_1,\dots,\boldsymbol{\alpha}_m\}=\Bigl\{\textstyle\sum_{i=1}^{m}\lambda_i\boldsymbol{\alpha}_i\Bigm|\lambda_i\in F\Bigr\},\qquad \dim\operatorname{Span}\{\boldsymbol{\alpha}_1,\dots,\boldsymbol{\alpha}_m\}=R(\boldsymbol{\alpha}_1,\dots,\boldsymbol{\alpha}_m)}

\section{维数、基与坐标}

\subsection{基与维数}
第4章中由线性运算建立的概念（线性组合、线性相关、最大无关组等）只涉及线性运算，因此对一般线性空间中的向量仍然适用，可直接引用。

\begin{definition}[基、维数]
在线性空间$V$中，如果存在$n$个向量$\boldsymbol{\alpha}_1,\boldsymbol{\alpha}_2,\dots,\boldsymbol{\alpha}_n$，满足
\begin{enumerate}
    \item[(i)] $\boldsymbol{\alpha}_1,\dots,\boldsymbol{\alpha}_n$线性无关；
    \item[(ii)] $V$中任一向量总可由$\boldsymbol{\alpha}_1,\dots,\boldsymbol{\alpha}_n$线性表示，
\end{enumerate}
那么$\boldsymbol{\alpha}_1,\dots,\boldsymbol{\alpha}_n$就称为线性空间$V$的一个\textbf{基}（基底），$n$称为线性空间$V$的\textbf{维数}，记作$\dim V=n$。维数为$n$的线性空间称为\textbf{$n$维线性空间}，记作$V_n$。
\end{definition}

只含一个零向量的线性空间没有基，规定其维数为$0$。$n$维与零维线性空间统称\textbf{有限维线性空间}；本书不讨论无限维线性空间。

由第4章的理论：若把线性空间看作向量组，则$V$的基就是该向量组的最大无关组，$V$的维数就是该向量组的秩。对$n$维线性空间$V_n$，若$\boldsymbol{\alpha}_1,\dots,\boldsymbol{\alpha}_n$为其一个基，则
\[
V_n=\operatorname{Span}\{\boldsymbol{\alpha}_1,\dots,\boldsymbol{\alpha}_n\},
\]
这就比较清楚地显示出$V_n$的构造。

\subsection{坐标}
\begin{definition}[坐标]
设$\boldsymbol{\alpha}_1,\dots,\boldsymbol{\alpha}_n$是线性空间$V_n$的一个基，对于任一向量$\boldsymbol{\alpha}\in V$，总有且仅有一组有序数$x_1,x_2,\dots,x_n\in F$，使
\[
\boldsymbol{\alpha}=x_1\boldsymbol{\alpha}_1+x_2\boldsymbol{\alpha}_2+\dots+x_n\boldsymbol{\alpha}_n,
\]
这组有序数就称为向量$\boldsymbol{\alpha}$在该基中的\textbf{坐标}，并记作$(x_1,x_2,\dots,x_n)$。
\end{definition}

取定一个基以后，$V_n$中的向量$\boldsymbol{\alpha}$与$\mathbb{R}^{n}$中的坐标向量$\boldsymbol{x}=(x_1,\dots,x_n)^{T}$之间就有一一对应的关系，且这个对应关系\textbf{保持线性运算}：
\[
\boldsymbol{\alpha}\leftrightarrow\boldsymbol{x},\quad \boldsymbol{\beta}\leftrightarrow\boldsymbol{y}
\ \Longrightarrow\ 
k\boldsymbol{\alpha}+l\boldsymbol{\beta}\leftrightarrow k\boldsymbol{x}+l\boldsymbol{y}.
\]
因此$V_n$与$\mathbb{R}^{n}$有相同的结构，称$V_n$与$\mathbb{R}^{n}$\textbf{同构}。一般地，两个线性空间之间若存在保持线性运算的一一对应，就说它们同构；维数相等的有限维线性空间都同构，即\textbf{线性空间的结构完全由它的维数决定}。

但要注意：$\mathbb{R}^{n}$中的内积概念在一般线性空间$V_n$中不一定有意义——同构只保证线性运算，不保证内积。

\subsection{求坐标的步骤}
\begin{enumerate}
    \item 设所求坐标为$\boldsymbol{x}=(x_1,\dots,x_n)^{T}$，按定义把$\boldsymbol{\alpha}=x_1\boldsymbol{\alpha}_1+\dots+x_n\boldsymbol{\alpha}_n$写成矩阵形式
    \[
    (\boldsymbol{\alpha}_1,\boldsymbol{\alpha}_2,\dots,\boldsymbol{\alpha}_n)\boldsymbol{x}=\boldsymbol{\alpha}.
    \]
    \item 若在$\mathbb{R}^{n}$中，令$A=(\boldsymbol{\alpha}_1,\dots,\boldsymbol{\alpha}_n)$，解方程组$A\boldsymbol{x}=\boldsymbol{\alpha}$（用初等行变换化$(A,\boldsymbol{\alpha})$为行最简形），所得解即坐标。
    \item 若题目只问"是否为基"，等价于判断$\boldsymbol{\alpha}_1,\dots,\boldsymbol{\alpha}_n$线性无关（$n$个向量在$n$维空间中），即$R(A)=n$或$|A|\ne0$。
\end{enumerate}

\corepoint{
（1）坐标是\textbf{相对于某个基}而言的，换基则坐标改变，写坐标必须说明"在基$\boldsymbol{\alpha}_1,\dots,\boldsymbol{\alpha}_n$下"。\\
（2）求坐标是解线性方程组，不是"把基按行排好直接读"；只有取自然基$\boldsymbol{e}_1,\dots,\boldsymbol{e}_n$时，坐标才恰好等于分量。\\
（3）多项式空间中求坐标，要先把向量用基"归并同类项"比较系数，本质仍是解方程组。
}

\section{基变换与坐标变换}

同一向量在不同基中的坐标不同，两者之间的关系由过渡矩阵刻画。

\subsection{基变换公式与过渡矩阵}
设$\boldsymbol{\alpha}_1,\dots,\boldsymbol{\alpha}_n$及$\boldsymbol{\beta}_1,\dots,\boldsymbol{\beta}_n$是线性空间$V_n$的两个基，把新基向量用旧基线性表示：
\[
\begin{cases}
\boldsymbol{\beta}_1=p_{11}\boldsymbol{\alpha}_1+p_{21}\boldsymbol{\alpha}_2+\dots+p_{n1}\boldsymbol{\alpha}_n,\\
\boldsymbol{\beta}_2=p_{12}\boldsymbol{\alpha}_1+p_{22}\boldsymbol{\alpha}_2+\dots+p_{n2}\boldsymbol{\alpha}_n,\\
\qquad\qquad\cdots\cdots\cdots\cdots\\
\boldsymbol{\beta}_n=p_{1n}\boldsymbol{\alpha}_1+p_{2n}\boldsymbol{\alpha}_2+\dots+p_{nn}\boldsymbol{\alpha}_n,
\end{cases}
\]
记$P=(p_{ij})_{n\times n}$，用矩阵形式写作
\[
(\boldsymbol{\beta}_1,\boldsymbol{\beta}_2,\dots,\boldsymbol{\beta}_n)
=(\boldsymbol{\alpha}_1,\boldsymbol{\alpha}_2,\dots,\boldsymbol{\alpha}_n)
\begin{pmatrix}p_{11}&p_{12}&\dots&p_{1n}\\p_{21}&p_{22}&\dots&p_{2n}\\\vdots&\vdots&\ddots&\vdots\\p_{n1}&p_{n2}&\dots&p_{nn}\end{pmatrix}
=(\boldsymbol{\alpha}_1,\dots,\boldsymbol{\alpha}_n)P .
\]
此式称为\textbf{基变换公式}，矩阵$P$称为由基$\boldsymbol{\alpha}_1,\dots,\boldsymbol{\alpha}_n$到基$\boldsymbol{\beta}_1,\dots,\boldsymbol{\beta}_n$的\textbf{过渡矩阵}。由于$\boldsymbol{\beta}_1,\dots,\boldsymbol{\beta}_n$线性无关，过渡矩阵$P$必可逆。

\subsection{坐标变换公式}
\begin{theorem}[坐标变换公式]
设$V_n$中的向量$\boldsymbol{\alpha}$在基$\boldsymbol{\alpha}_1,\dots,\boldsymbol{\alpha}_n$中的坐标为$(x_1,\dots,x_n)^{T}$，在基$\boldsymbol{\beta}_1,\dots,\boldsymbol{\beta}_n$中的坐标为$(x_1',\dots,x_n')^{T}$。若两个基满足基变换公式$(\boldsymbol{\beta}_1,\dots,\boldsymbol{\beta}_n)=(\boldsymbol{\alpha}_1,\dots,\boldsymbol{\alpha}_n)P$，则有
\[
\begin{pmatrix}x_1\\x_2\\\vdots\\x_n\end{pmatrix}
=P\begin{pmatrix}x_1'\\x_2'\\\vdots\\x_n'\end{pmatrix},
\qquad\text{或}\qquad
\begin{pmatrix}x_1'\\x_2'\\\vdots\\x_n'\end{pmatrix}
=P^{-1}\begin{pmatrix}x_1\\x_2\\\vdots\\x_n\end{pmatrix}.
\]
\end{theorem}

\textbf{证明}\quad 由$\boldsymbol{\alpha}=(\boldsymbol{\alpha}_1,\dots,\boldsymbol{\alpha}_n)\boldsymbol{x}=(\boldsymbol{\beta}_1,\dots,\boldsymbol{\beta}_n)\boldsymbol{x}'=(\boldsymbol{\alpha}_1,\dots,\boldsymbol{\alpha}_n)P\boldsymbol{x}'$，而$\boldsymbol{\alpha}_1,\dots,\boldsymbol{\alpha}_n$线性无关，故$\boldsymbol{x}=P\boldsymbol{x}'$。\hfill$\square$

定理的逆命题也成立：若任一向量的两种坐标都满足$\boldsymbol{x}=P\boldsymbol{x}'$，则两个基满足$(\boldsymbol{\beta})=(\boldsymbol{\alpha})P$。

\subsection{求过渡矩阵与坐标的方法}
\begin{enumerate}
    \item \textbf{代入法（定义法）}：把每个新基向量用旧基线性表示，解得系数，按"第$j$列是新基第$j$个向量的坐标"排成$P$。
    \item \textbf{初等变换法}：在$\mathbb{R}^{n}$中，令$A=(\boldsymbol{\alpha}_1,\dots,\boldsymbol{\alpha}_n)$，$B=(\boldsymbol{\beta}_1,\dots,\boldsymbol{\beta}_n)$，则$B=AP$，$P=A^{-1}B$。对$(A\mid B)$作初等行变换，当$A$化成$E$时，$B$就化成了$P$，即
    \[
    (A\mid B)\ \xrightarrow{\text{行变换}}\ (E\mid P).
    \]
    \item \textbf{坐标变换}：先用上法求$P$，再用$\boldsymbol{x}'=P^{-1}\boldsymbol{x}$（或解$P\boldsymbol{x}'=\boldsymbol{x}$）求新坐标。
\end{enumerate}

\keyformula{(\boldsymbol{\beta}_1,\dots,\boldsymbol{\beta}_n)=(\boldsymbol{\alpha}_1,\dots,\boldsymbol{\alpha}_n)P\quad(\text{基变换}),\qquad \boldsymbol{x}=P\boldsymbol{x}'\quad(\text{旧坐标}=\text{过渡矩阵}\times\text{新坐标})}

\keyformula{P=A^{-1}B,\qquad (A\mid B)\longrightarrow(E\mid P)}

\corepoint{
（1）方向极易记反：$P$是"由旧基\textbf{到}新基"的过渡矩阵，满足$(\text{新基})=(\text{旧基})P$；坐标变换则是$\boldsymbol{x}=P\boldsymbol{x}'$，即"旧坐标$=$过渡矩阵$\times$新坐标"，两者方向一致，不必再记$P^{-1}$在左边还是右边，直接用$\boldsymbol{x}=P\boldsymbol{x}'$推。\\
（2）$P$的列必须是新基向量在旧基下的坐标；若把$P$的行写成坐标，转置一下就错了。\\
（3）用$(A\mid B)\to(E\mid P)$时只能做\textbf{行变换}；若$A$不可逆说明$B$不是基，应回头检查。
}

\section{线性变换}

\subsection{映射与线性变换}
\begin{definition}[映射]
设有两个非空集合$A,B$，如果对于$A$中任一元素$\boldsymbol{\alpha}$，按照一定的规则，总有$B$中一个确定的元素$\boldsymbol{\beta}$和它对应，那么这个对应规则称为从$A$到$B$的\textbf{映射}，记作$T$，并记$\boldsymbol{\beta}=T(\boldsymbol{\alpha})$或$\boldsymbol{\beta}=T\boldsymbol{\alpha}$。$\boldsymbol{\beta}$称为$\boldsymbol{\alpha}$在映射$T$下的\textbf{像}，$\boldsymbol{\alpha}$称为$\boldsymbol{\beta}$的\textbf{原像}；$A$称为$T$的\textbf{定义域}，像的全体$T(A)=\{\boldsymbol{\beta}\mid\boldsymbol{\beta}=T(\boldsymbol{\alpha}),\ \boldsymbol{\alpha}\in A\}$称为\textbf{像集}。
\end{definition}

映射是函数概念的推广。

\begin{definition}[线性映射、线性变换]
设$V_n,U_m$分别是$n$维和$m$维线性空间，$T$是从$V_n$到$U_m$的映射。如果$T$满足
\begin{enumerate}
    \item[(i)] 任给$\boldsymbol{\alpha}_1,\boldsymbol{\alpha}_2\in V_n$，有$T(\boldsymbol{\alpha}_1+\boldsymbol{\alpha}_2)=T(\boldsymbol{\alpha}_1)+T(\boldsymbol{\alpha}_2)$；
    \item[(ii)] 任给$\boldsymbol{\alpha}\in V_n$，$\lambda\in F$，有$T(\lambda\boldsymbol{\alpha})=\lambda T(\boldsymbol{\alpha})$，
\end{enumerate}
那么$T$就称为从$V_n$到$U_m$的\textbf{线性映射}。$V_n$到自身的线性映射称为$V_n$中的\textbf{线性变换}。
\end{definition}

简言之，线性映射就是\textbf{保持线性组合}的映射：
\[
T(k_1\boldsymbol{\alpha}_1+\dots+k_m\boldsymbol{\alpha}_m)=k_1T(\boldsymbol{\alpha}_1)+\dots+k_mT(\boldsymbol{\alpha}_m).
\]

典型例子：多项式空间的求导运算$D$（$Dp=p'$）是线性变换；$\mathbb{R}^{n}$中由$n$阶矩阵$A$确定的变换$T(\boldsymbol{x})=A\boldsymbol{x}$是线性变换；平面上的旋转变换、把向量投影到某平面上的变换都是线性变换。而$T(p)=1$（把任何多项式映成常数$1$）不是线性变换，因为$T(p+q)=1\ne1+1=T(p)+T(q)$。

零变换$\Gamma_0$（把一切向量映成$\boldsymbol{0}$）与恒等变换$I$（$I(\boldsymbol{\alpha})=\boldsymbol{\alpha}$）都是线性变换。

\subsection{线性变换的基本性质}
设$T$是$V_n$中的线性变换，则
\begin{enumerate}
    \item $T(\boldsymbol{0})=\boldsymbol{0}$，$T(-\boldsymbol{\alpha})=-T(\boldsymbol{\alpha})$；
    \item 若$\boldsymbol{\beta}=k_1\boldsymbol{\alpha}_1+\dots+k_m\boldsymbol{\alpha}_m$，则$T(\boldsymbol{\beta})=k_1T(\boldsymbol{\alpha}_1)+\dots+k_mT(\boldsymbol{\alpha}_m)$（线性变换由它在一组基上的像唯一确定）；
    \item 若$\boldsymbol{\alpha}_1,\dots,\boldsymbol{\alpha}_m$线性相关，则$T(\boldsymbol{\alpha}_1),\dots,T(\boldsymbol{\alpha}_m)$也线性相关；
    \item 像集$T(V_n)$是$V_n$的子空间，称为$T$的\textbf{像空间}；
    \item 使$T(\boldsymbol{\alpha})=\boldsymbol{0}$的$\boldsymbol{\alpha}$的全体
    \[
    N_T=\ker T=\{\boldsymbol{\alpha}\in V_n\mid T(\boldsymbol{\alpha})=\boldsymbol{0}\}
    \]
    也是$V_n$的子空间，称为$T$的\textbf{核}。
\end{enumerate}

注意性质3的\textbf{逆命题不成立}：线性无关的向量组经线性变换后可能变得线性相关（例如投影变换把平面内两条不共线向量都投到同一条直线上）。这是因为线性变换可能"压平"空间。

\corepoint{
（1）判断一个映射是否为线性变换，只需验加法与数乘两条；常用来否定的是"常数项不为零"的映射（如$T(\boldsymbol{x})=A\boldsymbol{x}+\boldsymbol{b}$，$\boldsymbol{b}\ne\boldsymbol{0}$不是线性变换）。\\
（2）$T(\boldsymbol{0})=\boldsymbol{0}$是必要条件：若$T(\boldsymbol{0})\ne\boldsymbol{0}$，则$T$一定不是线性变换。\\
（3）像空间与核都是子空间，不要把"核"与"零向量"混淆：核一般含无穷多向量，只是其中包含零向量。
}

\section{线性变换的矩阵表示式}

\subsection{线性变换在给定基下的矩阵}
设$T$是$V_n$中的线性变换，$\boldsymbol{\varepsilon}_1,\boldsymbol{\varepsilon}_2,\dots,\boldsymbol{\varepsilon}_n$是$V_n$的一个基。由于$T(\boldsymbol{\varepsilon}_j)$仍在$V_n$中，故可由该基唯一地线性表示：
\[
\begin{cases}
T(\boldsymbol{\varepsilon}_1)=a_{11}\boldsymbol{\varepsilon}_1+a_{21}\boldsymbol{\varepsilon}_2+\dots+a_{n1}\boldsymbol{\varepsilon}_n,\\
T(\boldsymbol{\varepsilon}_2)=a_{12}\boldsymbol{\varepsilon}_1+a_{22}\boldsymbol{\varepsilon}_2+\dots+a_{n2}\boldsymbol{\varepsilon}_n,\\
\qquad\qquad\cdots\cdots\cdots\cdots\\
T(\boldsymbol{\varepsilon}_n)=a_{1n}\boldsymbol{\varepsilon}_1+a_{2n}\boldsymbol{\varepsilon}_2+\dots+a_{nn}\boldsymbol{\varepsilon}_n .
\end{cases}
\]
记
\[
A=\begin{pmatrix}
a_{11}&a_{12}&\dots&a_{1n}\\
a_{21}&a_{22}&\dots&a_{2n}\\
\vdots&\vdots&\ddots&\vdots\\
a_{n1}&a_{n2}&\dots&a_{nn}
\end{pmatrix},
\]
则上式可形式地写成
\[
T(\boldsymbol{\varepsilon}_1,\boldsymbol{\varepsilon}_2,\dots,\boldsymbol{\varepsilon}_n)
=\bigl(T(\boldsymbol{\varepsilon}_1),T(\boldsymbol{\varepsilon}_2),\dots,T(\boldsymbol{\varepsilon}_n)\bigr)
=(\boldsymbol{\varepsilon}_1,\boldsymbol{\varepsilon}_2,\dots,\boldsymbol{\varepsilon}_n)A .
\]

\begin{definition}[线性变换的矩阵]
上述矩阵$A$称为线性变换$T$在基$\boldsymbol{\varepsilon}_1,\boldsymbol{\varepsilon}_2,\dots,\boldsymbol{\varepsilon}_n$下的\textbf{矩阵}。
\end{definition}

矩阵$A$由基的像$T(\boldsymbol{\varepsilon}_1),\dots,T(\boldsymbol{\varepsilon}_n)$唯一确定；反之，给定一个$n$阶矩阵$A$，也就给定了基的像，从而唯一确定一个以$A$为矩阵的线性变换。因此在\textbf{取定一组基之后，线性变换与$n$阶矩阵一一对应}。

若$V_n$中向量$\boldsymbol{\alpha}$在该基下的坐标为$\boldsymbol{x}=(x_1,\dots,x_n)^{T}$，则$T(\boldsymbol{\alpha})$在同一基下的坐标为
\[
\boldsymbol{y}=A\boldsymbol{x},
\quad\text{即}\quad
T\bigl((\boldsymbol{\varepsilon}_1,\dots,\boldsymbol{\varepsilon}_n)\boldsymbol{x}\bigr)
=(\boldsymbol{\varepsilon}_1,\dots,\boldsymbol{\varepsilon}_n)A\boldsymbol{x}.
\]

\subsection{不同基下矩阵的关系}
\begin{theorem}[同一线性变换在不同基下的矩阵相似]
设线性空间$V_n$中取定两个基$\boldsymbol{\alpha}_1,\dots,\boldsymbol{\alpha}_n$和$\boldsymbol{\beta}_1,\dots,\boldsymbol{\beta}_n$，由基$\boldsymbol{\alpha}_1,\dots,\boldsymbol{\alpha}_n$到基$\boldsymbol{\beta}_1,\dots,\boldsymbol{\beta}_n$的过渡矩阵为$P$；$V_n$中的线性变换$T$在这两个基下的矩阵依次为$A$和$B$，则
\[
B=P^{-1}AP .
\]
\end{theorem}

\textbf{证明}\quad 由$(\boldsymbol{\beta}_1,\dots,\boldsymbol{\beta}_n)=(\boldsymbol{\alpha}_1,\dots,\boldsymbol{\alpha}_n)P$得
\[
\begin{aligned}
T(\boldsymbol{\beta}_1,\dots,\boldsymbol{\beta}_n)
&=T\bigl((\boldsymbol{\alpha}_1,\dots,\boldsymbol{\alpha}_n)P\bigr)
=\bigl(T(\boldsymbol{\alpha}_1,\dots,\boldsymbol{\alpha}_n)\bigr)P\\
&=(\boldsymbol{\alpha}_1,\dots,\boldsymbol{\alpha}_n)AP
=(\boldsymbol{\beta}_1,\dots,\boldsymbol{\beta}_n)P^{-1}AP,
\end{aligned}
\]
又$T(\boldsymbol{\beta}_1,\dots,\boldsymbol{\beta}_n)=(\boldsymbol{\beta}_1,\dots,\boldsymbol{\beta}_n)B$，而$\boldsymbol{\beta}_1,\dots,\boldsymbol{\beta}_n$线性无关，故$B=P^{-1}AP$。\hfill$\square$

定理表明：\textbf{同一线性变换在不同基下的矩阵是相似的，相似变换矩阵就是过渡矩阵$P$}。这是"相似"概念的几何意义——相似矩阵正是同一个线性变换在不同基下的矩阵。

\subsection{线性变换的秩}
\begin{definition}[线性变换的秩]
线性变换$T$的像空间$T(V_n)$的维数，称为线性变换$T$的\textbf{秩}。
\end{definition}

若$A$是$T$的矩阵，则$T$的秩就等于$R(A)$；若$R(A)=r$，则$T$的核的维数为
\[
\dim\ker T=n-r .
\]
（因为$\ker T$就是齐次方程组$A\boldsymbol{x}=\boldsymbol{0}$的解空间。）

\keyformula{T(\boldsymbol{\varepsilon}_1,\dots,\boldsymbol{\varepsilon}_n)=(\boldsymbol{\varepsilon}_1,\dots,\boldsymbol{\varepsilon}_n)A,\qquad T(\boldsymbol{\alpha})\ \text{的坐标}=A\boldsymbol{x}}

\keyformula{B=P^{-1}AP\ (\text{换基}),\qquad \operatorname{rank}T=R(A),\qquad \dim\ker T=n-R(A)}

\corepoint{
（1）求$T$在基$\boldsymbol{\varepsilon}_1,\dots,\boldsymbol{\varepsilon}_n$下的矩阵，关键是算$T(\boldsymbol{\varepsilon}_j)$并把它们\textbf{用该基}表示，系数按\textbf{列}排成$A$；若把系数按行排，得到的是$A^{T}$。\\
（2）同一变换换基后矩阵变为$P^{-1}AP$，不是$P^{T}AP$（那是合同变换），也不是$P^{-1}AP^{-1}$。\\
（3）"线性变换与矩阵一一对应"依赖于\textbf{取定一组基}；换基后对应同一个变换的矩阵不同，但它们相似。\\
（4）$\dim$像$+\dim$核$=n$是验算的有力工具。
}


% ---------------- 附录（教材无对应章，保留拓展内容） ----------------
\appendix
% 附录章标题显示为「附录 A」而不是「第 A 章」
\titleformat{\chapter}[display]
  {\sffamily\bfseries\color{yhbAccent}}
  {\normalsize\sffamily 附录\,\thechapter}
  {6pt}
  {\LARGE}
  [{\vspace{5pt}{\color{yhbAccent!55}\titlerule[1pt]}}]

\chapter{几何向量与二次曲面}
% file: 5.几何向量.tex
% 附录A 几何向量与二次曲面。线性代数教材无此章，内容以同济《高等数学》下册（第八版）
% 第八章「向量代数与空间解析几何」为主要依据，另补入解析几何中教材未展开的少数标准结论
% （点到平面/直线的距离公式、平面束），并写在明确的小节里。
% 只写知识点：定义、定理（写清条件）、核心公式、方法步骤、易错点；不写例题与习题。

\section{向量及其线性运算}

\subsection{向量的概念}

客观世界中有一类量，它们既有大小又有方向，例如位移、速度、加速度、力、力矩，这一类量叫做\textbf{向量}（或矢量）。数学上用一条\textbf{有向线段}表示向量：有向线段的长度表示向量的大小，有向线段的方向表示向量的方向。以 $A$ 为起点、$B$ 为终点的向量记作 $\overrightarrow{AB}$，也常用黑体字母 $\boldsymbol{a},\boldsymbol{b},\boldsymbol{r}$ 表示。

\begin{definition}
只考虑大小与方向、而不管起点位置的向量称为\textbf{自由向量}（本附录所讨论的向量均指自由向量）。若两个向量经平移后能完全重合，即它们的\textbf{模相等且方向相同}，则称这两个向量\textbf{相等}，记作 $\boldsymbol{a}=\boldsymbol{b}$。向量 $\boldsymbol{a}$ 的大小叫做它的\textbf{模}，记作 $|\boldsymbol{a}|$ 或 $|\overrightarrow{AB}|$。
\end{definition}

\begin{definition}
模等于 $1$ 的向量叫做\textbf{单位向量}；模等于 $0$ 的向量叫做\textbf{零向量}，记作 $\boldsymbol{0}$，规定零向量的方向任意。与 $\boldsymbol{a}$ 模相同、方向相反的向量叫做 $\boldsymbol{a}$ 的\textbf{负向量}，记作 $-\boldsymbol{a}$。
\end{definition}

\begin{definition}
设 $\boldsymbol{a},\boldsymbol{b}$ 是两个非零向量，任取空间一点 $O$，作 $\overrightarrow{OA}=\boldsymbol{a}$，$\overrightarrow{OB}=\boldsymbol{b}$，规定不超过 $\pi$ 的角 $\angle AOB$ 为 $\boldsymbol{a}$ 与 $\boldsymbol{b}$ 的\textbf{夹角}，记作 $(\boldsymbol{a},\boldsymbol{b})$ 或 $\langle\boldsymbol{a},\boldsymbol{b}\rangle$，于是 $0\le\langle\boldsymbol{a},\boldsymbol{b}\rangle\le\pi$。
\begin{enumerate}
\item 若 $\langle\boldsymbol{a},\boldsymbol{b}\rangle=0$ 或 $\pi$，称 $\boldsymbol{a}$ 与 $\boldsymbol{b}$ \textbf{平行}（共线），记作 $\boldsymbol{a}\parallel\boldsymbol{b}$；
\item 若 $\langle\boldsymbol{a},\boldsymbol{b}\rangle=\frac{\pi}{2}$，称 $\boldsymbol{a}$ 与 $\boldsymbol{b}$ \textbf{垂直}，记作 $\boldsymbol{a}\perp\boldsymbol{b}$。
\end{enumerate}
若 $\boldsymbol{a},\boldsymbol{b}$ 中有一个是零向量，规定它们的夹角可在 $[0,\pi]$ 内任意取值；因此零向量与任何向量都既平行又垂直。
\end{definition}

设有 $n\ (n\ge3)$ 个向量，把它们平移到同一起点，若各终点与这个公共起点在同一个平面上，就称这 $n$ 个向量\textbf{共面}。

\subsection{向量的加减法}

\begin{definition}[向量的加法]
设有两个向量 $\boldsymbol{a},\boldsymbol{b}$，任取一点 $A$，作 $\overrightarrow{AB}=\boldsymbol{a}$，再以 $B$ 为起点作 $\overrightarrow{BC}=\boldsymbol{b}$，则向量 $\overrightarrow{AC}=\boldsymbol{c}$ 叫做 $\boldsymbol{a}$ 与 $\boldsymbol{b}$ 的\textbf{和}，记作 $\boldsymbol{c}=\boldsymbol{a}+\boldsymbol{b}$。
\end{definition}

上述作法称为向量相加的\textbf{三角形法则}；当 $\boldsymbol{a},\boldsymbol{b}$ 不平行时，以 $\overrightarrow{AB}=\boldsymbol{a}$、$\overrightarrow{AD}=\boldsymbol{b}$ 为邻边作平行四边形，其对角线向量 $\overrightarrow{AC}=\boldsymbol{a}+\boldsymbol{b}$，这称为\textbf{平行四边形法则}（力学的合力法则）。

\begin{theorem}
向量的加法满足
\[
\boldsymbol{a}+\boldsymbol{b}=\boldsymbol{b}+\boldsymbol{a},\qquad
(\boldsymbol{a}+\boldsymbol{b})+\boldsymbol{c}=\boldsymbol{a}+(\boldsymbol{b}+\boldsymbol{c}).
\]
\end{theorem}

由交换律与结合律，$n$ 个向量相加可任意分组、任意换序，并有\textbf{多边形法则}：以前一个向量的终点作为次一个向量的起点相继作出这 $n$ 个向量，则以第一个向量的起点为起点、以最后一个向量的终点为终点的向量就是它们的和。

两个向量 $\boldsymbol{a}$ 与 $\boldsymbol{b}$ 的\textbf{差}规定为 $\boldsymbol{b}-\boldsymbol{a}=\boldsymbol{b}+(-\boldsymbol{a})$。把 $\boldsymbol{a},\boldsymbol{b}$ 平移到同一起点 $O$，则从 $\boldsymbol{a}$ 的终点指向 $\boldsymbol{b}$ 的终点的向量就是 $\boldsymbol{b}-\boldsymbol{a}$。

\keyformula{\big||\boldsymbol{a}|-|\boldsymbol{b}|\big|\ \le\ |\boldsymbol{a}\pm\boldsymbol{b}|\ \le\ |\boldsymbol{a}|+|\boldsymbol{b}|}

左端等号当且仅当 $\boldsymbol{a},\boldsymbol{b}$ 反向时成立，右端等号当且仅当 $\boldsymbol{a},\boldsymbol{b}$ 同向时成立。

\subsection{向量与数的乘法}

\begin{definition}
向量 $\boldsymbol{a}$ 与实数 $\lambda$ 的乘积 $\lambda\boldsymbol{a}$ 规定为一个向量：它的模 $|\lambda\boldsymbol{a}|=|\lambda|\,|\boldsymbol{a}|$；它的方向当 $\lambda>0$ 时与 $\boldsymbol{a}$ 相同，当 $\lambda<0$ 时与 $\boldsymbol{a}$ 相反，当 $\lambda=0$ 时 $\lambda\boldsymbol{a}=\boldsymbol{0}$。特别地 $1\boldsymbol{a}=\boldsymbol{a}$，$(-1)\boldsymbol{a}=-\boldsymbol{a}$。
\end{definition}

\begin{theorem}
数乘满足
\[
\lambda(\mu\boldsymbol{a})=(\lambda\mu)\boldsymbol{a},\qquad
(\lambda+\mu)\boldsymbol{a}=\lambda\boldsymbol{a}+\mu\boldsymbol{a},\qquad
\lambda(\boldsymbol{a}+\boldsymbol{b})=\lambda\boldsymbol{a}+\lambda\boldsymbol{b}.
\]
\end{theorem}

向量相加与数乘向量统称为\textbf{向量的线性运算}。取 $\lambda=1/|\boldsymbol{a}|$，得一个与 $\boldsymbol{a}$ 同方向、模为 $1$ 的向量，称为 $\boldsymbol{a}$ 的\textbf{单位化}：

\keyformula{\boldsymbol{e}_{\boldsymbol{a}}=\frac{\boldsymbol{a}}{|\boldsymbol{a}|}\qquad(\boldsymbol{a}\neq\boldsymbol{0})}

\begin{theorem}[两向量平行的充要条件]
设 $\boldsymbol{a}\neq\boldsymbol{0}$，则 $\boldsymbol{b}\parallel\boldsymbol{a}$ 的充分必要条件是：存在\textbf{唯一}的实数 $\lambda$，使 $\boldsymbol{b}=\lambda\boldsymbol{a}$。
\end{theorem}

必要性：若 $\boldsymbol{b}\parallel\boldsymbol{a}$，取 $\lambda=\pm\dfrac{|\boldsymbol{b}|}{|\boldsymbol{a}|}$（同向取正号、反向取负号）。唯一性：若 $\lambda\boldsymbol{a}=\mu\boldsymbol{a}$，则 $(\lambda-\mu)\boldsymbol{a}=\boldsymbol{0}$，而 $\boldsymbol{a}\neq\boldsymbol{0}$，故 $\lambda=\mu$。此定理是建立数轴的理论依据。

\subsection{空间直角坐标系与向量的坐标}

在空间取定一点 $O$ 和三个两两垂直的单位向量 $\boldsymbol{i},\boldsymbol{j},\boldsymbol{k}$，就确定了三条都以 $O$ 为原点的两两垂直的数轴，依次称为 $x$ 轴（横轴）、$y$ 轴（纵轴）、$z$ 轴（竖轴），统称\textbf{坐标轴}；它们构成\textbf{空间直角坐标系}，记作 $O$-$xyz$。当 $\boldsymbol{i},\boldsymbol{j},\boldsymbol{k}$ 的正向符合\textbf{右手规则}（右手四指从 $x$ 轴正向以 $\pi/2$ 转向 $y$ 轴正向时，大拇指的指向就是 $z$ 轴正向）时，称为右手坐标系。三条坐标轴中任意两条确定的平面统称\textbf{坐标面}，分别记作 $xOy$ 面、$yOz$ 面、$zOx$ 面；三个坐标面把空间分成八个部分，每一部分叫做一个\textbf{卦限}。

任给向量 $\boldsymbol{r}$，作 $\overrightarrow{OM}=\boldsymbol{r}$，以 $OM$ 为对角线、三条坐标轴为棱作长方体，则

\[
\boldsymbol{r}=\overrightarrow{OM}=\overrightarrow{OP}+\overrightarrow{OQ}+\overrightarrow{OR}
=x\boldsymbol{i}+y\boldsymbol{j}+z\boldsymbol{k}.
\]

上式称为向量 $\boldsymbol{r}$ 的\textbf{坐标分解式}，$x\boldsymbol{i},\,y\boldsymbol{j},\,z\boldsymbol{k}$ 称为 $\boldsymbol{r}$ 沿三个坐标轴方向的\textbf{分向量}，有序数 $x,y,z$ 称为 $\boldsymbol{r}$ 的\textbf{坐标}，记作 $\boldsymbol{r}=(x,y,z)$。

点 $M$ 与该点的向径 $\overrightarrow{OM}$ 有相同的坐标，故记号 $(x,y,z)$ 既可表示点 $M$，又可表示向量 $\overrightarrow{OM}$；但点是几何对象、向量是代数对象，\textbf{运算只对向量进行}。

坐标面与坐标轴上的点，其坐标有显著特征：$yOz$ 面上 $x=0$，$zOx$ 面上 $y=0$，$xOy$ 面上 $z=0$；$x$ 轴上 $y=z=0$，$y$ 轴上 $z=x=0$，$z$ 轴上 $x=y=0$；原点处 $x=y=z=0$。

\subsection{利用坐标作向量的线性运算}

设 $\boldsymbol{a}=(a_x,a_y,a_z)$，$\boldsymbol{b}=(b_x,b_y,b_z)$，$\lambda$ 为实数，则

\[
\boldsymbol{a}\pm\boldsymbol{b}=(a_x\pm b_x,\ a_y\pm b_y,\ a_z\pm b_z),\qquad
\lambda\boldsymbol{a}=(\lambda a_x,\ \lambda a_y,\ \lambda a_z).
\]

即对向量作加、减与数乘，只需对各个坐标分别作相应的数量运算。

设 $A(x_1,y_1,z_1)$、$B(x_2,y_2,z_2)$，则

\keyformula{\overrightarrow{AB}=(x_2-x_1,\ y_2-y_1,\ z_2-z_1)}

由平行的充要条件，当 $\boldsymbol{a}\neq\boldsymbol{0}$ 时

\keyformula{\boldsymbol{b}\parallel\boldsymbol{a}\iff\frac{b_x}{a_x}=\frac{b_y}{a_y}=\frac{b_z}{a_z}}

这里约定：某个分母为零时，相应的分子也必须为零。即\textbf{两向量平行等价于它们的坐标成比例}。

\textbf{分点坐标公式}\quad 设 $A(x_1,y_1,z_1)$、$B(x_2,y_2,z_2)$，点 $M$ 在直线 $AB$ 上且 $\overrightarrow{AM}=\lambda\overrightarrow{MB}$（$\lambda\neq-1$），则

\keyformula{\overrightarrow{OM}=\frac{\overrightarrow{OA}+\lambda\overrightarrow{OB}}{1+\lambda}}

即 $x=\dfrac{x_1+\lambda x_2}{1+\lambda}$，$y=\dfrac{y_1+\lambda y_2}{1+\lambda}$，$z=\dfrac{z_1+\lambda z_2}{1+\lambda}$。当 $\lambda=1$ 时得线段 $AB$ 的\textbf{中点}坐标 $\left(\dfrac{x_1+x_2}{2},\dfrac{y_1+y_2}{2},\dfrac{z_1+z_2}{2}\right)$。

\subsection{向量的模、方向角与投影}

由勾股定理，向量 $\boldsymbol{r}=(x,y,z)$ 的模

\keyformula{|\boldsymbol{r}|=\sqrt{x^{2}+y^{2}+z^{2}}}

从而点 $A(x_1,y_1,z_1)$ 与点 $B(x_2,y_2,z_2)$ 之间的距离为 $\overrightarrow{AB}$ 的模：

\keyformula{|AB|=\sqrt{(x_2-x_1)^{2}+(y_2-y_1)^{2}+(z_2-z_1)^{2}}}

\begin{definition}
非零向量 $\boldsymbol{r}$ 与三条坐标轴的夹角 $\alpha,\beta,\gamma$（$0\le\alpha,\beta,\gamma\le\pi$）称为 $\boldsymbol{r}$ 的\textbf{方向角}，$\cos\alpha,\cos\beta,\cos\gamma$ 称为 $\boldsymbol{r}$ 的\textbf{方向余弦}。
\end{definition}

\keyformula{\cos\alpha=\frac{x}{|\boldsymbol{r}|},\qquad
\cos\beta=\frac{y}{|\boldsymbol{r}|},\qquad
\cos\gamma=\frac{z}{|\boldsymbol{r}|}}

\keyformula{\cos^{2}\alpha+\cos^{2}\beta+\cos^{2}\gamma=1}

后一式表明：以方向余弦为坐标的向量正是与 $\boldsymbol{r}$ 同方向的单位向量，即
$\boldsymbol{e}_{\boldsymbol{r}}=(\cos\alpha,\ \cos\beta,\ \cos\gamma)$。

\begin{definition}
设数轴 $u$ 的单位向量为 $\boldsymbol{e}$，向量 $\boldsymbol{a}$ 与 $\boldsymbol{e}$ 的夹角为 $\varphi$，称 $|\boldsymbol{a}|\cos\varphi$ 为向量 $\boldsymbol{a}$ 在轴 $u$ 上的\textbf{投影}，记作 $\mathrm{Prj}_u\boldsymbol{a}$。
\end{definition}

投影具有线性性：

\[
\mathrm{Prj}_u(\boldsymbol{a}+\boldsymbol{b})=\mathrm{Prj}_u\boldsymbol{a}+\mathrm{Prj}_u\boldsymbol{b},\qquad
\mathrm{Prj}_u(\lambda\boldsymbol{a})=\lambda\,\mathrm{Prj}_u\boldsymbol{a}.
\]

向量的坐标就是它在各坐标轴上的投影：$a_x=\mathrm{Prj}_x\boldsymbol{a}$，$a_y=\mathrm{Prj}_y\boldsymbol{a}$，$a_z=\mathrm{Prj}_z\boldsymbol{a}$。投影是有符号的量，可正可负。

\corepoint{
（1）$\boldsymbol{a}\parallel\boldsymbol{b}$ 的坐标条件要写成“坐标成比例”，但\textbf{分母为零时必须单独讨论}：例如 $\boldsymbol{a}=(1,0,2)$、$\boldsymbol{b}=(2,0,4)$ 平行，此时 $\frac{0}{0}$ 无意义，只能写作 $b_x=2a_x,\ b_y=0,\ b_z=2a_z$。\\
（2）方向角的范围是 $[0,\pi]$，故 $\cos\alpha$ 可正可负；由 $\cos^{2}\alpha+\cos^{2}\beta+\cos^{2}\gamma=1$ 反求第三个方向余弦时，必须由点所在卦限（坐标符号）确定正负号，不能只取正根。\\
（3）记号 $(x,y,z)$ 兼指点与向量，须从上下文判断；对“点”不能作加减、数乘运算。
}

\section{数量积、向量积与混合积}

\subsection{两向量的数量积}

\begin{definition}
设 $\boldsymbol{a},\boldsymbol{b}$ 是两个向量，称 $|\boldsymbol{a}|\,|\boldsymbol{b}|\cos\langle\boldsymbol{a},\boldsymbol{b}\rangle$ 为 $\boldsymbol{a}$ 与 $\boldsymbol{b}$ 的\textbf{数量积}（内积、点积），记作 $\boldsymbol{a}\cdot\boldsymbol{b}$，即
\[
\boldsymbol{a}\cdot\boldsymbol{b}=|\boldsymbol{a}|\,|\boldsymbol{b}|\cos\langle\boldsymbol{a},\boldsymbol{b}\rangle .
\]
\end{definition}

物理背景：物体在恒力 $\boldsymbol{F}$ 作用下沿直线从 $M_1$ 移到 $M_2$，位移 $\boldsymbol{s}=\overrightarrow{M_1M_2}$，则力所做的功 $W=\boldsymbol{F}\cdot\boldsymbol{s}$。

设 $\boldsymbol{a}=(a_x,a_y,a_z)$，$\boldsymbol{b}=(b_x,b_y,b_z)$，由分配律与 $\boldsymbol{i}\cdot\boldsymbol{i}=\boldsymbol{j}\cdot\boldsymbol{j}=\boldsymbol{k}\cdot\boldsymbol{k}=1$、$\boldsymbol{i}\cdot\boldsymbol{j}=\boldsymbol{j}\cdot\boldsymbol{k}=\boldsymbol{k}\cdot\boldsymbol{i}=0$ 得坐标公式：

\keyformula{\boldsymbol{a}\cdot\boldsymbol{b}=a_xb_x+a_yb_y+a_zb_z}

\begin{property}
数量积满足下列运算律：
\begin{enumerate}
\item 交换律 $\boldsymbol{a}\cdot\boldsymbol{b}=\boldsymbol{b}\cdot\boldsymbol{a}$；
\item 分配律 $(\boldsymbol{a}+\boldsymbol{b})\cdot\boldsymbol{c}=\boldsymbol{a}\cdot\boldsymbol{c}+\boldsymbol{b}\cdot\boldsymbol{c}$；
\item 数乘结合律 $(\lambda\boldsymbol{a})\cdot\boldsymbol{b}=\lambda(\boldsymbol{a}\cdot\boldsymbol{b})$；
\item $\boldsymbol{a}\cdot\boldsymbol{a}=|\boldsymbol{a}|^{2}$，由此 $|\boldsymbol{a}|=\sqrt{\boldsymbol{a}\cdot\boldsymbol{a}}$。
\end{enumerate}
\end{property}

由定义与坐标公式即得两向量夹角的余弦：

\keyformula{\cos\langle\boldsymbol{a},\boldsymbol{b}\rangle=
\frac{\boldsymbol{a}\cdot\boldsymbol{b}}{|\boldsymbol{a}|\,|\boldsymbol{b}|}=
\frac{a_xb_x+a_yb_y+a_zb_z}
{\sqrt{a_x^{2}+a_y^{2}+a_z^{2}}\,\sqrt{b_x^{2}+b_y^{2}+b_z^{2}}}}

\begin{theorem}[垂直的充要条件]
设 $\boldsymbol{a},\boldsymbol{b}$ 为非零向量，则 $\boldsymbol{a}\perp\boldsymbol{b}\iff\boldsymbol{a}\cdot\boldsymbol{b}=0$。
\end{theorem}

数量积与投影的关系是：$\boldsymbol{a}\cdot\boldsymbol{b}=|\boldsymbol{a}|\,\mathrm{Prj}_{\boldsymbol{a}}\boldsymbol{b}=|\boldsymbol{b}|\,\mathrm{Prj}_{\boldsymbol{b}}\boldsymbol{a}$；特别当 $\boldsymbol{e}$ 为单位向量时 $\mathrm{Prj}_{\boldsymbol{e}}\boldsymbol{a}=\boldsymbol{a}\cdot\boldsymbol{e}$。

几何意义：数量积等于其中一个向量的模与另一个向量在这个向量方向上的投影的乘积。

\subsection{两向量的向量积}

\begin{definition}
对空间中任意两个向量 $\boldsymbol{a}$ 与 $\boldsymbol{b}$，规定它们的\textbf{向量积}（叉积、外积）$\boldsymbol{a}\times\boldsymbol{b}$ 如下：
\begin{enumerate}
\item 若 $\boldsymbol{a},\boldsymbol{b}$ 不平行，则 $\boldsymbol{a}\times\boldsymbol{b}$ 是一个向量，满足\\
\quad（i）$|\boldsymbol{a}\times\boldsymbol{b}|=|\boldsymbol{a}|\,|\boldsymbol{b}|\sin\langle\boldsymbol{a},\boldsymbol{b}\rangle$；\\
\quad（ii）$\boldsymbol{a}\times\boldsymbol{b}$ 垂直于 $\boldsymbol{a}$ 与 $\boldsymbol{b}$，且 $\boldsymbol{a},\boldsymbol{b},\boldsymbol{a}\times\boldsymbol{b}$ 依次构成\textbf{右手系}（右手四指从 $\boldsymbol{a}$ 转向 $\boldsymbol{b}$ 时，大拇指指向 $\boldsymbol{a}\times\boldsymbol{b}$ 的方向）；
\item 若 $\boldsymbol{a}\parallel\boldsymbol{b}$，规定 $\boldsymbol{a}\times\boldsymbol{b}=\boldsymbol{0}$。
\end{enumerate}
注意其中 $\boldsymbol{a}$ 与 $\boldsymbol{b}$ 的次序不能颠倒。
\end{definition}

坐标公式可按“$\boldsymbol{i},\boldsymbol{j},\boldsymbol{k}$ 作第一行、$\boldsymbol{a}$ 的坐标作第二行、$\boldsymbol{b}$ 的坐标作第三行”的三阶行列式展开：

\keyformula{\boldsymbol{a}\times\boldsymbol{b}=
\begin{vmatrix}
\boldsymbol{i}&\boldsymbol{j}&\boldsymbol{k}\\
a_x&a_y&a_z\\
b_x&b_y&b_z
\end{vmatrix}
=\begin{vmatrix}
a_y&a_z\\ b_y&b_z
\end{vmatrix}\boldsymbol{i}
-\begin{vmatrix}
a_x&a_z\\ b_x&b_z
\end{vmatrix}\boldsymbol{j}
+\begin{vmatrix}
a_x&a_y\\ b_x&b_y
\end{vmatrix}\boldsymbol{k}}

即 $\boldsymbol{a}\times\boldsymbol{b}=(a_yb_z-a_zb_y,\ a_zb_x-a_xb_z,\ a_xb_y-a_yb_x)$。展开时第二项带负号，这是最容易出错的地方。

\begin{property}
向量积满足：
\begin{enumerate}
\item 反交换律 $\boldsymbol{a}\times\boldsymbol{b}=-\boldsymbol{b}\times\boldsymbol{a}$，特别地 $\boldsymbol{a}\times\boldsymbol{a}=\boldsymbol{0}$；
\item $\boldsymbol{a}\times\boldsymbol{b}\perp\boldsymbol{a}$，$\boldsymbol{a}\times\boldsymbol{b}\perp\boldsymbol{b}$；
\item 分配律 $(\boldsymbol{a}+\boldsymbol{b})\times\boldsymbol{c}=\boldsymbol{a}\times\boldsymbol{c}+\boldsymbol{b}\times\boldsymbol{c}$；
\item 数乘结合 $(\lambda\boldsymbol{a})\times\boldsymbol{b}=\lambda(\boldsymbol{a}\times\boldsymbol{b})=\boldsymbol{a}\times(\lambda\boldsymbol{b})$；
\item $\boldsymbol{a}\parallel\boldsymbol{b}\iff\boldsymbol{a}\times\boldsymbol{b}=\boldsymbol{0}$（$\boldsymbol{a},\boldsymbol{b}$ 为非零向量）。
\end{enumerate}
\end{property}

几何意义：$|\boldsymbol{a}\times\boldsymbol{b}|$ 等于以 $\boldsymbol{a},\boldsymbol{b}$ 为邻边的平行四边形的面积；以 $\boldsymbol{a},\boldsymbol{b}$ 为邻边的三角形面积为 $\frac12|\boldsymbol{a}\times\boldsymbol{b}|$。物理上力矩 $\boldsymbol{M}=\boldsymbol{r}\times\boldsymbol{F}$，刚体绕定轴以角速度 $\boldsymbol{\omega}$ 转动时，点 $M$（向径 $\boldsymbol{r}$）的线速度 $\boldsymbol{v}=\boldsymbol{\omega}\times\boldsymbol{r}$。

\subsection{向量的混合积}

\begin{definition}
已知三个向量 $\boldsymbol{a},\boldsymbol{b},\boldsymbol{c}$，先作 $\boldsymbol{a}$ 与 $\boldsymbol{b}$ 的向量积 $\boldsymbol{a}\times\boldsymbol{b}$，再与 $\boldsymbol{c}$ 作数量积 $(\boldsymbol{a}\times\boldsymbol{b})\cdot\boldsymbol{c}$，所得的数量叫做这三个向量的\textbf{混合积}，记作 $[\boldsymbol{a}\boldsymbol{b}\boldsymbol{c}]$。
\end{definition}

\keyformula{[\boldsymbol{a}\boldsymbol{b}\boldsymbol{c}]=(\boldsymbol{a}\times\boldsymbol{b})\cdot\boldsymbol{c}=
\begin{vmatrix}
a_x&a_y&a_z\\
b_x&b_y&b_z\\
c_x&c_y&c_z
\end{vmatrix}}

\begin{property}
混合积满足：
\begin{enumerate}
\item \textbf{轮换对称性} $[\boldsymbol{a}\boldsymbol{b}\boldsymbol{c}]=[\boldsymbol{b}\boldsymbol{c}\boldsymbol{a}]=[\boldsymbol{c}\boldsymbol{a}\boldsymbol{b}]$；
\item 任意对换两个向量，混合积变号：$[\boldsymbol{a}\boldsymbol{b}\boldsymbol{c}]=-[\boldsymbol{b}\boldsymbol{a}\boldsymbol{c}]$；
\item $[\boldsymbol{a}\boldsymbol{b}\boldsymbol{c}]=0$ 的充要条件是 $\boldsymbol{a},\boldsymbol{b},\boldsymbol{c}$ \textbf{共面}（含其中有零向量或有两者平行的情况）。
\end{enumerate}
\end{property}

几何意义：$|[\boldsymbol{a}\boldsymbol{b}\boldsymbol{c}]|$ 等于以 $\boldsymbol{a},\boldsymbol{b},\boldsymbol{c}$ 为棱的平行六面体的体积；以它们为棱的四面体体积为 $\frac16|[\boldsymbol{a}\boldsymbol{b}\boldsymbol{c}]|$。混合积的\textbf{符号}由 $\boldsymbol{a},\boldsymbol{b},\boldsymbol{c}$ 构成右手系（为正）还是左手系（为负）决定。

\corepoint{
（1）数量积的结果是\textbf{数}，向量积的结果是\textbf{向量}，混合积的结果是\textbf{数}；写 $\boldsymbol{a}\cdot\boldsymbol{b}\cdot\boldsymbol{c}$ 没有意义，必须写成 $(\boldsymbol{a}\cdot\boldsymbol{b})\boldsymbol{c}$ 或 $[\boldsymbol{a}\boldsymbol{b}\boldsymbol{c}]$。\\
（2）向量积\textbf{不满足结合律}：一般地 $(\boldsymbol{a}\times\boldsymbol{b})\times\boldsymbol{c}\neq\boldsymbol{a}\times(\boldsymbol{b}\times\boldsymbol{c})$，只有分配律可用。\\
（3）向量积坐标公式的第二个分量带负号，务必按行列式展开，不要凭记忆写 $(a_yb_z-a_zb_y,\ a_zb_x-a_xb_z,\ \dots)$ 时把符号弄反。\\
（4）用混合积判断共面、求体积时，结论是“混合积为零”与“取绝对值”，符号本身只表示左右手系。
}

\section{平面及其方程}

\subsection{平面的点法式方程}

\begin{definition}
垂直于平面 $\pi$ 的非零向量 $\boldsymbol{n}=(A,B,C)$ 称为平面 $\pi$ 的\textbf{法向量}。
\end{definition}

设平面 $\pi$ 过点 $M_0(x_0,y_0,z_0)$，且以 $\boldsymbol{n}=(A,B,C)\neq\boldsymbol{0}$ 为法向量。点 $M(x,y,z)$ 在 $\pi$ 上的充要条件是 $\overrightarrow{M_0M}\perp\boldsymbol{n}$，即 $\boldsymbol{n}\cdot\overrightarrow{M_0M}=0$，于是

\keyformula{A(x-x_0)+B(y-y_0)+C(z-z_0)=0}

称其为平面 $\pi$ 的\textbf{点法式方程}。几何上“一点 + 一个法向量”唯一确定一张平面。

\subsection{平面的一般方程}

把点法式展开，得

\[
Ax+By+Cz+D=0\qquad(A,B,C\ \text{不全为零}),
\]

称为平面的\textbf{一般方程}；其法向量就是 $\boldsymbol{n}=(A,B,C)$。反之，任何一个三元一次方程都表示一张平面。三元一次方程中三个系数 $A,B,C$ 的取值决定了平面的位置，常见特殊情形如下。

\begin{table}[htbp]
\centering
\renewcommand{\tablename}{表}
\caption{平面一般方程 $Ax+By+Cz+D=0$ 的特殊位置}
\begin{tabular}{@{}ll@{}}
\toprule
方程特征 & 图形 \\
\midrule
$D=0$ & 过原点 \\
$A=0$ & 平行于 $x$ 轴 \\
$A=D=0$ & 通过 $x$ 轴 \\
$A=B=0$ & 平行于 $xOy$ 面（即 $z=-\frac{D}{C}$）\\
$A=B=D=0$ & 就是 $xOy$ 面（$z=0$）\\
\bottomrule
\end{tabular}
\end{table}

若平面与三条坐标轴分别交于 $P(a,0,0)$、$Q(0,b,0)$、$R(0,0,c)$，且 $abc\neq0$，则平面方程可写成

\keyformula{\frac{x}{a}+\frac{y}{b}+\frac{z}{c}=1}

称为平面的\textbf{截距式方程}，$a,b,c$ 依次称为平面在 $x$ 轴、$y$ 轴、$z$ 轴上的\textbf{截距}。截距式便于快速作图，但过原点或平行于某坐标轴的平面没有截距式。

\textbf{三点式}\quad 过不共线三点 $M_1,M_2,M_3$ 的平面，取法向量 $\boldsymbol{n}=\overrightarrow{M_1M_2}\times\overrightarrow{M_1M_3}$，再按点法式写出即可；三点共线的判据是 $\overrightarrow{M_1M_2}\times\overrightarrow{M_1M_3}=\boldsymbol{0}$。

\subsection{两平面的夹角}

\begin{definition}
两平面的\textbf{夹角}规定为它们法向量的夹角（取锐角或直角），即 $\theta\in\left[0,\frac{\pi}{2}\right]$。
\end{definition}

设 $\pi_1:A_1x+B_1y+C_1z+D_1=0$，$\pi_2:A_2x+B_2y+C_2z+D_2=0$，法向量分别为 $\boldsymbol{n}_1,\boldsymbol{n}_2$，则

\keyformula{\cos\theta=\frac{|\boldsymbol{n}_1\cdot\boldsymbol{n}_2|}{|\boldsymbol{n}_1|\,|\boldsymbol{n}_2|}=
\frac{|A_1A_2+B_1B_2+C_1C_2|}
{\sqrt{A_1^{2}+B_1^{2}+C_1^{2}}\,\sqrt{A_2^{2}+B_2^{2}+C_2^{2}}}}

两平面的位置关系由法向量与常数项共同决定：

\begin{table}[htbp]
\centering
\renewcommand{\tablename}{表}
\caption{两平面位置关系的判据}
\begin{tabular}{@{}ll@{}}
\toprule
位置关系 & 充要条件 \\
\midrule
$\pi_1\perp\pi_2$ & $A_1A_2+B_1B_2+C_1C_2=0$ \\
$\pi_1\parallel\pi_2$（不重合） & $\dfrac{A_1}{A_2}=\dfrac{B_1}{B_2}=\dfrac{C_1}{C_2}\neq\dfrac{D_1}{D_2}$ \\
$\pi_1$ 与 $\pi_2$ 重合 & $\dfrac{A_1}{A_2}=\dfrac{B_1}{B_2}=\dfrac{C_1}{C_2}=\dfrac{D_1}{D_2}$ \\
\bottomrule
\end{tabular}
\end{table}

\subsection{点到平面的距离}

设 $P_0(x_0,y_0,z_0)$ 是平面 $\pi:Ax+By+Cz+D=0$ 外一点，在 $\pi$ 上任取一点 $P_1$，则 $\overrightarrow{P_1P_0}$ 在法向量 $\boldsymbol{n}$ 上的投影的绝对值就是所求距离：

\keyformula{d=\frac{|Ax_0+By_0+Cz_0+D|}{\sqrt{A^{2}+B^{2}+C^{2}}}}

推导要点：$d=\left|\mathrm{Prj}_{\boldsymbol{n}}\overrightarrow{P_1P_0}\right|=\dfrac{|\boldsymbol{n}\cdot\overrightarrow{P_1P_0}|}{|\boldsymbol{n}|}$，而 $\boldsymbol{n}\cdot\overrightarrow{P_1P_0}=A(x_0-x_1)+B(y_0-y_1)+C(z_0-z_1)=Ax_0+By_0+Cz_0+D$（用了 $P_1$ 在平面上）。

由此可得两\textbf{平行}平面 $Ax+By+Cz+D_1=0$ 与 $Ax+By+Cz+D_2=0$ 之间的距离

\keyformula{d=\frac{|D_1-D_2|}{\sqrt{A^{2}+B^{2}+C^{2}}}}

\corepoint{
（1）点到平面距离公式中分子要加\textbf{绝对值}，否则得到的是有符号的“离差”，不能当距离。\\
（2）用平行平面距离公式前，必须先把两方程化成 $x,y,z$ 系数\textbf{完全相同}的形式：若给的是 $2x+4y-2z+1=0$ 与 $x+2y-z+3=0$，须先把后者乘 $2$（得 $2x+4y-2z+6=0$）再套公式。\\
（3）平面一般方程的三个系数不能同时为零；若只有一个系数为零（如 $A=0$），平面仍存在，只是平行于相应的坐标轴，不要误判为“退化”。
}

\section{空间直线及其方程}

\subsection{空间直线的一般方程}

空间直线可看作两张不平行的平面的交线。设

\[
\pi_1:A_1x+B_1y+C_1z+D_1=0,\qquad
\pi_2:A_2x+B_2y+C_2z+D_2=0,
\]

且 $\boldsymbol{n}_1=(A_1,B_1,C_1)$ 与 $\boldsymbol{n}_2=(A_2,B_2,C_2)$ 不平行（即 $A_1:B_1:C_1\neq A_2:B_2:C_2$），则方程组

\[
\begin{cases}
A_1x+B_1y+C_1z+D_1=0,\\
A_2x+B_2y+C_2z+D_2=0
\end{cases}
\]

表示一条空间直线，称为空间直线的\textbf{一般方程}。同一条直线可写成无数种一般方程（取不同的两张过它的平面），故一般方程不唯一。

\subsection{对称式方程与参数方程}

\begin{definition}
平行于直线 $L$ 的非零向量 $\boldsymbol{s}=(m,n,p)$ 称为 $L$ 的\textbf{方向向量}，$m,n,p$ 称为 $L$ 的\textbf{方向数}。
\end{definition}

设直线 $L$ 过点 $M_0(x_0,y_0,z_0)$，方向向量为 $\boldsymbol{s}=(m,n,p)$，则点 $M(x,y,z)$ 在 $L$ 上的充要条件是 $\overrightarrow{M_0M}\parallel\boldsymbol{s}$，即存在 $t$ 使 $\overrightarrow{M_0M}=t\boldsymbol{s}$。消去 $t$ 得

\keyformula{\frac{x-x_0}{m}=\frac{y-y_0}{n}=\frac{z-z_0}{p}}

称为直线的\textbf{对称式方程}（点向式方程）。约定：某个分母为零时，相应的分子也为零，例如 $m=0$ 时该式应理解为 $x=x_0$ 与 $\dfrac{y-y_0}{n}=\dfrac{z-z_0}{p}$ 联立。

保留参数 $t$ 即得\textbf{参数方程}

\keyformula{\begin{cases}
x=x_0+mt,\\
y=y_0+nt,\\
z=z_0+pt
\end{cases}\qquad(-\infty<t<+\infty)}

过两点 $M_1(x_1,y_1,z_1)$、$M_2(x_2,y_2,z_2)$ 的直线，取 $\boldsymbol{s}=\overrightarrow{M_1M_2}$，得\textbf{两点式}

\keyformula{\frac{x-x_1}{x_2-x_1}=\frac{y-y_1}{y_2-y_1}=\frac{z-z_1}{z_2-z_1}}

\subsection{由一般方程求对称式的步骤}

\begin{enumerate}
\item 求方向向量：$\boldsymbol{s}=\boldsymbol{n}_1\times\boldsymbol{n}_2$，它同时垂直于两个平面的法向量，故平行于交线；
\item 求直线上一点：在一般方程中令某个坐标（通常取使计算最简的那个，如 $z=0$）为常数，解出其余两个坐标，得 $M_0$；
\item 写出对称式；若 $\boldsymbol{s}$ 的某个分量为零，按约定写成“分子为零”的形式，不要写成除以零。
\end{enumerate}

\section{位置关系与距离}

\subsection{两直线的位置关系}

设
\[
L_1:\frac{x-x_1}{m_1}=\frac{y-y_1}{n_1}=\frac{z-z_1}{p_1},\qquad
L_2:\frac{x-x_2}{m_2}=\frac{y-y_2}{n_2}=\frac{z-z_2}{p_2},
\]
方向向量分别为 $\boldsymbol{s}_1=(m_1,n_1,p_1)$、$\boldsymbol{s}_2=(m_2,n_2,p_2)$，$M_1\in L_1$，$M_2\in L_2$。两直线的位置关系由 $\boldsymbol{s}_1,\boldsymbol{s}_2,\overrightarrow{M_1M_2}$ 三者的关系完全确定：

\begin{enumerate}
\item $L_1$ 与 $L_2$ \textbf{重合}：$\boldsymbol{s}_1\parallel\boldsymbol{s}_2$ 且 $M_1\in L_2$（即 $\overrightarrow{M_1M_2}\parallel\boldsymbol{s}_2$）；
\item $L_1\parallel L_2$（不重合）：$\boldsymbol{s}_1\parallel\boldsymbol{s}_2$ 但 $\overrightarrow{M_1M_2}$ 不平行于 $\boldsymbol{s}_2$；
\item $L_1$ 与 $L_2$ \textbf{相交}于一点：$\boldsymbol{s}_1,\boldsymbol{s}_2$ 不平行，且混合积 $[\boldsymbol{s}_1\boldsymbol{s}_2\overrightarrow{M_1M_2}]=0$（三向量共面）；
\item $L_1$ 与 $L_2$ \textbf{异面}：混合积 $[\boldsymbol{s}_1\boldsymbol{s}_2\overrightarrow{M_1M_2}]\neq0$（三向量不共面）。
\end{enumerate}

即：\textbf{先看方向向量是否平行，再用混合积是否为零区分“相交”与“异面”}。

\subsection{两直线的夹角}

两直线的夹角 $\theta$ 规定为它们方向向量的夹角，取锐角或直角，$\theta\in\left[0,\frac{\pi}{2}\right]$：

\keyformula{\cos\theta=\frac{|\boldsymbol{s}_1\cdot\boldsymbol{s}_2|}{|\boldsymbol{s}_1|\,|\boldsymbol{s}_2|}=
\frac{|m_1m_2+n_1n_2+p_1p_2|}
{\sqrt{m_1^{2}+n_1^{2}+p_1^{2}}\,\sqrt{m_2^{2}+n_2^{2}+p_2^{2}}}}

由此：$L_1\perp L_2\iff m_1m_2+n_1n_2+p_1p_2=0$；$L_1\parallel L_2\iff\dfrac{m_1}{m_2}=\dfrac{n_1}{n_2}=\dfrac{p_1}{p_2}$（分母为零时按约定理解）。

\subsection{直线与平面的位置关系}

设直线 $L:\dfrac{x-x_0}{m}=\dfrac{y-y_0}{n}=\dfrac{z-z_0}{p}$，平面 $\pi:Ax+By+Cz+D=0$，方向向量 $\boldsymbol{s}=(m,n,p)$，法向量 $\boldsymbol{n}=(A,B,C)$。

\begin{enumerate}
\item $L$ \textbf{在平面} $\pi$ 上：$\boldsymbol{s}\cdot\boldsymbol{n}=0$ 且 $Ax_0+By_0+Cz_0+D=0$；
\item $L\parallel\pi$（不落在平面上）：$\boldsymbol{s}\cdot\boldsymbol{n}=0$ 而 $Ax_0+By_0+Cz_0+D\neq0$；
\item $L$ 与 $\pi$ \textbf{相交}于一点：$\boldsymbol{s}\cdot\boldsymbol{n}\neq0$，交点的求法是把 $L$ 的参数方程代入 $\pi$ 的方程解出 $t$，再回代；
\item $L\perp\pi$：$\boldsymbol{s}\parallel\boldsymbol{n}$，即 $\dfrac{m}{A}=\dfrac{n}{B}=\dfrac{p}{C}$。
\end{enumerate}

\begin{definition}
直线 $L$ 与它在平面 $\pi$ 上的\textbf{投影}所成的夹角 $\varphi$（$0\le\varphi\le\frac{\pi}{2}$）称为直线与平面的\textbf{夹角}。
\end{definition}

由于 $\boldsymbol{s}$ 与 $\boldsymbol{n}$ 的夹角与 $\varphi$ 互余（或为其补角），故用\textbf{正弦}：

\keyformula{\sin\varphi=\frac{|\boldsymbol{s}\cdot\boldsymbol{n}|}{|\boldsymbol{s}|\,|\boldsymbol{n}|}=
\frac{|Am+Bn+Cp|}
{\sqrt{A^{2}+B^{2}+C^{2}}\,\sqrt{m^{2}+n^{2}+p^{2}}}}

\corepoint{
（1）直线与平面夹角用 $\sin$，两直线夹角与两平面夹角用 $\cos$，三者都取绝对值使夹角落在 $\left[0,\frac{\pi}{2}\right]$；弄混是最高频的错误。\\
（2）“$\boldsymbol{s}\cdot\boldsymbol{n}=0$”只能说明直线平行于平面\textbf{或落在平面上}，必须再代一个点检验，才能区分这两种情形。\\
（3）两直线“方向向量不平行 + 混合积为零”才是相交；只验证混合积为零而忽略平行情形，会把重合或平行的直线误判为相交。
}

\subsection{点到直线的距离}

设点 $P(x_0,y_0,z_0)$ 与直线 $L:\dfrac{x-x_1}{m}=\dfrac{y-y_1}{n}=\dfrac{z-z_1}{p}$（$M_1(x_1,y_1,z_1)\in L$，方向向量 $\boldsymbol{s}=(m,n,p)$）。以 $\overrightarrow{M_1P}$ 与 $\boldsymbol{s}$ 为邻边的平行四边形，其面积为 $|\boldsymbol{s}\times\overrightarrow{M_1P}|$，而该面积又等于 $|\boldsymbol{s}|\cdot d$，故

\keyformula{d=\frac{|\boldsymbol{s}\times\overrightarrow{M_1P}|}{|\boldsymbol{s}|}}

若 $L$ 以一般方程给出，先用 $\boldsymbol{s}=\boldsymbol{n}_1\times\boldsymbol{n}_2$ 求方向向量，并在 $L$ 上任取一点作 $M_1$。

\subsection{两平行直线与两异面直线间的距离}

\textbf{两平行直线间的距离}\quad 设 $L_1\parallel L_2$，公共方向向量 $\boldsymbol{s}$，$M_1\in L_1$，$M_2\in L_2$，则

\keyformula{d=\frac{|\boldsymbol{s}\times\overrightarrow{M_1M_2}|}{|\boldsymbol{s}|}}

与点到直线距离公式同形：在 $L_2$ 上任取一点，化为“点到直线”的情形。

\textbf{两异面直线间的距离}\quad 设
\[
L_1:\frac{x-x_1}{m_1}=\frac{y-y_1}{n_1}=\frac{z-z_1}{p_1},\qquad
L_2:\frac{x-x_2}{m_2}=\frac{y-y_2}{n_2}=\frac{z-z_2}{p_2},
\]
$\boldsymbol{s}_1=(m_1,n_1,p_1)$，$\boldsymbol{s}_2=(m_2,n_2,p_2)$，$M_1(x_1,y_1,z_1)$，$M_2(x_2,y_2,z_2)$。两异面直线的公垂线方向为 $\boldsymbol{s}_1\times\boldsymbol{s}_2$，$\overrightarrow{M_1M_2}$ 在该方向上的投影的绝对值就是 $L_1$ 与 $L_2$ 之间的距离：

\keyformula{d=\frac{|[\boldsymbol{s}_1\boldsymbol{s}_2\overrightarrow{M_1M_2}]|}{|\boldsymbol{s}_1\times\boldsymbol{s}_2|}}

注意：此式要求 $\boldsymbol{s}_1\times\boldsymbol{s}_2\neq\boldsymbol{0}$（即两直线不平行）；若两直线平行，公垂线方向不唯一，应改用上一式。

\corepoint{
（1）求距离的三条公式本质相同，都是“面积（或体积）除以底”：点到直线用平行四边形面积 $|\boldsymbol{s}\times\overrightarrow{M_1P}|$，异面直线用平行六面体体积 $|[\boldsymbol{s}_1\boldsymbol{s}_2\overrightarrow{M_1M_2}]|$。记住这个几何来源，就不必死记公式。\\
（2）异面直线距离公式的分母是 $|\boldsymbol{s}_1\times\boldsymbol{s}_2|$ 而不是 $|\boldsymbol{s}_1|\,|\boldsymbol{s}_2|$，两者相差 $\sin\langle\boldsymbol{s}_1,\boldsymbol{s}_2\rangle$。\\
（3）所有距离公式最后都要取绝对值；用混合积算体积时若得负数，只说明三向量成左手系，距离仍取绝对值。
}

\section{平面束}

\subsection{平面束的方程}

通过给定直线 $L$ 的所有平面的全体，称为通过 $L$ 的\textbf{平面束}。设 $L$ 由一般方程

\[
\pi_1:A_1x+B_1y+C_1z+D_1=0,\qquad
\pi_2:A_2x+B_2y+C_2z+D_2=0
\]

给出，且 $\boldsymbol{n}_1=(A_1,B_1,C_1)$ 与 $\boldsymbol{n}_2=(A_2,B_2,C_2)$ 不平行（即 $\pi_1,\pi_2$ 确实交于一条直线）。对任意实数 $\lambda$，方程

\[
A_1x+B_1y+C_1z+D_1+\lambda\bigl(A_2x+B_2y+C_2z+D_2\bigr)=0
\]

仍是 $x,y,z$ 的一次方程，因而表示一张平面；又因为 $L$ 上每点同时满足两个括号为零，所以 $L$ 上的点都满足它。于是它表示过 $L$ 的平面。反之，除 $\pi_2$ 以外的任一过 $L$ 的平面都可由某个 $\lambda$ 表示，故

\keyformula{\pi_1+\lambda\pi_2=0\quad(\lambda\in\mathbb{R})}

称为过直线 $L$ 的\textbf{平面束方程}（单参数形式）。若需要把 $\pi_2$ 本身也包含进来，改用两个参数 $\lambda,\mu$ 不同时为零的齐次组合

\keyformula{\lambda\bigl(A_1x+B_1y+C_1z+D_1\bigr)+\mu\bigl(A_2x+B_2y+C_2z+D_2\bigr)=0}

此时表示过 $L$ 的\textbf{全部}平面。

\subsection{用平面束求平面方程的步骤}

\begin{enumerate}
\item 写出平面束方程 $\pi_1+\lambda\pi_2=0$，其中 $\lambda$ 待定；
\item 把题设条件翻译成关于 $\lambda$ 的方程：过某点 $\Rightarrow$ 代入该点坐标为零；垂直于某平面 $\Rightarrow$ 两法向量数量积为零；平行于某向量 $\Rightarrow$ 法向量与该向量数量积为零；与某平面成定角 $\Rightarrow$ 用法向量夹角余弦公式；
\item 解出 $\lambda$ 后代回，化简即得所求平面；
\item 若步骤 2 得到的方程无解或恒成立，说明所求平面正是 $\pi_2$，需单独取出 $\pi_2$ 检验。
\end{enumerate}

\corepoint{
（1）单参数形式 $\pi_1+\lambda\pi_2=0$ \textbf{永远漏掉} $\pi_2$ 这一张平面；若所求平面的条件恰好只被 $\pi_2$ 满足，用单参数式会“无解”，此时应写双参数形式或直接验证 $\pi_2$。\\
（2）平面束只对\textbf{相交}的两平面成立。若 $\boldsymbol{n}_1\parallel\boldsymbol{n}_2$，则 $\pi_1+\lambda\pi_2=0$ 表示与这两平面\textbf{平行}的一族平面，不是过同一直线的平面束，判据是有没有“过同一条直线”。\\
（3）$\pi_1,\pi_2$ 的取法不唯一（同一条直线可用不同两张平面表示），取计算最简便的一组，例如让其中一个平面的某个系数为零。
}

\section{常见曲面}

\subsection{曲面方程的概念}

\begin{definition}
若曲面 $S$ 上每一点的坐标都满足方程 $F(x,y,z)=0$，而不在 $S$ 上的点的坐标都不满足这个方程，则称 $F(x,y,z)=0$ 为曲面 $S$ 的\textbf{方程}，$S$ 为该方程的\textbf{图形}。
\end{definition}

曲面研究有两个基本问题：\textbf{已知曲面求方程}——用轨迹法，把几何条件“翻译”成坐标的等式；\textbf{已知方程求图形}——用截痕法，靠截线形状推断曲面形状。空间曲线则看作两张曲面的交线，其一般方程为两个方程联立的方程组。

\subsection{旋转曲面}

\begin{definition}
一条平面曲线 $C$ 绕它所在平面内的一条定直线 $L$ 旋转一周所成的曲面，叫做\textbf{旋转曲面}；$C$ 称为\textbf{母线}，$L$ 称为\textbf{旋转轴}。
\end{definition}

设母线 $C$ 在 $yOz$ 面上，方程为 $f(y,z)=0$。把 $C$ 绕 $z$ 轴旋转，则旋转过程中 $z$ 坐标不变，而点到 $z$ 轴的距离 $\sqrt{x^{2}+y^{2}}$ 保持等于原来的 $|y|$，故旋转曲面的方程为

\keyformula{f\bigl(\pm\sqrt{x^{2}+y^{2}},\ z\bigr)=0}

同理，$yOz$ 面上的曲线 $f(y,z)=0$ 绕 $y$ 轴旋转所得曲面为 $f\bigl(y,\ \pm\sqrt{x^{2}+z^{2}}\bigr)=0$。

记忆口诀：\textbf{与旋转轴同名的变量不动，另一个变量换成另外两个变量平方和的平方根}（正负号成对出现，平方后往往自动消去）。例如 $yOz$ 面上的抛物线 $y^{2}=2pz$ 绕 $z$ 轴旋转，得旋转抛物面 $x^{2}+y^{2}=2pz$；$yOz$ 面上的椭圆 $\dfrac{y^{2}}{b^{2}}+\dfrac{z^{2}}{c^{2}}=1$ 绕 $z$ 轴旋转，得旋转椭球面 $\dfrac{x^{2}+y^{2}}{b^{2}}+\dfrac{z^{2}}{c^{2}}=1$。

\subsection{柱面}

\begin{definition}
平行于定直线并沿定曲线 $C$ 移动的直线 $L$ 所形成的曲面叫做\textbf{柱面}，$C$ 称为\textbf{准线}，$L$ 称为\textbf{母线}。
\end{definition}

\textbf{判据}\quad 若曲面方程中\emph{只含两个变量}（例如只含 $x,y$），则它表示柱面，母线平行于\emph{方程中不出现的那个变量}所对应的坐标轴（此例中母线平行于 $z$ 轴），该柱面在 $xOy$ 面上的截痕就是它的准线。

常见柱面：$\dfrac{x^{2}}{a^{2}}+\dfrac{y^{2}}{b^{2}}=1$ 为椭圆柱面；$\dfrac{x^{2}}{a^{2}}-\dfrac{y^{2}}{b^{2}}=1$ 为双曲柱面；$y^{2}=2px$ 为抛物柱面；$x^{2}+y^{2}=R^{2}$ 为圆柱面。它们的母线都平行于 $z$ 轴。

\subsection{二次曲面与截痕法}

\begin{definition}
三元二次方程所表示的曲面叫做\textbf{二次曲面}。通过适当的坐标变换（平移、旋转），二次曲面可化为若干种\textbf{标准方程}之一。
\end{definition}

\textbf{截痕法}是识别二次曲面类型的基本方法：

\begin{enumerate}
\item 判断对称性：方程中 $x,y,z$ 都以平方项（或一次项）出现时，曲面关于相应的坐标面、坐标轴、原点对称；
\item 求与三个坐标面的截痕：分别令 $x=0$、$y=0$、$z=0$，看所得平面曲线的类型；
\item 用平行平面截割：令 $z=h$（或 $x=h$、$y=h$），观察截线类型随 $h$ 的变化（椭圆的大小、双曲线的实轴方向、是否退化为直线或点）；
\item 综合对称性、截线类型与变化趋势，定出曲面名称与大致图形（是否有界、是否分叶、是否有直纹）。
\end{enumerate}

\subsection{二次曲面标准方程对照}

\begingroup
\renewcommand{\tablename}{表}
\setlength{\LTcapwidth}{\linewidth}
\begin{longtable}{@{}llp{0.42\textwidth}@{}}
\caption{常见二次曲面的标准方程与截痕特征}\\
\toprule
名称 & 标准方程 & 截痕特征 \\
\midrule
\endhead
\bottomrule
\endlastfoot
椭球面 & $\dfrac{x^{2}}{a^{2}}+\dfrac{y^{2}}{b^{2}}+\dfrac{z^{2}}{c^{2}}=1$ & 三组平行平面截得椭圆，有界（$|x|\le a$ 等） \\
单叶双曲面 & $\dfrac{x^{2}}{a^{2}}+\dfrac{y^{2}}{b^{2}}-\dfrac{z^{2}}{c^{2}}=1$ & 平行 $xOy$ 面截得椭圆且半轴随 $|h|$ 增大；平行 $yOz$ 面截得双曲线 \\
双叶双曲面 & $\dfrac{x^{2}}{a^{2}}-\dfrac{y^{2}}{b^{2}}-\dfrac{z^{2}}{c^{2}}=1$ & 分两叶（$|x|\ge a$）；$|x|>a$ 时截得椭圆 \\
椭圆抛物面 & $\dfrac{x^{2}}{2p}+\dfrac{y^{2}}{2q}=z\ (p,q\ \text{同号})$ & 平行 $xOy$ 面截得椭圆，平行 $yOz$ 面截得抛物线；无界 \\
双曲抛物面 & $\dfrac{x^{2}}{2p}-\dfrac{y^{2}}{2q}=z\ (p,q\ \text{同号})$ & 马鞍形；$z=h\neq0$ 截得双曲线，$z=0$ 截得两条相交直线 \\
二次锥面 & $\dfrac{x^{2}}{a^{2}}+\dfrac{y^{2}}{b^{2}}-\dfrac{z^{2}}{c^{2}}=0$ & 二次齐次；顶点在原点，$z=h\neq0$ 截得椭圆 \\
\end{longtable}
\endgroup

表中六个是真正“弯曲”的二次曲面；椭圆柱面、双曲柱面、抛物柱面这三类（也属二次曲面）的标准方程与母线方向见上一小节。椭球面、单叶双曲面、双曲抛物面、二次锥面还具有\textbf{直纹性}（曲面可由直线族生成），这是工程上（冷却塔、薄壳屋顶）常用的性质。用二次型的正交变换系统判别曲面类型，见第5章「拓展：空间中的曲面与二次曲面」一节。

\corepoint{
（1）旋转曲面公式中“与旋转轴同名的变量不动”，绕 $z$ 轴旋转时是 $y$ 换成 $\pm\sqrt{x^{2}+y^{2}}$，不是 $z$ 换；写错会得到完全不同的曲面。\\
（2）“方程只含两个变量 $\Rightarrow$ 柱面”这条判据里的“缺少的变量”决定母线方向：缺 $z$ 则母线平行 $z$ 轴。\\
（3）$\dfrac{x^{2}}{a^{2}}+\dfrac{y^{2}}{b^{2}}-\dfrac{z^{2}}{c^{2}}=1$ 与 $=0$ 只差一个常数项，图形却分别是单叶双曲面与二次锥面，判断时必须看清右端常数。\\
（4）截痕法只作定性判断；若要精确判定类型，应把二次型化为标准形，由特征值的符号与个数（惯性指数）分类。
}


\chapter{拓展与补充}
% file: 11.拓展.tex
% 作为主文档 1.main.tex 的子文件（% file: 11.拓展.tex
% 作为主文档 1.main.tex 的子文件（% file: 11.拓展.tex
% 作为主文档 1.main.tex 的子文件（\input{11.拓展.tex}），不要写 \documentclass、\begin{document}、\end{document}
% 附录B 拓展与补充：整理原残稿（伴随矩阵的秩、降阶公式），并把原稿空缺的几处
% 按同一主题写实（秩与线性相关性、生成子空间、二次曲面判别、向量方法）。
% 与教材重复的部分只作指引，不重复推导。不写例题与习题。

\section{矩阵秩的补充结论}
\noindent\textit{本节的结论第3、5章都已用到，但教材未单独成节，这里集中补出，并说明条件与推导思路。}

\subsection{伴随矩阵的秩}

\begin{theorem}
设 $A$ 为 $n$ 阶方阵（$n\ge2$），$A^{*}$ 为其伴随矩阵，则
\[
r(A^{*})=
\begin{cases}
n, & r(A)=n,\\[1mm]
1, & r(A)=n-1,\\[1mm]
0, & r(A)\le n-2.
\end{cases}
\]
\end{theorem}

推导思路（三段都要用 $AA^{*}=A^{*}A=|A|E$）：

\begin{enumerate}
\item $r(A)=n$：$A$ 可逆，$|A|\neq0$，由 $A^{*}=|A|A^{-1}$ 知 $A^{*}$ 可逆，故 $r(A^{*})=n$；
\item $r(A)=n-1$：$|A|=0$，故 $AA^{*}=O$，由 $r(A)+r(A^{*})\le n$ 得 $r(A^{*})\le1$；又此时 $A$ 有非零的 $n-1$ 阶子式，即 $A^{*}\neq O$，故 $r(A^{*})=1$；
\item $r(A)\le n-2$：所有 $n-1$ 阶子式为零，$A^{*}=O$，故 $r(A^{*})=0$。
\end{enumerate}

\corepoint{
（1）中间一档最容易记成 $n-1$：$r(A)=n-1$ 时 $A^{*}$ 的秩\textbf{恒为 $1$}，与 $n$ 无关。\\
（2）结论的三个分界，正是“最高阶非零子式的阶数”这一秩的定义的直接体现；$n=1$ 时约定 $A^{*}=1$，公式不适用。
}

\subsection{降阶公式}

\begin{theorem}
设 $A$ 是 $m\times n$ 矩阵，$B$ 是 $n\times m$ 矩阵，$m>n$，$\lambda$ 为任意数，则
\[
|\lambda E_m-AB|=\lambda^{m-n}\,|\lambda E_n-BA|.
\]
\end{theorem}

即 $AB$（$m$ 阶）与 $BA$（$n$ 阶）的特征多项式只差因子 $\lambda^{m-n}$，故它们的\textbf{非零特征值完全相同（含重数）}。又 $r(AB)\le n<m$，$AB$ 必奇异，$0$ 是 $AB$ 的特征值，而 $BA$ 未必以 $0$ 为特征值——差别正好被 $\lambda^{m-n}$ 吸收。

\begin{corollary}
取 $\lambda=1$ 得 $|E_m-AB|=|E_n-BA|$，可把高阶行列式化为低阶行列式计算。
\end{corollary}

\corepoint{
（1）公式要求 $m>n$，且乘法方向不能颠倒：$AB$ 是 $m$ 阶、$BA$ 是 $n$ 阶。\\
（2）$\lambda=0$ 时两边同为 $0$，等式成立但不提供信息；求特征值须取 $\lambda\neq0$ 讨论。$m=n$ 时退化为恒等式。
}

\section{线性代数与解析几何的接口}

\subsection{秩与线性相关性}

矩阵的秩、向量组的秩、线性相关性三者是同一件事的三种说法。

\begin{theorem}
设 $A=(\boldsymbol{a}_1,\dots,\boldsymbol{a}_s)$ 是以向量组为列的矩阵，则
\begin{enumerate}
\item $\boldsymbol{a}_1,\dots,\boldsymbol{a}_s$ 线性相关 $\iff r(A)<s$；线性无关 $\iff r(A)=s$；
\item $r(\boldsymbol{a}_1,\dots,\boldsymbol{a}_s)=r(A)$，即极大线性无关组所含向量的个数；
\item 当向量个数等于维数（$s=n$）时，线性相关 $\iff|A|=0$。
\end{enumerate}
\end{theorem}

搬到三维空间，就得到解析几何中一串“行列式判零”的判据。设 $\boldsymbol{a},\boldsymbol{b},\boldsymbol{c}$ 为三维向量，则

\keyformula{\boldsymbol{a},\boldsymbol{b},\boldsymbol{c}\ \text{共面}\iff[\boldsymbol{a}\boldsymbol{b}\boldsymbol{c}]=0\iff
\begin{vmatrix}
a_x&a_y&a_z\\ b_x&b_y&b_z\\ c_x&c_y&c_z
\end{vmatrix}=0\iff r(\boldsymbol{a},\boldsymbol{b},\boldsymbol{c})<3}

因为混合积正是以三向量为行的三阶行列式（见附录A），而“共面”就是“线性相关”。同理：三点共线 $\iff$ 两方向向量线性相关（秩为 $1$）；四点 $M_0,M_1,M_2,M_3$ 共面 $\iff\overrightarrow{M_0M_1},\overrightarrow{M_0M_2},\overrightarrow{M_0M_3}$ 的行列式为零；两直线异面 $\iff r(\boldsymbol{s}_1,\boldsymbol{s}_2,\overrightarrow{M_1M_2})=3$。

同一套语言也解释线性方程组：$r(A)=r(A\,|\,\boldsymbol{b})=3$ 即三张平面交于一点，$=2$ 即交于一条直线，$=1$ 即三面重合，$r(A)<r(A\,|\,\boldsymbol{b})$ 即无公共点。克拉默法则“系数行列式非零”的几何含义，就是三个法向量不共面。

\subsection{生成子空间}

\begin{definition}
设 $\boldsymbol{a}_1,\dots,\boldsymbol{a}_s$ 是 $n$ 维向量，其全体线性组合所成的集合
\[
\mathrm{span}\{\boldsymbol{a}_1,\dots,\boldsymbol{a}_s\}
=\bigl\{k_1\boldsymbol{a}_1+\cdots+k_s\boldsymbol{a}_s\ \big|\ k_i\in\mathbb{R}\bigr\}
\]
称为由它们\textbf{生成}（张成）的向量空间，也称\textbf{生成子空间}。
\end{definition}

\begin{theorem}
$\mathrm{span}\{\boldsymbol{a}_1,\dots,\boldsymbol{a}_s\}$ 是 $\mathbb{R}^{n}$ 的子空间，维数等于向量组的秩，任一极大线性无关组就是它的一组基；它是包含 $\boldsymbol{a}_1,\dots,\boldsymbol{a}_s$ 的\textbf{最小}子空间。
\end{theorem}

几何上：一个非零向量生成过原点的\textbf{直线}，两个不平行向量生成过原点的\textbf{平面}，三个不共面向量生成整个 $\mathbb{R}^{3}$。所以“三点共线”“四向量共面”本质上是生成子空间的维数降到 $1$ 或 $2$。

\corepoint{
（1）生成子空间一定\textbf{过原点}：不含原点的直线、平面都不是子空间。解析几何中的 $Ax+By+Cz+D=0$（$D\neq0$）不是向量空间，只有 $Ax+By+Cz=0$ 才是；谈“平面的基”须先平移到过原点。\\
（2）维数为 $r$ 不代表可以随便挑 $r$ 个向量作基，必须挑线性无关的一组；齐次方程组 $A\boldsymbol{x}=\boldsymbol{0}$ 的解空间维数为 $n-r(A)$，它与 $A$ 的行空间互为正交补。
}

\section{二次型与二次曲面判别对照}

第5章已用正交变换把二次型化为标准形。设 $f(\boldsymbol{x})=\boldsymbol{x}^{T}A\boldsymbol{x}$（$A^{T}=A$，$\boldsymbol{x}=(x,y,z)^{T}$），$A$ 的特征值为 $\lambda_1,\lambda_2,\lambda_3$，正、负特征值个数记为 $p,q$（$p+q=r(A)$）。由惯性定理，$p,q$ 与正交变换的选法无关。经 $\boldsymbol{x}=P\boldsymbol{y}$ 后 $f=1$ 化为 $\lambda_1x'^{2}+\lambda_2y'^{2}+\lambda_3z'^{2}=1$。

\begin{table}[htbp]
\centering
\renewcommand{\tablename}{表}
\caption{由特征值符号判别 $f(\boldsymbol{x})=1$ 与 $f(\boldsymbol{x})=0$ 的图形}
\begin{tabular}{@{}cll@{}}
\toprule
特征值符号 $(p,q)$ & $f(\boldsymbol{x})=1$ 的图形 & $f(\boldsymbol{x})=0$ 的图形 \\
\midrule
$(3,0)$（正定） & 椭球面 & 一个点（原点） \\
$(2,1)$ & 单叶双曲面 & 二次锥面 \\
$(1,2)$ & 双叶双曲面 & 二次锥面 \\
$(2,0)$ & 椭圆柱面 & 一条直线 \\
$(1,1)$ & 双曲柱面 & 两张相交平面 \\
$(1,0)$ & 两张平行平面 & 一张平面（二重） \\
$(0,q)$，$q\ge1$ & 无实图形 & 一个点或一张平面 \\
\bottomrule
\end{tabular}
\end{table}

若二次型退化（$r(A)\le2$）而方程中还含\textbf{一次项}，配方后会出现一个“抛物”型变量，得到\textbf{椭圆抛物面}（$p=2$）或\textbf{双曲抛物面}（$p=1,\,q=1$）。抛物面不可能由纯二次型 $f=1$ 得到，这是它与其他五类曲面的本质区别。

判别步骤：写出对称矩阵 $A$，求特征值（或用顺序主子式定符号）$\rightarrow$ 由 $(p,q)$ 与常数项定大类型 $\rightarrow$ 若有一次项先配方平移 $\rightarrow$ 必要时用截痕法核对（见附录A末节）。

\corepoint{
（1）判别依据是特征值的\textbf{符号}，不是数值：$2x^{2}+3y^{2}+z^{2}=1$ 与 $x^{2}+y^{2}+z^{2}=1$ 都是椭球面。\\
（2）只数“平方项有两个”就断言是柱面是常见错误：要看有无一次项。抛物面必须有一次项；$r(A)$ 则决定曲面“退化到几维”。
}

\section{向量方法在立体几何中的用法}

解析几何与线性代数合起来，给立体几何提供了一条“不算几何、只算代数”的通路。固定套路是三步：\textbf{建系、取点、算向量}。原点优先取对称中心或直角顶点，坐标轴取三条两两垂直的棱，并保证是右手系；把各关键点的坐标写出，未知的用参数表示。原则是让尽量多的点落在坐标轴或坐标面上。

\textbf{证明平行与垂直}：两直线平行 $\iff$ 方向向量成比例；线面平行 $\iff\boldsymbol{s}\cdot\boldsymbol{n}=0$ 且线上一点不在面内；线面垂直 $\iff\boldsymbol{s}\parallel\boldsymbol{n}$；面面平行 $\iff$ 两法向量成比例；各种垂直 $\iff$ 相应向量数量积为零；四点共面 $\iff$ 以其中一点为起点的三向量混合积为零。

\textbf{求法向量}：平面内有两个不共线向量 $\boldsymbol{a},\boldsymbol{b}$ 时取 $\boldsymbol{n}=\boldsymbol{a}\times\boldsymbol{b}$；已知三点则取 $\boldsymbol{n}=\overrightarrow{M_1M_2}\times\overrightarrow{M_1M_3}$。也可设 $\boldsymbol{n}=(x,y,z)$ 由 $\boldsymbol{n}\cdot\boldsymbol{a}=\boldsymbol{n}\cdot\boldsymbol{b}=0$ 解出，此时不要习惯性令 $x=1$（若法向量第一分量恰为零就会漏解），应令绝对值最大的分量为 $1$。

\textbf{求角}：异面直线所成角 $\cos\theta=\dfrac{|\boldsymbol{s}_1\cdot\boldsymbol{s}_2|}{|\boldsymbol{s}_1||\boldsymbol{s}_2|}$；直线与平面所成角 $\sin\varphi=\dfrac{|\boldsymbol{s}\cdot\boldsymbol{n}|}{|\boldsymbol{s}||\boldsymbol{n}|}$，两者都取 $\left[0,\frac{\pi}{2}\right]$。二面角先求两半平面的法向量，余弦的绝对值等于 $\dfrac{|\boldsymbol{n}_1\cdot\boldsymbol{n}_2|}{|\boldsymbol{n}_1||\boldsymbol{n}_2|}$，\textbf{锐钝须回原图判断}。各种距离见附录A「位置关系与距离」一节。

\corepoint{
（1）向量方法只解决“算”，不解决“选路”：二面角是锐是钝、点是否在平面内，都须看图或另行验证，公式本身不会告诉你。\\
（2）建系必须是右手系，$z$ 轴取反会使叉积、混合积全部变号，距离不变但角度判断会错；法向量可取任意非零倍数，不必“归一化”到唯一形式。
}
），不要写 \documentclass、\begin{document}、\end{document}
% 附录B 拓展与补充：整理原残稿（伴随矩阵的秩、降阶公式），并把原稿空缺的几处
% 按同一主题写实（秩与线性相关性、生成子空间、二次曲面判别、向量方法）。
% 与教材重复的部分只作指引，不重复推导。不写例题与习题。

\section{矩阵秩的补充结论}
\noindent\textit{本节的结论第3、5章都已用到，但教材未单独成节，这里集中补出，并说明条件与推导思路。}

\subsection{伴随矩阵的秩}

\begin{theorem}
设 $A$ 为 $n$ 阶方阵（$n\ge2$），$A^{*}$ 为其伴随矩阵，则
\[
r(A^{*})=
\begin{cases}
n, & r(A)=n,\\[1mm]
1, & r(A)=n-1,\\[1mm]
0, & r(A)\le n-2.
\end{cases}
\]
\end{theorem}

推导思路（三段都要用 $AA^{*}=A^{*}A=|A|E$）：

\begin{enumerate}
\item $r(A)=n$：$A$ 可逆，$|A|\neq0$，由 $A^{*}=|A|A^{-1}$ 知 $A^{*}$ 可逆，故 $r(A^{*})=n$；
\item $r(A)=n-1$：$|A|=0$，故 $AA^{*}=O$，由 $r(A)+r(A^{*})\le n$ 得 $r(A^{*})\le1$；又此时 $A$ 有非零的 $n-1$ 阶子式，即 $A^{*}\neq O$，故 $r(A^{*})=1$；
\item $r(A)\le n-2$：所有 $n-1$ 阶子式为零，$A^{*}=O$，故 $r(A^{*})=0$。
\end{enumerate}

\corepoint{
（1）中间一档最容易记成 $n-1$：$r(A)=n-1$ 时 $A^{*}$ 的秩\textbf{恒为 $1$}，与 $n$ 无关。\\
（2）结论的三个分界，正是“最高阶非零子式的阶数”这一秩的定义的直接体现；$n=1$ 时约定 $A^{*}=1$，公式不适用。
}

\subsection{降阶公式}

\begin{theorem}
设 $A$ 是 $m\times n$ 矩阵，$B$ 是 $n\times m$ 矩阵，$m>n$，$\lambda$ 为任意数，则
\[
|\lambda E_m-AB|=\lambda^{m-n}\,|\lambda E_n-BA|.
\]
\end{theorem}

即 $AB$（$m$ 阶）与 $BA$（$n$ 阶）的特征多项式只差因子 $\lambda^{m-n}$，故它们的\textbf{非零特征值完全相同（含重数）}。又 $r(AB)\le n<m$，$AB$ 必奇异，$0$ 是 $AB$ 的特征值，而 $BA$ 未必以 $0$ 为特征值——差别正好被 $\lambda^{m-n}$ 吸收。

\begin{corollary}
取 $\lambda=1$ 得 $|E_m-AB|=|E_n-BA|$，可把高阶行列式化为低阶行列式计算。
\end{corollary}

\corepoint{
（1）公式要求 $m>n$，且乘法方向不能颠倒：$AB$ 是 $m$ 阶、$BA$ 是 $n$ 阶。\\
（2）$\lambda=0$ 时两边同为 $0$，等式成立但不提供信息；求特征值须取 $\lambda\neq0$ 讨论。$m=n$ 时退化为恒等式。
}

\section{线性代数与解析几何的接口}

\subsection{秩与线性相关性}

矩阵的秩、向量组的秩、线性相关性三者是同一件事的三种说法。

\begin{theorem}
设 $A=(\boldsymbol{a}_1,\dots,\boldsymbol{a}_s)$ 是以向量组为列的矩阵，则
\begin{enumerate}
\item $\boldsymbol{a}_1,\dots,\boldsymbol{a}_s$ 线性相关 $\iff r(A)<s$；线性无关 $\iff r(A)=s$；
\item $r(\boldsymbol{a}_1,\dots,\boldsymbol{a}_s)=r(A)$，即极大线性无关组所含向量的个数；
\item 当向量个数等于维数（$s=n$）时，线性相关 $\iff|A|=0$。
\end{enumerate}
\end{theorem}

搬到三维空间，就得到解析几何中一串“行列式判零”的判据。设 $\boldsymbol{a},\boldsymbol{b},\boldsymbol{c}$ 为三维向量，则

\keyformula{\boldsymbol{a},\boldsymbol{b},\boldsymbol{c}\ \text{共面}\iff[\boldsymbol{a}\boldsymbol{b}\boldsymbol{c}]=0\iff
\begin{vmatrix}
a_x&a_y&a_z\\ b_x&b_y&b_z\\ c_x&c_y&c_z
\end{vmatrix}=0\iff r(\boldsymbol{a},\boldsymbol{b},\boldsymbol{c})<3}

因为混合积正是以三向量为行的三阶行列式（见附录A），而“共面”就是“线性相关”。同理：三点共线 $\iff$ 两方向向量线性相关（秩为 $1$）；四点 $M_0,M_1,M_2,M_3$ 共面 $\iff\overrightarrow{M_0M_1},\overrightarrow{M_0M_2},\overrightarrow{M_0M_3}$ 的行列式为零；两直线异面 $\iff r(\boldsymbol{s}_1,\boldsymbol{s}_2,\overrightarrow{M_1M_2})=3$。

同一套语言也解释线性方程组：$r(A)=r(A\,|\,\boldsymbol{b})=3$ 即三张平面交于一点，$=2$ 即交于一条直线，$=1$ 即三面重合，$r(A)<r(A\,|\,\boldsymbol{b})$ 即无公共点。克拉默法则“系数行列式非零”的几何含义，就是三个法向量不共面。

\subsection{生成子空间}

\begin{definition}
设 $\boldsymbol{a}_1,\dots,\boldsymbol{a}_s$ 是 $n$ 维向量，其全体线性组合所成的集合
\[
\mathrm{span}\{\boldsymbol{a}_1,\dots,\boldsymbol{a}_s\}
=\bigl\{k_1\boldsymbol{a}_1+\cdots+k_s\boldsymbol{a}_s\ \big|\ k_i\in\mathbb{R}\bigr\}
\]
称为由它们\textbf{生成}（张成）的向量空间，也称\textbf{生成子空间}。
\end{definition}

\begin{theorem}
$\mathrm{span}\{\boldsymbol{a}_1,\dots,\boldsymbol{a}_s\}$ 是 $\mathbb{R}^{n}$ 的子空间，维数等于向量组的秩，任一极大线性无关组就是它的一组基；它是包含 $\boldsymbol{a}_1,\dots,\boldsymbol{a}_s$ 的\textbf{最小}子空间。
\end{theorem}

几何上：一个非零向量生成过原点的\textbf{直线}，两个不平行向量生成过原点的\textbf{平面}，三个不共面向量生成整个 $\mathbb{R}^{3}$。所以“三点共线”“四向量共面”本质上是生成子空间的维数降到 $1$ 或 $2$。

\corepoint{
（1）生成子空间一定\textbf{过原点}：不含原点的直线、平面都不是子空间。解析几何中的 $Ax+By+Cz+D=0$（$D\neq0$）不是向量空间，只有 $Ax+By+Cz=0$ 才是；谈“平面的基”须先平移到过原点。\\
（2）维数为 $r$ 不代表可以随便挑 $r$ 个向量作基，必须挑线性无关的一组；齐次方程组 $A\boldsymbol{x}=\boldsymbol{0}$ 的解空间维数为 $n-r(A)$，它与 $A$ 的行空间互为正交补。
}

\section{二次型与二次曲面判别对照}

第5章已用正交变换把二次型化为标准形。设 $f(\boldsymbol{x})=\boldsymbol{x}^{T}A\boldsymbol{x}$（$A^{T}=A$，$\boldsymbol{x}=(x,y,z)^{T}$），$A$ 的特征值为 $\lambda_1,\lambda_2,\lambda_3$，正、负特征值个数记为 $p,q$（$p+q=r(A)$）。由惯性定理，$p,q$ 与正交变换的选法无关。经 $\boldsymbol{x}=P\boldsymbol{y}$ 后 $f=1$ 化为 $\lambda_1x'^{2}+\lambda_2y'^{2}+\lambda_3z'^{2}=1$。

\begin{table}[htbp]
\centering
\renewcommand{\tablename}{表}
\caption{由特征值符号判别 $f(\boldsymbol{x})=1$ 与 $f(\boldsymbol{x})=0$ 的图形}
\begin{tabular}{@{}cll@{}}
\toprule
特征值符号 $(p,q)$ & $f(\boldsymbol{x})=1$ 的图形 & $f(\boldsymbol{x})=0$ 的图形 \\
\midrule
$(3,0)$（正定） & 椭球面 & 一个点（原点） \\
$(2,1)$ & 单叶双曲面 & 二次锥面 \\
$(1,2)$ & 双叶双曲面 & 二次锥面 \\
$(2,0)$ & 椭圆柱面 & 一条直线 \\
$(1,1)$ & 双曲柱面 & 两张相交平面 \\
$(1,0)$ & 两张平行平面 & 一张平面（二重） \\
$(0,q)$，$q\ge1$ & 无实图形 & 一个点或一张平面 \\
\bottomrule
\end{tabular}
\end{table}

若二次型退化（$r(A)\le2$）而方程中还含\textbf{一次项}，配方后会出现一个“抛物”型变量，得到\textbf{椭圆抛物面}（$p=2$）或\textbf{双曲抛物面}（$p=1,\,q=1$）。抛物面不可能由纯二次型 $f=1$ 得到，这是它与其他五类曲面的本质区别。

判别步骤：写出对称矩阵 $A$，求特征值（或用顺序主子式定符号）$\rightarrow$ 由 $(p,q)$ 与常数项定大类型 $\rightarrow$ 若有一次项先配方平移 $\rightarrow$ 必要时用截痕法核对（见附录A末节）。

\corepoint{
（1）判别依据是特征值的\textbf{符号}，不是数值：$2x^{2}+3y^{2}+z^{2}=1$ 与 $x^{2}+y^{2}+z^{2}=1$ 都是椭球面。\\
（2）只数“平方项有两个”就断言是柱面是常见错误：要看有无一次项。抛物面必须有一次项；$r(A)$ 则决定曲面“退化到几维”。
}

\section{向量方法在立体几何中的用法}

解析几何与线性代数合起来，给立体几何提供了一条“不算几何、只算代数”的通路。固定套路是三步：\textbf{建系、取点、算向量}。原点优先取对称中心或直角顶点，坐标轴取三条两两垂直的棱，并保证是右手系；把各关键点的坐标写出，未知的用参数表示。原则是让尽量多的点落在坐标轴或坐标面上。

\textbf{证明平行与垂直}：两直线平行 $\iff$ 方向向量成比例；线面平行 $\iff\boldsymbol{s}\cdot\boldsymbol{n}=0$ 且线上一点不在面内；线面垂直 $\iff\boldsymbol{s}\parallel\boldsymbol{n}$；面面平行 $\iff$ 两法向量成比例；各种垂直 $\iff$ 相应向量数量积为零；四点共面 $\iff$ 以其中一点为起点的三向量混合积为零。

\textbf{求法向量}：平面内有两个不共线向量 $\boldsymbol{a},\boldsymbol{b}$ 时取 $\boldsymbol{n}=\boldsymbol{a}\times\boldsymbol{b}$；已知三点则取 $\boldsymbol{n}=\overrightarrow{M_1M_2}\times\overrightarrow{M_1M_3}$。也可设 $\boldsymbol{n}=(x,y,z)$ 由 $\boldsymbol{n}\cdot\boldsymbol{a}=\boldsymbol{n}\cdot\boldsymbol{b}=0$ 解出，此时不要习惯性令 $x=1$（若法向量第一分量恰为零就会漏解），应令绝对值最大的分量为 $1$。

\textbf{求角}：异面直线所成角 $\cos\theta=\dfrac{|\boldsymbol{s}_1\cdot\boldsymbol{s}_2|}{|\boldsymbol{s}_1||\boldsymbol{s}_2|}$；直线与平面所成角 $\sin\varphi=\dfrac{|\boldsymbol{s}\cdot\boldsymbol{n}|}{|\boldsymbol{s}||\boldsymbol{n}|}$，两者都取 $\left[0,\frac{\pi}{2}\right]$。二面角先求两半平面的法向量，余弦的绝对值等于 $\dfrac{|\boldsymbol{n}_1\cdot\boldsymbol{n}_2|}{|\boldsymbol{n}_1||\boldsymbol{n}_2|}$，\textbf{锐钝须回原图判断}。各种距离见附录A「位置关系与距离」一节。

\corepoint{
（1）向量方法只解决“算”，不解决“选路”：二面角是锐是钝、点是否在平面内，都须看图或另行验证，公式本身不会告诉你。\\
（2）建系必须是右手系，$z$ 轴取反会使叉积、混合积全部变号，距离不变但角度判断会错；法向量可取任意非零倍数，不必“归一化”到唯一形式。
}
），不要写 \documentclass、\begin{document}、\end{document}
% 附录B 拓展与补充：整理原残稿（伴随矩阵的秩、降阶公式），并把原稿空缺的几处
% 按同一主题写实（秩与线性相关性、生成子空间、二次曲面判别、向量方法）。
% 与教材重复的部分只作指引，不重复推导。不写例题与习题。

\section{矩阵秩的补充结论}
\noindent\textit{本节的结论第3、5章都已用到，但教材未单独成节，这里集中补出，并说明条件与推导思路。}

\subsection{伴随矩阵的秩}

\begin{theorem}
设 $A$ 为 $n$ 阶方阵（$n\ge2$），$A^{*}$ 为其伴随矩阵，则
\[
r(A^{*})=
\begin{cases}
n, & r(A)=n,\\[1mm]
1, & r(A)=n-1,\\[1mm]
0, & r(A)\le n-2.
\end{cases}
\]
\end{theorem}

推导思路（三段都要用 $AA^{*}=A^{*}A=|A|E$）：

\begin{enumerate}
\item $r(A)=n$：$A$ 可逆，$|A|\neq0$，由 $A^{*}=|A|A^{-1}$ 知 $A^{*}$ 可逆，故 $r(A^{*})=n$；
\item $r(A)=n-1$：$|A|=0$，故 $AA^{*}=O$，由 $r(A)+r(A^{*})\le n$ 得 $r(A^{*})\le1$；又此时 $A$ 有非零的 $n-1$ 阶子式，即 $A^{*}\neq O$，故 $r(A^{*})=1$；
\item $r(A)\le n-2$：所有 $n-1$ 阶子式为零，$A^{*}=O$，故 $r(A^{*})=0$。
\end{enumerate}

\corepoint{
（1）中间一档最容易记成 $n-1$：$r(A)=n-1$ 时 $A^{*}$ 的秩\textbf{恒为 $1$}，与 $n$ 无关。\\
（2）结论的三个分界，正是“最高阶非零子式的阶数”这一秩的定义的直接体现；$n=1$ 时约定 $A^{*}=1$，公式不适用。
}

\subsection{降阶公式}

\begin{theorem}
设 $A$ 是 $m\times n$ 矩阵，$B$ 是 $n\times m$ 矩阵，$m>n$，$\lambda$ 为任意数，则
\[
|\lambda E_m-AB|=\lambda^{m-n}\,|\lambda E_n-BA|.
\]
\end{theorem}

即 $AB$（$m$ 阶）与 $BA$（$n$ 阶）的特征多项式只差因子 $\lambda^{m-n}$，故它们的\textbf{非零特征值完全相同（含重数）}。又 $r(AB)\le n<m$，$AB$ 必奇异，$0$ 是 $AB$ 的特征值，而 $BA$ 未必以 $0$ 为特征值——差别正好被 $\lambda^{m-n}$ 吸收。

\begin{corollary}
取 $\lambda=1$ 得 $|E_m-AB|=|E_n-BA|$，可把高阶行列式化为低阶行列式计算。
\end{corollary}

\corepoint{
（1）公式要求 $m>n$，且乘法方向不能颠倒：$AB$ 是 $m$ 阶、$BA$ 是 $n$ 阶。\\
（2）$\lambda=0$ 时两边同为 $0$，等式成立但不提供信息；求特征值须取 $\lambda\neq0$ 讨论。$m=n$ 时退化为恒等式。
}

\section{线性代数与解析几何的接口}

\subsection{秩与线性相关性}

矩阵的秩、向量组的秩、线性相关性三者是同一件事的三种说法。

\begin{theorem}
设 $A=(\boldsymbol{a}_1,\dots,\boldsymbol{a}_s)$ 是以向量组为列的矩阵，则
\begin{enumerate}
\item $\boldsymbol{a}_1,\dots,\boldsymbol{a}_s$ 线性相关 $\iff r(A)<s$；线性无关 $\iff r(A)=s$；
\item $r(\boldsymbol{a}_1,\dots,\boldsymbol{a}_s)=r(A)$，即极大线性无关组所含向量的个数；
\item 当向量个数等于维数（$s=n$）时，线性相关 $\iff|A|=0$。
\end{enumerate}
\end{theorem}

搬到三维空间，就得到解析几何中一串“行列式判零”的判据。设 $\boldsymbol{a},\boldsymbol{b},\boldsymbol{c}$ 为三维向量，则

\keyformula{\boldsymbol{a},\boldsymbol{b},\boldsymbol{c}\ \text{共面}\iff[\boldsymbol{a}\boldsymbol{b}\boldsymbol{c}]=0\iff
\begin{vmatrix}
a_x&a_y&a_z\\ b_x&b_y&b_z\\ c_x&c_y&c_z
\end{vmatrix}=0\iff r(\boldsymbol{a},\boldsymbol{b},\boldsymbol{c})<3}

因为混合积正是以三向量为行的三阶行列式（见附录A），而“共面”就是“线性相关”。同理：三点共线 $\iff$ 两方向向量线性相关（秩为 $1$）；四点 $M_0,M_1,M_2,M_3$ 共面 $\iff\overrightarrow{M_0M_1},\overrightarrow{M_0M_2},\overrightarrow{M_0M_3}$ 的行列式为零；两直线异面 $\iff r(\boldsymbol{s}_1,\boldsymbol{s}_2,\overrightarrow{M_1M_2})=3$。

同一套语言也解释线性方程组：$r(A)=r(A\,|\,\boldsymbol{b})=3$ 即三张平面交于一点，$=2$ 即交于一条直线，$=1$ 即三面重合，$r(A)<r(A\,|\,\boldsymbol{b})$ 即无公共点。克拉默法则“系数行列式非零”的几何含义，就是三个法向量不共面。

\subsection{生成子空间}

\begin{definition}
设 $\boldsymbol{a}_1,\dots,\boldsymbol{a}_s$ 是 $n$ 维向量，其全体线性组合所成的集合
\[
\mathrm{span}\{\boldsymbol{a}_1,\dots,\boldsymbol{a}_s\}
=\bigl\{k_1\boldsymbol{a}_1+\cdots+k_s\boldsymbol{a}_s\ \big|\ k_i\in\mathbb{R}\bigr\}
\]
称为由它们\textbf{生成}（张成）的向量空间，也称\textbf{生成子空间}。
\end{definition}

\begin{theorem}
$\mathrm{span}\{\boldsymbol{a}_1,\dots,\boldsymbol{a}_s\}$ 是 $\mathbb{R}^{n}$ 的子空间，维数等于向量组的秩，任一极大线性无关组就是它的一组基；它是包含 $\boldsymbol{a}_1,\dots,\boldsymbol{a}_s$ 的\textbf{最小}子空间。
\end{theorem}

几何上：一个非零向量生成过原点的\textbf{直线}，两个不平行向量生成过原点的\textbf{平面}，三个不共面向量生成整个 $\mathbb{R}^{3}$。所以“三点共线”“四向量共面”本质上是生成子空间的维数降到 $1$ 或 $2$。

\corepoint{
（1）生成子空间一定\textbf{过原点}：不含原点的直线、平面都不是子空间。解析几何中的 $Ax+By+Cz+D=0$（$D\neq0$）不是向量空间，只有 $Ax+By+Cz=0$ 才是；谈“平面的基”须先平移到过原点。\\
（2）维数为 $r$ 不代表可以随便挑 $r$ 个向量作基，必须挑线性无关的一组；齐次方程组 $A\boldsymbol{x}=\boldsymbol{0}$ 的解空间维数为 $n-r(A)$，它与 $A$ 的行空间互为正交补。
}

\section{二次型与二次曲面判别对照}

第5章已用正交变换把二次型化为标准形。设 $f(\boldsymbol{x})=\boldsymbol{x}^{T}A\boldsymbol{x}$（$A^{T}=A$，$\boldsymbol{x}=(x,y,z)^{T}$），$A$ 的特征值为 $\lambda_1,\lambda_2,\lambda_3$，正、负特征值个数记为 $p,q$（$p+q=r(A)$）。由惯性定理，$p,q$ 与正交变换的选法无关。经 $\boldsymbol{x}=P\boldsymbol{y}$ 后 $f=1$ 化为 $\lambda_1x'^{2}+\lambda_2y'^{2}+\lambda_3z'^{2}=1$。

\begin{table}[htbp]
\centering
\renewcommand{\tablename}{表}
\caption{由特征值符号判别 $f(\boldsymbol{x})=1$ 与 $f(\boldsymbol{x})=0$ 的图形}
\begin{tabular}{@{}cll@{}}
\toprule
特征值符号 $(p,q)$ & $f(\boldsymbol{x})=1$ 的图形 & $f(\boldsymbol{x})=0$ 的图形 \\
\midrule
$(3,0)$（正定） & 椭球面 & 一个点（原点） \\
$(2,1)$ & 单叶双曲面 & 二次锥面 \\
$(1,2)$ & 双叶双曲面 & 二次锥面 \\
$(2,0)$ & 椭圆柱面 & 一条直线 \\
$(1,1)$ & 双曲柱面 & 两张相交平面 \\
$(1,0)$ & 两张平行平面 & 一张平面（二重） \\
$(0,q)$，$q\ge1$ & 无实图形 & 一个点或一张平面 \\
\bottomrule
\end{tabular}
\end{table}

若二次型退化（$r(A)\le2$）而方程中还含\textbf{一次项}，配方后会出现一个“抛物”型变量，得到\textbf{椭圆抛物面}（$p=2$）或\textbf{双曲抛物面}（$p=1,\,q=1$）。抛物面不可能由纯二次型 $f=1$ 得到，这是它与其他五类曲面的本质区别。

判别步骤：写出对称矩阵 $A$，求特征值（或用顺序主子式定符号）$\rightarrow$ 由 $(p,q)$ 与常数项定大类型 $\rightarrow$ 若有一次项先配方平移 $\rightarrow$ 必要时用截痕法核对（见附录A末节）。

\corepoint{
（1）判别依据是特征值的\textbf{符号}，不是数值：$2x^{2}+3y^{2}+z^{2}=1$ 与 $x^{2}+y^{2}+z^{2}=1$ 都是椭球面。\\
（2）只数“平方项有两个”就断言是柱面是常见错误：要看有无一次项。抛物面必须有一次项；$r(A)$ 则决定曲面“退化到几维”。
}

\section{向量方法在立体几何中的用法}

解析几何与线性代数合起来，给立体几何提供了一条“不算几何、只算代数”的通路。固定套路是三步：\textbf{建系、取点、算向量}。原点优先取对称中心或直角顶点，坐标轴取三条两两垂直的棱，并保证是右手系；把各关键点的坐标写出，未知的用参数表示。原则是让尽量多的点落在坐标轴或坐标面上。

\textbf{证明平行与垂直}：两直线平行 $\iff$ 方向向量成比例；线面平行 $\iff\boldsymbol{s}\cdot\boldsymbol{n}=0$ 且线上一点不在面内；线面垂直 $\iff\boldsymbol{s}\parallel\boldsymbol{n}$；面面平行 $\iff$ 两法向量成比例；各种垂直 $\iff$ 相应向量数量积为零；四点共面 $\iff$ 以其中一点为起点的三向量混合积为零。

\textbf{求法向量}：平面内有两个不共线向量 $\boldsymbol{a},\boldsymbol{b}$ 时取 $\boldsymbol{n}=\boldsymbol{a}\times\boldsymbol{b}$；已知三点则取 $\boldsymbol{n}=\overrightarrow{M_1M_2}\times\overrightarrow{M_1M_3}$。也可设 $\boldsymbol{n}=(x,y,z)$ 由 $\boldsymbol{n}\cdot\boldsymbol{a}=\boldsymbol{n}\cdot\boldsymbol{b}=0$ 解出，此时不要习惯性令 $x=1$（若法向量第一分量恰为零就会漏解），应令绝对值最大的分量为 $1$。

\textbf{求角}：异面直线所成角 $\cos\theta=\dfrac{|\boldsymbol{s}_1\cdot\boldsymbol{s}_2|}{|\boldsymbol{s}_1||\boldsymbol{s}_2|}$；直线与平面所成角 $\sin\varphi=\dfrac{|\boldsymbol{s}\cdot\boldsymbol{n}|}{|\boldsymbol{s}||\boldsymbol{n}|}$，两者都取 $\left[0,\frac{\pi}{2}\right]$。二面角先求两半平面的法向量，余弦的绝对值等于 $\dfrac{|\boldsymbol{n}_1\cdot\boldsymbol{n}_2|}{|\boldsymbol{n}_1||\boldsymbol{n}_2|}$，\textbf{锐钝须回原图判断}。各种距离见附录A「位置关系与距离」一节。

\corepoint{
（1）向量方法只解决“算”，不解决“选路”：二面角是锐是钝、点是否在平面内，都须看图或另行验证，公式本身不会告诉你。\\
（2）建系必须是右手系，$z$ 轴取反会使叉积、混合积全部变号，距离不变但角度判断会错；法向量可取任意非零倍数，不必“归一化”到唯一形式。
}


\end{document}
